% MAT670 -- Packings, Lattices and Configurations -- shared preamble
\documentclass[11pt,openany]{book}
\usepackage[margin=1.1in,marginparwidth=0.9in]{geometry}
\usepackage{amsmath,amssymb,amsthm,mathtools}
\usepackage[T1]{fontenc}
\usepackage{enumitem}
\usepackage{xcolor}
\usepackage{tikz}
\usetikzlibrary{calc,arrows.meta,patterns,decorations.pathmorphing}
\usepackage{graphicx}
\usepackage{booktabs}
\usepackage[colorlinks=true,linkcolor=blue!50!black,citecolor=green!40!black,urlcolor=blue!50!black]{hyperref}

% ---------- theorem environments ----------
\theoremstyle{plain}
\newtheorem{theorem}{Theorem}[chapter]
\newtheorem{lemma}[theorem]{Lemma}
\newtheorem{proposition}[theorem]{Proposition}
\newtheorem{corollary}[theorem]{Corollary}
\newtheorem{conjecture}[theorem]{Conjecture}
\newtheorem{claim}[theorem]{Claim}
\theoremstyle{definition}
\newtheorem{definition}[theorem]{Definition}
\newtheorem{example}[theorem]{Example}
\newtheorem{exercise}[theorem]{Exercise}
\newtheorem{question}[theorem]{Question}
\newtheorem{problem}[theorem]{Problem}
\theoremstyle{remark}
\newtheorem{remark}[theorem]{Remark}

% ---------- editorial markup (first-pass) ----------
% \added{...}   material NOT in the handwritten notes (gap filled by editor)
% \fixed{...}   an error in the notes that was corrected; argument says what changed
% \unclear{...} a reading of the handwriting that is uncertain, or a gap not resolved
% \orig{...}    short note on what the original page said / showed (e.g. "sketch omitted")
\definecolor{addcol}{RGB}{0,100,160}
\definecolor{flagcol}{RGB}{190,60,0}
\newcommand{\added}[1]{{\color{addcol}#1}}
\newcommand{\fixed}[1]{\marginpar{\raggedright\scriptsize\color{flagcol}\textbf{Fix:} #1}}
\newcommand{\unclear}[1]{\marginpar{\raggedright\scriptsize\color{flagcol}\textbf{?} #1}}
\newcommand{\orig}[1]{\marginpar{\raggedright\scriptsize\color{gray}#1}}
\newenvironment{addedblock}{\par\color{addcol}\noindent\textsf{\small[Added]}\ }{\par}

% ---------- math macros ----------
\newcommand{\R}{\mathbb{R}}
\newcommand{\Z}{\mathbb{Z}}
\newcommand{\Q}{\mathbb{Q}}
\newcommand{\C}{\mathbb{C}}
\newcommand{\N}{\mathbb{N}}
\newcommand{\T}{\mathbb{T}}
\newcommand{\Sph}{\mathbb{S}}          % sphere S^{n-1}
\newcommand{\Lat}{\Lambda}             % a lattice
\newcommand{\vol}{\operatorname{vol}}
\newcommand{\covol}{\operatorname{covol}}
\newcommand{\dist}{\operatorname{dist}}
\newcommand{\diam}{\operatorname{diam}}
\newcommand{\conv}{\operatorname{conv}}
\newcommand{\interior}{\operatorname{int}}
\newcommand{\GL}{\mathrm{GL}}
\newcommand{\SL}{\mathrm{SL}}
\newcommand{\SO}{\mathrm{SO}}
\newcommand{\Orth}{\mathrm{O}}
\newcommand{\Conf}{\operatorname{Conf}}
\newcommand{\Hess}{\operatorname{Hess}}
\newcommand{\ind}{\operatorname{ind}}
\newcommand{\abs}[1]{\left\lvert #1 \right\rvert}
\newcommand{\norm}[1]{\left\lVert #1 \right\rVert}
\newcommand{\ip}[2]{\langle #1, #2\rangle}
\newcommand{\dens}{\delta}             % packing density
\newcommand{\cent}{\theta}             % center density
\newcommand{\kiss}{\tau}               % kissing number
\newcommand{\Ball}{B}                  % ball


\title{\textbf{Packings, Lattices and Configurations}\\[0.6em]
\large Lecture notes for MAT670, TU Graz, Summer Term 2016\\[0.3em]
\normalsize First typeset draft from the handwritten notes}
\author{Woden Kusner}
\date{Draft of \today}

\begin{document}
\frontmatter
\maketitle

\chapter*{About this draft}
\addcontentsline{toc}{chapter}{About this draft}

These notes are a first typeset pass at the handwritten notes for the course
\emph{Packings, Lattices and Configurations} (MAT670), taught in eleven two-hour lectures at TU Graz
in the summer term of 2016. The course description read:
\begin{quote}
This course will introduce a variety of different tools and models to explore some of the theory of
geometric configuration spaces that underlie physical systems. \dots\ By the end of the course, you
should have some familiarity with the geometry and topology of configuration spaces, Morse theory, the
geometry of lattices, maximal density arrangements, random packings and statistical mechanics.
\end{quote}

The structure, order, examples and choice of proofs follow the handwritten notes. Telegraphic notes have
been written out as prose, sketches have been redrawn, and gaps have been filled. The colour code is:

\begin{itemize}
\item Black text is a cleaned-up rendering of what is on the page.
\item \added{Blue text} was \emph{not} in the handwritten notes. It covers completed definitions and
  proofs, missing hypotheses, references, and remarks on results that appeared after 2016.
  Longer insertions are marked \added{\textsf{\small[Added]}}.
\item Orange margin notes \textbf{Fix:} mark places where the notes contained an error that has been
  corrected, saying what changed. \textbf{?} marks a doubtful reading of the handwriting,
  or a step that is still open.
\item Grey margin notes describe the original page (crossed-out text, omitted doodles, and so on).
\end{itemize}

\paragraph{The interlude.} The chapter ``Interlude: From Dynkin diagrams to Lie algebras and Lie groups'',
placed after Lecture~6, is entirely new. It is set in black for readability but is all added material.

\paragraph{Sources.} The scans used are \texttt{MAT670\_1}, \texttt{\_2}, \texttt{\_3},
\texttt{\_46}, \texttt{\_78}, \texttt{\_9X} and \texttt{\_XI} on the course web page.
\texttt{MAT670\_emsemble} is a bundled copy of Lectures 1--3.
\texttt{MAT670\_classification} is an earlier scan of Lectures 4--5, the first 31 pages of
\texttt{MAT670\_46} minus one recap page. Both were transcribed independently and used only as a
cross-check.

\paragraph{Known open items in this draft.}
\begin{itemize}
\item Lecture 3: the ``improved'' radial Blichfeldt gauge is stated in the notes without proof. No proof
  has been supplied here; a numerical search in low dimensions found no counterexample.
\item Lecture 10: the deformation theorem for $0<\theta<\pi/N$ had no proof in the notes. The proof given
  in blue is the editor's and should be checked. Pages 24--26 of the scan are mostly headings.
  The notes themselves mark this lecture ``needs to be rewritten, out of order / missing pieces''.
\item Lecture 11: the count of locally jammed configurations in the open-ended channel may be off at the
  ends, giving $\binom{\ell-k-1}{k}$ rather than $\binom{\ell-k+1}{k}$. It is flagged but not changed.
\end{itemize}

\tableofcontents

\mainmatter
\part{Packing bounds}
% ============================================================
% MAT670 -- Lecture 1 (source: MAT670_1, pp. 1--17)
% ============================================================
\chapter{Lecture 1: Packing problems}
\label{L1:chap}

\section{Course format}
\label{L1:sec:admin}
\orig{1.0 ``Admin'': contact, room, breaks omitted}
Evaluation is by short quizzes and a project: a talk of at least 20 minutes
together with a written summary. Updates are posted by email, through TU Online,
and on the course page \texttt{wkusner.github.io/MAT670/}.

\section{Mathematics of packings, lattices and configurations}
\label{L1:sec:overview}

The subject of this course sits at the meeting point of several areas:
\begin{itemize}[nosep]
  \item classical convex and discrete geometry, and the geometry of numbers;
  \item applied topology and Morse theory;
  \item statistical mechanics, combinatorics and information theory.
\end{itemize}
\orig{p.\,1: a diagram linking these areas; ``Rigidity'' crossed out there}
We will take various ideas from each of these, all of them dealing with
\emph{configurations} of points or bodies. Related topics include rigidity theory
and configuration spaces of linkages, information theory (codes), and the geometry
of numbers. Packing problems are my main motivation.

Packing problems are typically \emph{easy to state but hard to solve}. The same
is true of the geometric inequalities and the configuration spaces we will meet
along the way.

\begin{remark}
\added{Some landmarks, for orientation: the densest packing of congruent balls
in $\R^3$ (the Kepler conjecture, density $\pi/\sqrt{18}$) was proved by Hales
(announced 1998, published 2005; formal verification by the Flyspeck project 2014/2017).
In dimension $8$ the $E_8$ lattice packing was proved optimal by Viazovska (2016),
and in dimension $24$ the Leech lattice packing by
Cohn--Kumar--Miller--Radchenko--Viazovska (2016/2017). These dimension $8$ and $24$
results appeared during the course itself. In every other dimension
$d\ge 4$ the problem is still open.}
\end{remark}

\section{Packing problems}
\label{L1:sec:packing}

\subsection{Packing congruent bodies in a box}

We start in two dimensions, with the problem of packing circles (disks) in a box.
Let $B\subset\R^2$ be a box (a square, say) and $K$ a disk. We want to place $N$
congruent copies of $K$ in $B$ so as to maximize the \emph{density}, or volume fraction,
\[
  \frac{\vol\bigl(\bigcup_{i=1}^N K_i\bigr)}{\vol(B)} .
\]

\begin{center}
\begin{tikzpicture}[scale=0.8]
  \draw[thick] (0,0) rectangle (4,4);
  \foreach \x/\y in {1/1,3/1,1/3,3/3} \draw[fill=blue!10] (\x,\y) circle (1);
  \node[right] at (4,2) {$B$};
\end{tikzpicture}
\\[2pt]
{\small \added{Four congruent disks packed in a square $B$ (this $2\times2$ arrangement is optimal for $N=4$).}}
\end{center}

\begin{definition}[Packing in a box]
\label{L1:def:packing-box}
A collection $K_1,\dots,K_N$ of subsets of $\R^n$ is a \emph{packing of $K$ in $B$} if
\begin{enumerate}[nosep]
  \item the $K_i$ are congruent to $K$: $K_i=\varphi_i K$ with
        $\varphi_i\in\R^n\rtimes\SO(n)$ a rigid motion\added{;}
  \item $K_i\subset B$ for all $i$\added{;}
  \item the interiors are pairwise disjoint: $\interior K_i\cap\interior K_j=\emptyset$
        for $i\neq j$.
\end{enumerate}
\end{definition}
\orig{p.\,3 writes $\R^2\rtimes\SO(2)$}

\begin{example}
\label{L1:ex:radius}
Two equivalent formulations of the disk-in-a-box problem:
\begin{enumerate}[nosep]
  \item Find the largest scale $r$ such that $\bigcup_{i=1}^N \varphi_i(rK)\subset B$ is
        a packing of $B$\added{, where $K$ is the unit disk}. (Here $rK$ is in Minkowski notation:
        $rK=\{rx : x\in K\}$.)
  \item Find a configuration $\{x_1,\dots,x_N\}$ of $N$ points in $B$ such that
        \[
          d(x_i,x_j)\ \ge\ 2r \quad\text{for } i\ne j,
          \qquad\text{and}\qquad d(x_i,\partial B)\ \ge\ r \quad\text{for all } i,
        \]
        with $r$ as large as possible.
\end{enumerate}
\fixed{notes: $d(x_i,x_j)>r$; disks of radius $r$ need centres $2r$ apart}
\added{The points $x_i$ are the centres of the disks $\varphi_i(rK)$.}
\end{example}

These are hard problems, and there is an infinite variety of them: we can change
$N$, $K$ and $B$, the dimension, add other constraints, and so on.

One case of particular interest is $B=\R^2$: the circle packing problem in the
plane. Other natural choices are $B=\R^d$, or $B=\Sph^2$.
\added{(For $B=\R^d$ the number of bodies is infinite and ``density'' must be
defined by a limit; this is done next.)}

\subsection{Packings of the plane}
\label{L1:subsec:plane}
\orig{numbered ``1.2'' a second time in the notes}

Now take $B=\R^2$. We have an infinite collection of congruent disks $K_i$, now
normalized to have unit radius, with $\interior K_i\cap\interior K_j=\emptyset$ for $i\ne j$.
Write $P=\bigcup_i K_i$.

\begin{definition}[Upper and lower density]
\label{L1:def:density}
The \emph{upper density} of the packing $P$ is
\[
  \dens^+(P)\;=\;\limsup_{R\to\infty}\frac{\vol(P\cap R\,\mathbb{D}^2)}{\vol(R\,\mathbb{D}^2)},
\]
where $\mathbb{D}^2$ is the closed unit disk \added{centred at the origin}.
The \emph{lower density} $\dens^-(P)$ is defined in the same way with $\liminf$.
In general dimension,
\[
  \dens^+(P)\;=\;\limsup_{R\to\infty}\frac{\vol(P\cap R\,\mathbb{D}^n)}{\vol(R\,\mathbb{D}^n)}
\]
\added{with $\mathbb{D}^n$ the unit ball in $\R^n$}.
\end{definition}

\begin{remark}
These quantities need not be nice in general \added{(e.g.\ $\dens^+(P)\neq\dens^-(P)$
is possible, and then ``the'' density of $P$ is not defined)}. However:
\begin{itemize}[nosep]
  \item they do \emph{not} depend on the choice of origin \added{(moving the centre
        changes the numerator only in a boundary layer of volume $O(R^{n-1})$)};
  \item they \emph{do} depend, in general, on the shape of the averaging region
        $\mathbb{D}^n$ \added{(for periodic packings, though, any reasonable shape gives the same value)}.
\end{itemize}
\end{remark}

Does a packing of maximal density exist? We might look at this in more detail later.
\added{(It does: the supremum of $\dens^+$ over all packings is attained, by a compactness
argument; see e.g.\ Groemer's theorem in Conway--Sloane, \emph{Sphere Packings, Lattices
and Groups}, Ch.~1.)}
\unclear{existence deferred in the notes}

\subsection{Bounds on the density}

How can we compute bounds on the density? Clearly $\dens^+(P)\le 1$, which is a bad bound.

We can also construct lower bounds using highly structured packings. For example,
the unit balls centred at the points of $(2\Z)^d$ form a packing. Since this structure
is periodic, we can work with the fundamental domain.

\begin{exercise}[Easy]
\label{L1:ex:periodic-density}
Why can we compute the density on a fundamental domain? \added{Show that for a periodic
packing, with period lattice $\Lat$ and $m$ balls of radius $1$ centred in a fundamental
domain $F$ of $\Lat$, the density exists (that is, $\dens^+=\dens^-$) and equals
$m\,\vol(\Ball^d)/\vol(F)$.}
\end{exercise}

Recall the volume of the ball of radius $r$ in $\R^d$:
\[
  V_d(r)\;=\;\vol\bigl(\Ball^d(r)\bigr)\;=\;\frac{\pi^{d/2}}{\Gamma\!\left(\frac d2+1\right)}\,r^d .
\]
This is easier to remember in even dimension: $V_{2k}(r)=\dfrac{\pi^k}{k!}\,r^{2k}$.
For $(2\Z)^d$ the fundamental domain is the cube $[0,2)^d$ of volume $2^d$, containing one
ball, so the density is
\[
  \dens\bigl((2\Z)^d\bigr)\;=\;\frac{V_d(1)}{2^d}\;=\;\frac{\pi^{d/2}}{\Gamma\!\left(\frac d2+1\right)2^d}
  \;\overset{d=2k}{=}\;\frac{\pi^{k}}{k!\,2^{2k}} .
\]
\orig{notes write $(\tfrac d2)!$ for $\Gamma(\tfrac d2+1)$}

We could improve this by considering better periodic systems.

\begin{exercise}[Easy]
\label{L1:ex:periodic-approx}
Show that periodic packings can approximate a densest packing: \added{the supremum of
$\dens^+(P)$ over all packings $P$ equals the supremum over periodic packings.
(Hint: cut a large cube out of a near-optimal packing, remove the balls that meet its
boundary, and repeat periodically.)}
\end{exercise}

But which periodic systems should we use?

\begin{exercise}
\label{L1:ex:root-lattices}
Consider the lattices coming from the root systems $A_1\times\cdots\times A_1$
($d$ factors), \dots, $A_d$. Compute their packing densities and compare.
\unclear{p.\,6: ``Consider $A^1$, $A_1\times^d A_1$ \dots\ $A_d$? root systems''}
\added{(For reference: $A_1^d$ gives $\Z^d$, which is $(2\Z)^d$ after scaling, with
density $V_d(1)/2^d$. The lattice $A_d$ has minimal norm $\sqrt2$ and covolume
$\sqrt{d+1}$, hence density $V_d(1)\,2^{-d/2}/\sqrt{d+1}$. This gives
$\pi/\sqrt{12}$ for $d=2$ and $\pi/\sqrt{18}$ (FCC) for $d=3$.)}
\end{exercise}

\section{A better bound (non-constructive)}
\label{L1:sec:saturation}
\orig{\S1.3 in the notes}

\begin{definition}[Saturated packing]
\label{L1:def:saturated}
A packing $P=\bigcup_i K_i$ of congruent copies of $K$ in $\R^n$ is \emph{saturated} if
$P\cup\varphi K$ is not a packing for any $\varphi\in\R^n\rtimes\SO(n)$; that is,
no additional disk (ball) can be added.
\end{definition}

\begin{itemize}[nosep]
  \item Saturated packings certainly exist: \added{start from any packing and keep adding
        balls}. This is an ``infinite algorithm''. \added{It can be made rigorous with
        Zorn's lemma, or by adding balls greedily along an enumeration of a countable
        dense set of centres.}
  \item Saturation does not decrease $\dens^+(P)$ \added{since we only add balls}.
\end{itemize}
So for lower bounds on the best possible density we may assume that $P$ is saturated.

This is a surprisingly useful idea. Let $K$ be the unit ball, $K_i=\varphi_i K$ with centre $c_i$.

\begin{proposition}
\label{L1:prop:saturated-bound}
If $P$ is a saturated packing of unit balls in $\R^d$, then $\dens^+(P)\ge 2^{-d}$.
\end{proposition}

\begin{proof}
If $P$ is saturated, then for every $x\in\R^d$ there is an $i$ with $d(x,K_i)<1$.
\added{Otherwise $|x-c_i|\ge 2$ for all $i$, and the unit ball centred at $x$ could be added.}
Hence the doubled balls
\[
  2K_i\;:=\;\varphi_i\bigl(2\,\varphi_i^{-1}(K_i)\bigr)\;=\;\Ball(c_i,2)
\]
form a covering of $\R^d$. \added{Each $2K_i$ has $2^d$ times the volume of $K_i$, so the
fraction of $R\,\mathbb{D}^d$ covered by $\bigcup_i 2K_i$ is, up to a boundary term that
vanishes as $R\to\infty$, at most $2^d$ times the fraction covered by $P$. Since the
doubled balls cover everything,}
\[
  2^d\,\dens^+(P)\ \ge\ 1,\qquad\text{i.e.}\qquad \dens^+(P)\ \ge\ \frac1{2^d}. \qedhere
\]
\end{proof}
\orig{boxed correction: dilate about the centre, $(2\varphi_i^{-1}K)\varphi_i$}

\begin{remark}
\added{Compare with $(2\Z)^d$, of density $V_d(1)/2^d$. Since $V_d(1)<1$ exactly for
$d\ge13$, and $V_d(1)\to0$ superexponentially, the non-constructive bound $2^{-d}$ is
the better one in high dimensions.}
\end{remark}

We may go over some slight improvements in the future, but none seem satisfying.
\begin{remark}
\added{Known improvements of the $2^{-d}$ lower bound include the bound of
Minkowski and Hlawka
($2\zeta(d)\,2^{-d}$), Rogers and Ball ($c\,d\,2^{-d}$), Venkatesh 2013
($c\,d\log\log d\,2^{-d}$ along a sparse sequence of $d$), and, after the course,
Campos, Jenssen, Michelen and Sahasrabudhe 2023 ($c\,d\log d\,2^{-d}$) and
Klartag 2025 ($c\,d^2\,2^{-d}$). All of them are only polynomial improvements over $2^{-d}$.}
\end{remark}

The problem is that there is much more freedom in dimensions $d>3$!
\unclear{``$d>3$! even.'': perhaps ``even $d=4$''}
For example, take $(2\Z)^4$. The main diagonal of the cube $[0,2]^4$ has length
$2\sqrt4=4$, so the centre $(1,1,1,1)$ is at distance $2$ from every vertex.
\added{So a unit ball fits into every such hole. Adding them gives the
$D_4$ packing $(2\Z)^4\cup\bigl((2\Z)^4+(1,1,1,1)\bigr)$, which has twice the density,
$\pi^2/16$.}

\section{Upper bounds in dimension 2}
\label{L1:sec:upper2}
\orig{\S1.4}

The lower bound in $d=2$ is constructive: it comes from the hexagonal (triangular) lattice packing.

\begin{center}
\begin{tikzpicture}[scale=0.55]
  \clip (-0.6,-1.2) rectangle (10.6,6.4);
  \foreach \j in {-1,0,1,2,3,4} {
    \foreach \i in {-3,...,6} {
      \pgfmathsetmacro{\x}{2*\i+\j}
      \pgfmathsetmacro{\y}{sqrt(3)*\j}
      \draw[fill=blue!8] (\x,\y) circle (1);
      \fill (\x,\y) circle (1.5pt);
    }
  }
  \draw[very thick,red!70!black] (2,0) -- (4,0) -- (3,{sqrt(3)}) -- cycle;
  \draw[->,thick] (0,0) -- (2,0) node[midway,below] {$v_1$};
  \draw[->,thick] (0,0) -- (1,{sqrt(3)}) node[midway,left] {$v_2$};
\end{tikzpicture}
\\[2pt]
{\small Hexagonal packing of unit disks, basis $v_1=(2,0)$, $v_2=(1,\sqrt3)$ \added{(angle $\pi/3$)}, and an equilateral triangle of side $2$.}
\end{center}

\begin{example}[Hexagonal packing]
\label{L1:ex:hex}
Consider the equilateral triangle of side $2$ spanned by three mutually tangent unit disks.
Its height is $\sqrt{2^2-1^2}=\sqrt3$, so its area is $\sqrt3$. It contains three
sectors of angle $\pi/3$, i.e.\ half a unit disk, of area $\pi/2$. The plane is tiled by
such triangles, so
\[
  \dens_{\mathrm{hex}}\;=\;\frac{\pi/2}{\sqrt3}\;=\;\frac{\pi}{2\sqrt3}\;=\;\frac{\pi}{\sqrt{12}}\;\approx\;0.9069 .
\]
\end{example}

\subsection{Lattice packings}

For lattices, we want to maximize the minimum distance between two points of a unit lattice
\added{(a lattice $\Lat\subset\R^2$ of covolume $1$)}.
\orig{crossed out: ``minimize the det of the span for a unit lattice''}
This minimum distance can be arbitrarily small, but it is bounded above.

\begin{theorem}[\added{Lagrange}]
\label{L1:thm:lattice2}
Every lattice packing of congruent disks in $\R^2$ has density at most $\pi/\sqrt{12}$.
\added{Equality holds only for the hexagonal lattice.}
\end{theorem}

\begin{proof}
Scale so that $\Lat$ has covolume $1$, and let $2r$ be its minimal distance, so that disks of
radius $r$ centred at $\Lat$ form a packing of density $\pi r^2$. Let $v$ be a shortest
vector, $|v|=2r$. The lattice is the union of translates of the row $\Z v$, which lie on
parallel lines. Since the covolume is $1$, adjacent lines are at distance $h=1/(2r)$.
Note that this is the height of a fundamental parallelogram; the second basis vector
need not be orthogonal to $v$.

\begin{center}
\begin{tikzpicture}[scale=1.1]
  \draw (-0.5,0) -- (6.5,0);
  \draw (-0.5,1.9) -- (6.5,1.9);
  \foreach \x in {0,2,4,6} \fill (\x,0) circle (1.6pt);
  \fill (3,1.9) circle (1.6pt);
  \draw[dashed] (2,0) -- (3,1.9) -- (4,0);
  \draw[dotted] (3,0) -- (3,1.9);
  \draw[gray] (2,0) ++(0:2) arc (0:180:2);
  \node[below] at (3,0) {$2r$};
  \draw[<->] (6.3,0) -- (6.3,1.9) node[midway,right] {$h=\frac1{2r}$};
  \node[below] at (2,0) {\footnotesize $0$};
  \node[below] at (4,0) {\footnotesize $v$};
\end{tikzpicture}
\end{center}

\added{Translating along the row by multiples of $v$, we may take a lattice point $w$ in the
adjacent row whose orthogonal projection onto the first line lies within distance $r$ of
some lattice point there.} Since $|w-u|\ge 2r$ for all lattice points $u\ne w$,
\[
  h^2+r^2\;\ge\;(2r)^2,\qquad\text{so}\qquad \frac1{2r}\;=\;h\;\ge\;\frac{\sqrt3}{2}\,(2r)\;=\;\sqrt3\,r .
\]
Hence $r^2\le \frac1{2\sqrt3}$, i.e.\ $2r\le\sqrt{2/\sqrt3}$, and
\[
  \dens\;=\;\frac{\pi r^2}{1}\;\le\;\frac{\pi}{2\sqrt3}\;=\;\frac{\pi}{\sqrt{12}} .
\]
\added{Equality forces $h=\sqrt3 r$ with offset exactly $r$, which is the hexagonal lattice.}
\end{proof}
\orig{p.\,9 is an unnumbered duplicate of this computation}

\section{Partitioning space: Voronoi cells and Delaunay triangulations}
\label{L1:sec:voronoi}
\orig{\S1.5}

For packings that are not lattices, we need some analogous method to partition space.
Classically, Dirichlet--Voronoi diagrams work well.

\begin{definition}[Voronoi cell]
\label{L1:def:voronoi}
Given a packing with bodies $P_i$, the \emph{(Dirichlet--)Voronoi cell} associated with $P_i$ is
\[
  D_i\;=\;\{x\in\R^2 : d(x,P_i)<d(x,P_j)\ \text{for all } j\neq i\}.
\]
\end{definition}

When the $P_i$ are congruent disks and $d$ is the standard distance, this is the same as using
the distance to the centres of the disks. \added{Indeed $d(x,P_i)=|x-c_i|-1$ for $x\notin P_i$.}

Voronoi cells have nice properties:
\begin{itemize}[nosep]
  \item each cell is an intersection of half-spaces determined by the centres
        \added{($D_i=\bigcap_{j\ne i}\{x:|x-c_i|<|x-c_j|\}$)}, so it is convex;
  \item the cells partition space up to a boundary term \added{(a set of measure zero)}.
\end{itemize}

\begin{center}
\begin{tikzpicture}[scale=0.6]
  \clip (-3.2,-2.6) rectangle (5.2,4.3);
  \foreach \j in {-2,...,3} {
    \foreach \i in {-3,...,4} {
      \pgfmathsetmacro{\x}{2*\i+\j}
      \pgfmathsetmacro{\y}{sqrt(3)*\j}
      \draw[gray!60] (\x,\y) circle (1);
      \draw[blue!70!black,thick] (\x,\y) ++(30:1.1547) -- ++(150:1.1547) -- ++(210:1.1547)
         -- ++(270:1.1547) -- ++(330:1.1547) -- ++(30:1.1547) -- cycle;
      \draw[red!70!black,densely dashed] (\x,\y) -- ++(2,0);
      \draw[red!70!black,densely dashed] (\x,\y) -- ++(1,{sqrt(3)});
      \draw[red!70!black,densely dashed] (\x,\y) -- ++(-1,{sqrt(3)});
      \fill (\x,\y) circle (1.3pt);
    }
  }
\end{tikzpicture}
\\[2pt]
{\small \added{Hexagonal packing: Voronoi cells (solid hexagons) and Delaunay triangles (dashed).}}
\end{center}

The dual notion is the \emph{Delaunay triangulation}.
\begin{itemize}[nosep]
  \item It is not unique in general (e.g.\ four points at the vertices of a square).
  \item It is harder to define directly. Its characterization: the circumcircles of the
        Delaunay triangles are empty of other points (vertices).
  \item The circumcentre of a Delaunay triangle is a Voronoi vertex.
\end{itemize}
\begin{addedblock}
Precisely: for a locally finite set $X\subset\R^2$ (here, the centres), a triangle with vertices
in $X$ is \emph{Delaunay} if its open circumdisk contains no point of $X$. A Delaunay triangulation is
a triangulation of (the convex hull of) $X$ by Delaunay triangles. It is unique when no four
points of $X$ are concyclic.
\end{addedblock}

\subsection{Existence via lifting}

Consider a collection of points in the plane and lift them to the paraboloid
$z=x^2+y^2$.

\begin{lemma}
\label{L1:lem:lifting}
Let $a,b,c,d$ be points in the $(x,y)$-plane, and lift them to
$\hat a,\hat b,\hat c,\hat d$ on $z=x^2+y^2$ \added{(i.e.\ $\hat p=(p,|p|^2)$)}. Then $d$ lies
inside the circumcircle of $a,b,c$ if and only if $\hat d$ lies in the lower half-space of the
plane $\langle\hat a,\hat b,\hat c\rangle$ \added{spanned by $\hat a,\hat b,\hat c$}.
\end{lemma}

\begin{proof}
A plane tangent to the paraboloid at the point above $(p,q)$ has the form
\[
  z\;=\;2px+2qy-(p^2+q^2).
\]
Shifting it up by $h^2$ gives
\[
  z\;=\;2px+2qy-(p^2+q^2)+h^2 .
\]
\fixed{signs: notes have $+(p^2+q^2)$ and $(x+p)^2$, $(y+q)^2$}
It meets the paraboloid where
\[
  x^2+y^2\;=\;2px+2qy-(p^2+q^2)+h^2
  \quad\Longleftrightarrow\quad (x-p)^2+(y-q)^2\;=\;h^2 .
\]
So the plane through $(p,q,p^2+q^2)$, parallel to the tangent plane there, at height $h^2$
above it, projects its intersection with the paraboloid onto the circle of radius $h$
about $(p,q)$ in the $(x,y)$-plane.
\begin{addedblock}
Every non-vertical plane $z=\alpha x+\beta y+\gamma$ meeting the paraboloid in more than one point
is of this form, with $p=\alpha/2$, $q=\beta/2$ and $h^2=\gamma+p^2+q^2>0$. Apply this to the plane
through $\hat a,\hat b,\hat c$. The corresponding circle passes through $a,b,c$, so it is their
circumcircle. A point $\hat x$ of the paraboloid lies strictly below this plane iff
$x^2+y^2<2px+2qy-(p^2+q^2)+h^2$, i.e.\ iff $(x-p)^2+(y-q)^2<h^2$, i.e.\ iff $x$ lies strictly
inside the circumcircle.
\end{addedblock}
\end{proof}

\begin{corollary}[Existence of Delaunay triangulations]
\label{L1:cor:delaunay-exists}
The existence of the convex hull of the lifted points implies the existence of a Delaunay triangulation.
\added{More precisely: project the faces of the lower convex hull of $\{\hat x:x\in X\}$ to the plane.
By Lemma~\ref{L1:lem:lifting}, a lower face has no lifted point below its plane, so its projection has
an empty circumcircle. Triangulating any non-triangular faces (which come from concyclic points)
gives a Delaunay triangulation. For the infinite, locally finite point sets of a packing, apply this
to large finite pieces; the triangles near a fixed point stabilize.}
\end{corollary}

\section{Thue's theorem}
\label{L1:sec:thue}
\orig{\S1.6; a first ``1.6 Thm'' on p.\,13 is abandoned}

Let $P$ be a saturated packing of unit disks in $\R^2$ and $DT$ a Delaunay triangulation of its set of
centres. Every edge of $DT$ has length at least $2$, since the disks do not overlap.
\added{Moreover, every triangle of $DT$ has circumradius $R<2$. Its open circumdisk contains no
centres, so the circumcentre is at distance $\ge R$ from every centre. If $R\ge2$, a unit disk centred
there could be added, contradicting saturation.}

\begin{lemma}
\label{L1:lem:angle}
The largest angle $\theta$ of a triangle $ABC\in DT$ of a saturated packing satisfies
\[
  \frac\pi3\;\le\;\theta\;<\;\frac{2\pi}3 .
\]
\end{lemma}
\unclear{upper inequality written ambiguously; the proof gives strict ``$<$''}

\begin{proof}
The lower bound is obvious, since $\theta$ is the largest angle of a triangle.
Assume $\theta\ge 2\pi/3$. Let $A$ be the smallest angle. Then $A\le\frac12(\pi-\theta)\le\pi/6$,
so $\sin A\le\frac12$. By the circumradius formula, and since the side $BC$ opposite $A$ has
length at least $2$,
\[
  R\;=\;\frac12\,\frac{|BC|}{\sin A}\;\ge\;\frac12\cdot\frac{2}{1/2}\;=\;2 .
\]
\fixed{notes: ``$=4$''; $\frac12\cdot\frac{2}{1/2}=2$}
\added{This contradicts $R<2$.}
\end{proof}

\begin{lemma}
\label{L1:lem:triangle-density}
The density in a single Delaunay triangle is at most $\pi/\sqrt{12}$, with equality
if and only if the triangle is equilateral \added{(of side $2$)}.
\end{lemma}
\begin{addedblock}
Here the \emph{density} of a triangle $T\in DT$ means $\frac{\pi/2}{\operatorname{area}(T)}$:
the three sectors of the unit disks at the vertices of $T$ have angles summing to $\pi$, hence total area $\pi/2$.
\end{addedblock}

\begin{proof}
Let $B$ be the largest angle of $ABC$. Then
\[
  \operatorname{area}(ABC)\;=\;\tfrac12\,|AB|\,|CB|\,\sin B
  \;\ge\;\tfrac12\cdot2\cdot2\cdot\min_{B\in[\pi/3,\,2\pi/3)}\sin B
  \;=\;\tfrac12\cdot2\cdot2\cdot\tfrac{\sqrt3}{2}\;=\;\sqrt3,
\]
using Lemma~\ref{L1:lem:angle}, and the minimum is attained at $B=\pi/3$. So $\operatorname{area}(T)\ge\sqrt3$
for every $T\in DT$, and the density of $T$ is at most $\frac{\pi/2}{\sqrt3}=\frac{\pi}{2\sqrt3}=\frac{\pi}{\sqrt{12}}$.
\added{Equality requires $|AB|=|CB|=2$ and $B=\pi/3$, i.e.\ an equilateral triangle of side $2$.}
\end{proof}

\begin{theorem}[\added{Thue; Fejes T\'oth}]
\label{L1:thm:thue}
Every packing $P$ of unit disks in the plane satisfies $\dens^+(P)\le\pi/\sqrt{12}$.
\end{theorem}

\begin{proof}
\added{Saturating $P$ does not decrease $\dens^+$ (\S\ref{L1:sec:saturation}), so we may assume $P$ is saturated.}
For a finite union of Delaunay triangles, the density is
\[
  \sum_{T\in DT}\frac{\operatorname{area}(T)}{\operatorname{area}(DT)}\cdot\dens(T)
  \;\le\;\frac{\pi}{\sqrt{12}},
\]
a weighted average of the densities of the triangles, each at most $\pi/\sqrt{12}$
by Lemma~\ref{L1:lem:triangle-density}. Hence any finite union of Delaunay triangles has density
at most $\pi/\sqrt{12}$.

For a large disk $R\,\mathbb{D}^2$, most of its area is covered by the Delaunay triangles inside it,
so the error in the limit is a boundary term.
\begin{addedblock}
Every Delaunay triangle has circumradius $<2$, so diameter $<4$. Let $U_R$ be the union of the
triangles meeting $R\,\mathbb{D}^2$. Then $R\,\mathbb{D}^2\subset U_R\subset(R+4)\mathbb{D}^2$. The sectors
in the triangles of $U_R$ account for every disk centred in $R\,\mathbb{D}^2$. The disks meeting
$R\,\mathbb{D}^2$ but centred outside it lie in the annulus $(R+1)\mathbb{D}^2\setminus(R-1)\mathbb{D}^2$, of area $O(R)$. So
\[
  \vol(P\cap R\,\mathbb{D}^2)\;\le\;\frac{\pi}{\sqrt{12}}\,\vol(U_R)+O(R)
  \;\le\;\frac{\pi}{\sqrt{12}}\,\pi(R+4)^2+O(R).
\]
Dividing by $\pi R^2$ and letting $R\to\infty$ gives $\dens^+(P)\le\pi/\sqrt{12}$.
\end{addedblock}
\end{proof}
\orig{p.\,16 sketch: large concentric disks}

\begin{remark}
\added{Combined with Example~\ref{L1:ex:hex}, this shows that the hexagonal packing is the
densest packing of congruent disks in the plane. The statement goes back to Thue (1892, 1910).
The first proof generally accepted as complete is due to L.~Fejes T\'oth (1940); the argument above
follows his Delaunay-triangle approach. For lattice packings the result is Lagrange's
(Theorem~\ref{L1:thm:lattice2}).}
\end{remark}

% MAT670 -- Lecture 2 (transcribed from MAT670_2, handwritten pp. 17--31)
\chapter{Lecture 2: The Fejes T\'oth inequality and Blichfeldt's bound}\label{L2:chap}

\section{Last time: Thue's theorem in the plane}\label{L2:sec:thue}
\orig{Notes pp.~17--31 (MAT670\_2).}
Last time we gave a modern take on Thue's result in $\R^2$:

\begin{theorem}[Thue]\label{L2:thm:thue}
The density of a packing of the plane by congruent circles is bounded by the density of the regular
hexagonal packing\added{, namely $\pi/\sqrt{12}\approx 0.9069$}.
\end{theorem}

\begin{center}
\begin{tikzpicture}[scale=0.9]
  \draw[thick] (0,0) circle (0.5);
  \foreach \k in {0,...,5} { \draw[thick] ({\k*60}:1) circle (0.5); }
\end{tikzpicture}
\end{center}

The idea of the proof was to use \emph{saturation} and the \emph{Delaunay triangulation}.

\begin{remark}
In the proof we left out the error analysis of the limits, etc.
\end{remark}

\begin{remark}
Rogers showed a density bound coming from the regular simplex.
\added{Precisely (C.~A.~Rogers, 1958): the density of any packing of unit balls in $\R^n$ is at most the
fraction of a regular simplex of edge length $2$ that is covered by the $n+1$ unit balls centred at its
vertices. For $n=2$ this is exactly $\pi/\sqrt{12}$, so Rogers' bound generalizes Thue's theorem.}
\end{remark}

% ---------------------------------------------------------------------------
\section{The Fejes T\'oth inequality for \texorpdfstring{$\Sph^2$}{S2}}\label{L2:sec:FT}

All distances on $\Sph^2$ are spherical (angular) distances.

\begin{theorem}[Fejes T\'oth\added{, 1943}]\label{L2:thm:FT}
For any $n\ge 3$ points on $\Sph^2$ there exists a pair with spherical distance
\begin{equation}\label{L2:eq:FT}
d \;\le\; \arccos\!\left(\frac{\cot^2(\omega)-1}{2}\right),
\qquad \omega=\frac{n}{n-2}\cdot\frac{\pi}{6}. \tag{$*$}
\end{equation}
\end{theorem}
\fixed{Notes have $\cot(\omega)$; must be $\cot^2(\omega)$.}

\begin{addedblock}
Check: for $n=12$, $\omega=\pi/5$, and $(\cot^2(\pi/5)-1)/2 = 1/\sqrt5$, so
$d\le\arccos(1/\sqrt5)\approx 63.43^\circ$, the edge of the regular icosahedron. Similarly $n=3,4,6$ give
$120^\circ$, $\arccos(-1/3)\approx109.47^\circ$ and $90^\circ$ (equilateral triangle on a great circle,
regular tetrahedron, regular octahedron). These are exactly the cases in which the bound is attained.
\end{addedblock}

\begin{remark}\label{L2:rem:capdensity}
The inequality can be rewritten as a density result for spherical caps of diameter $d$ (i.e.\ angular
radius $d/2$): a cap of angular radius $r$ has area $2\pi(1-\cos r)$, so the density of $n$ non-overlapping
caps of diameter $d$ is
\[
\frac{n\cdot 2\pi\bigl(1-\cos(d/2)\bigr)}{4\pi}
\added{\ \le\ \frac n2\Bigl(1-\cos\tfrac{d_n}{2}\Bigr)}
\;\longrightarrow\; \frac{\pi}{\sqrt{12}}
\added{\quad (n\to\infty),}
\]
\added{where $d_n$ is the right-hand side of \eqref{L2:eq:FT}.} So in the limit we recover Thue's bound.
\end{remark}

\begin{exercise}
Verify the limit in Remark~\ref{L2:rem:capdensity}. \added{(Hint: as $n\to\infty$, $d_n\to 0$ and
$\frac{\sqrt3}{4}d_n^2\sim \frac{4\pi}{2n}$; see also Exercise~\ref{L2:ex:Lhuilier}.)}
\end{exercise}

The ``flat'' values also converge, but it is not so nice to take the limit.
\unclear{``Flat values'': presumably the planar triangles of the convex hull instead of spherical ones.}

A natural question: is there a good \emph{lower} bound, i.e.\ can we construct a sequence of sets of
points that achieves this bound?
\begin{addedblock}
Equality in \eqref{L2:eq:FT} holds only for $n=3,4,6,12$. Asymptotically the bound is sharp: there are
configurations (e.g.\ suitable subdivisions of the icosahedron) whose cap density tends to
$\pi/\sqrt{12}$.
\end{addedblock}

\subsection{A lemma on spherical triangles}

To prove the Fejes T\'oth theorem we need the following lemma.

\begin{lemma}\label{L2:lem:triangle}
Let $ABC$ be a spherical triangle and let $\overline{AB}$ be its shortest side. Let $ABC'$ be the
equilateral spherical triangle drawn on $\overline{AB}$\added{, with $C'$ on the same side of the great
circle $AB$ as $C$}. If
\[
\operatorname{Area}(ABC) < \operatorname{Area}(ABC'),
\]
then the spherical radius of the circumcircle of $ABC$ (the circle in which the plane through $A,B,C$
meets $\Sph^2$) is greater than the length of $\overline{AB}$.
\end{lemma}

\begin{proof}
Write $a$ for the length of $\overline{AB}$.

\emph{Step 1: draw the picture.} Let $C'$ be such that $ABC'$ is equilateral, with $C'$ on the same side
of the great circle $AB$ as $C$. Let $\alpha,\beta,\gamma$ be the circles of radius $a$ centred at
$A,B,C'$ respectively\added{; thus $\alpha$ passes through $B$ and $C'$, $\beta$ through $A$ and $C'$,
and $\gamma$ through $A$ and $B$}. Let $A'$ be the second intersection point of $\alpha$ and $\gamma$
\added{(so that $AC'A'$ is equilateral)}, and $B'$ the second intersection point of $\beta$ and $\gamma$.
Let $-A,-B$ denote the antipodes of $A,B$.

\begin{figure}[ht]
\centering
\begin{tikzpicture}[scale=0.95]
  \begin{scope}
    \clip (-4.3,-2.5) rectangle (4.3,5.3);
    % admissible region for C: C' side of AB, below the Lexell circle, outside alpha and beta
    \fill[gray!18] (0,1.470) circle (3.522);
    \fill[white] (0,2.885) circle (2.885);
    \fill[white] (-2.683,-3.482) circle (4.396);
    \fill[white] (2.683,-3.482) circle (4.396);
    % circles
    \draw[thick] (0,1.470) circle (3.522);                 % great circle AB
    \draw[thick,dashed] (0,2.885) circle (2.885);          % Lexell circle
    \draw (-2.683,-3.482) circle (4.396);                  % alpha
    \draw (2.683,-3.482) circle (4.396);                   % beta
    \draw (0,0) circle (2.241);                            % gamma
  \end{scope}
  % equilateral triangle ABC' (great-circle arcs)
  \draw[red!70!black,densely dotted,thick] (-1.368,-1.775) -- (-1.122,-1.868) -- (-0.875,-1.941) -- (-0.626,-1.995) -- (-0.376,-2.031) -- (-0.125,-2.049) -- (0.125,-2.049) -- (0.376,-2.031) -- (0.626,-1.995) -- (0.875,-1.941) -- (1.122,-1.868) -- (1.368,-1.775)
     -- (1.212,-1.573) -- (1.067,-1.384) -- (0.930,-1.207) -- (0.800,-1.038) -- (0.676,-0.877) -- (0.557,-0.722) -- (0.441,-0.572) -- (0.328,-0.426) -- (0.218,-0.283) -- (0.109,-0.141) -- (0,0)
     -- (-0.109,-0.141) -- (-0.218,-0.283) -- (-0.328,-0.426) -- (-0.441,-0.572) -- (-0.557,-0.722) -- (-0.676,-0.877) -- (-0.800,-1.038) -- (-0.930,-1.207) -- (-1.067,-1.384) -- (-1.212,-1.573) -- cycle;
  % points
  \foreach \p/\n/\pos in {(-1.368,-1.775)/A/below left, (1.368,-1.775)/B/below right, (0,0)/C'/above,
      (-2.065,0.870)/A'/left, (2.065,0.870)/B'/right, (2.790,3.620)/-A/right, (-2.790,3.620)/-B/left}
      { \fill \p circle (1.6pt) node[\pos,font=\small] {$\n$}; }
  \fill (2.7,1.2) circle (1.6pt) node[right,font=\small] {$C$};
  % labels
  \node[font=\small] at (-4.05,1.0) {$\alpha$};
  \node[font=\small] at (4.05,1.0) {$\beta$};
  \node[font=\small] at (-1.75,1.65) {$\gamma$};
  \node[font=\small,align=center] at (0,4.25) {Lexell\\circle};
  \draw[thin,gray] (3.47,2.1) -- (3.95,2.4) node[right,font=\small,black] {great circle $AB$};
\end{tikzpicture}
\caption{The configuration of Lemma~\ref{L2:lem:triangle} (side $a=70^\circ$), drawn in stereographic
projection centred at $C'$, so circles on $\Sph^2$ appear as circles. Shaded: the points between the great circle $AB$ and the Lexell circle that lie outside $\alpha$ and
$\beta$ (the possible positions of $C$); they avoid the disc bounded by $\gamma$.}
\label{L2:fig:lexell}
\end{figure}
\orig{Hand sketch p.~19 has the great circle $AB$ as a straight line; redrawn exactly here.}

\emph{Step 2: show the picture is ``correct''.} Note that the spherical triangles $ABA'$, $ABB'$ and $ABC'$
all have equal area. \added{(For $ABA'$: it is isosceles with legs $a$ and apex angle $2\theta$, where
$\theta$ is the angle of $ABC'$; a direct computation with the formula
$\tan\frac E2=\frac{\sin\varphi\,\tan\frac b2\tan\frac c2}{1+\cos\varphi\,\tan\frac b2\tan\frac c2}$ for the
area $E$ of a triangle with sides $b,c$ and included angle $\varphi$, together with
$\tan^2\frac a2=1-2\cos\theta$, shows that both areas agree.)} Hence the circle in which the plane through
$\{A',B',C'\}$ meets the sphere is the locus of points $X$ with
$\operatorname{Area}(ABX)=\operatorname{Area}(ABC')$ (on that side of $AB$): this is the
\emph{Lexell circle theorem}. \added{(Lexell: for fixed $A,B$, the locus of apexes $X$ of triangles $ABX$
of fixed area is an arc of a small circle through $-A$ and $-B$.)}

By assumption, $\operatorname{Area}(ABC)<\operatorname{Area}(ABC')$, so $C$ lies between the Lexell circle
and the great circle $AB$. As $\overline{AB}$ is the shortest edge of $ABC$, $C$ does not lie in
\added{the open discs bounded by} $\alpha$ or $\beta$. It follows \added{(see
Figure~\ref{L2:fig:lexell})} that $C$ does not lie in \added{the disc bounded by} $\gamma$.
\added{(The part of the disc of $\gamma$ on the $C'$ side of $AB$ and below the Lexell circle is covered by
the discs of $\alpha$ and $\beta$.)}
\unclear{This covering step is read off the picture in the notes.}

But $C$ is on the $C'$ side of $AB$. \added{The circles through $A$ and $B$ form a pencil whose centres lie
on the perpendicular bisector of $AB$; on the $C'$ side, the region enclosed by such a circle grows with its
radius. Since $\gamma$ is the member of the pencil centred at $C'$ and $C$ lies outside it,} the
circumcircle of $ABC$ has radius greater than the radius of $\gamma$, which is the length of
$\overline{AB}$.
\end{proof}

\subsection{Proof of the Fejes T\'oth theorem}

The case $n=3$ is trivial\added{: the three sides of a spherical triangle have total length at most
$2\pi$, so one of them is at most $2\pi/3=\arccos(-1/2)$}. So let $n\ge 4$ and let $p_1,\dots,p_n\in\Sph^2$.

Without loss of generality the convex hull $\conv\{p_1,\dots,p_n\}$ contains the origin $0$
\added{in its interior}. Otherwise we could flow the points from the pole towards the equator
\added{of a hemisphere containing them} and increase all distances.
\begin{addedblock}
If all points lie in an open hemisphere with pole $N$, replace the colatitude $\phi$ (angle from $N$) of
each point by $\lambda\phi$, $\lambda>1$ close to $1$. From
$\cos d=\cos\phi_1\cos\phi_2+\sin\phi_1\sin\phi_2\cos\Delta$ one checks
$\frac{\partial}{\partial\lambda}\cos d\le-(\phi_1-\phi_2)\sin(\phi_1-\phi_2)\le0$, with strict
inequality for distinct points, so all mutual distances increase.
\end{addedblock}
\unclear{Degenerate case (points in a closed, not open, hemisphere) needs a small perturbation argument.}

\emph{Step 1: triangulate the convex hull.}
\added{The convex hull is a polytope whose vertices are among the $p_i$; after a small perturbation (which
changes distances only slightly) we may assume it is simplicial, with all $n$ points as vertices.}
By the Euler characteristic, it has $2n-4$ faces.

\begin{exercise}\label{L2:ex:euler}
Show that a triangulated convex polytope with $n$ vertices has $2n-4$ faces.
\emph{Solution sketch:} $v-e+f=2$ and $3f=2e$ for triangulations; hence $2n-2e+2f=4$, so
$2n-3f+2f=4$, i.e.\ $f=2n-4$.
\end{exercise}

The triangulation of $\conv\{p_1,\dots,p_n\}$ gives a \emph{spherical net} by radial projection to the
sphere: the edges become geodesic arcs\added{, and since $0$ is in the interior of the hull the $2n-4$
spherical triangles tile $\Sph^2$}.

\begin{center}
\begin{tikzpicture}[scale=1.1]
  % tetrahedron
  \coordinate (P1) at (-1.0,-0.45); \coordinate (P2) at (0.45,-0.6);
  \coordinate (P3) at (0.85,-0.15); \coordinate (P4) at (0.05,0.75);
  \draw[thick] (P1)--(P2)--(P3)--(P4)--cycle (P4)--(P2);
  \draw[dashed] (P1)--(P3);
  \draw[->] (1.4,0.05) -- (2.3,0.05);
  % spherical net
  \begin{scope}[xshift=3.6cm]
    \draw (0,0) circle (1);
    \draw[thick] (-0.83,-0.38) to[bend right=18] (0.40,-0.80);
    \draw[thick] (0.40,-0.80) to[bend right=18] (0.83,-0.28);
    \draw[thick] (0.83,-0.28) to[bend right=22] (0.05,0.96);
    \draw[thick] (0.05,0.96) to[bend right=22] (-0.83,-0.38);
    \draw[thick] (0.05,0.96) to[bend left=10] (0.40,-0.80);
    \draw[dashed] (-0.83,-0.38) to[bend left=12] (0.83,-0.28);
    \foreach \p in {(-0.83,-0.38),(0.40,-0.80),(0.83,-0.28),(0.05,0.96)} \fill \p circle (1.3pt);
  \end{scope}
\end{tikzpicture}
\end{center}

\emph{Step 2.} Now, for $\{p_1,\dots,p_n\}$, suppose that
\[
\operatorname{length}\bigl(\widehat{p_ip_j}\bigr) > d \qquad\text{for all } i\ne j,
\]
where $d=d_n$ is the right-hand side of \eqref{L2:eq:FT}.

\begin{exercise}[L'Huilier]\label{L2:ex:Lhuilier}
Show that $d$ in \eqref{L2:eq:FT} is the side length of an equilateral spherical triangle of area
$\dfrac{4\pi}{2n-4}$. \added{(One can use L'Huilier's formula, the spherical analogue of Heron's formula.
Alternatively: the angle $\theta$ of such a triangle satisfies $3\theta-\pi=\frac{2\pi}{n-2}$, i.e.\
$\theta=2\omega$, and the spherical law of cosines for angles gives
$\cos d=\frac{\cos\theta}{1-\cos\theta}=\frac{\cot^2\omega-1}{2}$.)}
\end{exercise}

Then the triangle $p_ip_jp_k$ of minimal area in the spherical net satisfies the hypotheses of the Lemma:
its area is small, but its edges are all longer than $d$. Here ``small'' means: less than or equal to
the area of the equilateral triangle of side length $d$\added{, since $2n-4$ triangles of total area $4\pi$
tile the sphere. As the area of an equilateral triangle increases with its side, and the shortest side of
$p_ip_jp_k$ is longer than $d$, the hypothesis of Lemma~\ref{L2:lem:triangle} holds with strict
inequality.}

Hence the circumcircle of $p_ip_jp_k$ has radius larger than $d$. But the net construction implies that
the circumcircle is empty\added{: the face $p_ip_jp_k$ lies in a supporting plane of the convex hull, so no
$p_l$ lies in the open cap cut off by that plane (and this cap is smaller than a hemisphere, since $0$ is on
the other side)}.

So we may place a new point at the centre \added{of this cap; it has distance greater than $d$ from every
$p_l$}. This new collection of $n+1$ points satisfies the same inequalities for the \emph{same}~$d$!
\added{Moreover the new hull still contains $0$ in its interior and has $2(n+1)-4$ faces, so its smallest
triangle again has area less than that of the equilateral triangle of side $d$, and we can repeat
indefinitely.}

This is absurd, e.g.\ by a volume bound: $\vol(\Sph^2)=4\pi$, and \added{the caps of radius $d/2$ around the
points are pairwise disjoint, so} we would have
\[
\sum_{i=1}^{N}\operatorname{area}\bigl(\operatorname{cap}(d/2)\bigr) < 4\pi \qquad\text{for all } N>0,
\]
\added{which is impossible since each cap has positive area.} \qed

% ---------------------------------------------------------------------------
\section{Higher-dimensional sphere packing: Blichfeldt's method}\label{L2:sec:blichfeldt}

Consider an early method of Blichfeldt (1929) to get density bounds in higher dimensions.
\added{(H.~F.~Blichfeldt, \emph{The minimum value of quadratic forms, and the closest packing of spheres},
Math.\ Ann.\ 101 (1929).)}

\begin{definition}\label{L2:def:convexbody}
A \emph{convex body} $C$ in $\R^n$ is a compact convex subset of $\R^n$ with non-empty interior.
\end{definition}
\orig{Crossed out: ``(so think $C=\mathbb{B}^n$)''.}

Also consider collections of isometries $\{\varphi_i\}_{i=1}^\infty$ of $\R^n$ such that
$\{\varphi_iC\}_{i=1}^\infty$ is a packing \added{(i.e.\ the interiors of the $\varphi_iC$ are pairwise
disjoint)}.

For now, $C=\Ball^n$ \added{is the closed unit ball}, and a packing is of the form
$\{\varphi_i\Ball^n\}_{i=1}^\infty$.

Consider replacing the ball with its characteristic function
\[
\chi(x)=\begin{cases}1, & |x|\le1,\\ 0, & |x|>1.\end{cases}
\]
So it has mass
\[
I(\chi)=\int_{\R^n}\chi(x)\,dx \;\added{=\vol(\Ball^n)},
\]
and a packing corresponds to the function
\[
x\longmapsto\sum_{i=1}^\infty\chi(\varphi_i^{-1}x)
\]
\added{(the characteristic function of $\bigcup_i\varphi_i\Ball^n$, up to the measure-zero set of boundary
points).} For a general convex body $C$ this works as well.

Furthermore, we can now replace $\chi$ with some other function $f$, and in general associate to a packing
the function
\[
x\longmapsto\sum_{i=1}^\infty f(\varphi_i^{-1}x).
\]
A density may be defined similarly to the packing density, where each object now has mass
\[
I(f)=\int_{\R^n}f(x)\,dx
\]
spread over $\operatorname{supp}(f)$. So:
\[
\Delta_f \;=\; \dens_{\text{packing}}\cdot\frac{I(f)}{\vol(C)}.
\]
\added{Here $\dens_{\text{packing}}$ is the density of the packing $\{\varphi_iC\}$ (Lecture~1), i.e.\ the
average number of copies per unit volume times $\vol(C)$; $\Delta_f$ is the average of
$\sum_i f(\varphi_i^{-1}x)$ over $x\in\R^n$.}

We can think of this as spreading out and renormalizing the characteristic function in some (nice,
hopefully) way: ``cut up the sphere and spread it out.''

\textbf{Insight.} There are functions $f$ for which the sums $\sum_{i=1}^\infty f(\varphi_i^{-1}x)$ are
uniformly pointwise bounded over $\R^n$ \emph{and} over all collections of isometries $\{\varphi_i\}$ such
that $\{\varphi_iC\}$ is a packing!

% ---------------------------------------------------------------------------
\section{Blichfeldt gauges}\label{L2:sec:gauge}

\begin{definition}\label{L2:def:gauge}
Given a convex body $C$ in $\R^n$, a \added{non-negative, integrable} function $f$ on $\R^n$ is a
\emph{Blichfeldt gauge} for $C$ if for any collection $\{\varphi_i\}_{i=1}^\infty$ of isometries of
$\R^n$ such that $\{\varphi_iC\}_{i=1}^\infty$ is a packing,
\[
\sum_{i=1}^\infty f(\varphi_i^{-1}x)\le 1 \qquad\text{for all } x\in\R^n.
\]
\end{definition}
\fixed{Notes: ``isometries of $\R^d$''; should be $\R^n$.}
\orig{Margin: ``add non-negative condition''.}

\begin{lemma}\label{L2:lem:gauge}
If $f$ is a Blichfeldt gauge for $C$, then the density of any packing of $\R^n$ by copies of $C$ satisfies
\[
\dens(C)\;\le\;\frac{\vol(C)}{I(f)}. \tag{$**$}
\]
\end{lemma}
\fixed{Notes have strict $<$; equality can occur (e.g.\ $f=\chi_{\interior C}$, $C$ a cube).}

\begin{proof}[Proof sketch]
By the definition, $\Delta_f\le 1$ \added{(the average of a function bounded by $1$ is at most $1$)}.
Since $\Delta_f=\dens\cdot I(f)/\vol(C)$, this gives $(**)$. \added{The limiting argument that makes this
precise is carried out for balls in Section~\ref{L2:sec:bound}; the same argument works for any $C$ and any
gauge with bounded support.}
\end{proof}

% ---------------------------------------------------------------------------
\section{Back to \texorpdfstring{$C=\Ball^n$}{C = B\^{}n}}\label{L2:sec:radial}

We can consider radially symmetric functions by an averaging argument: there is no benefit to anisotropy,
since $\SO(n)$ kills it.
\begin{addedblock}
If $f$ is a gauge for $\Ball^n$, so is $\bar f(x)=\int_{\SO(n)}f(gx)\,dg$ (Haar probability measure):
for each $g$, $\{\varphi_ig^{-1}\Ball^n\}=\{\varphi_i\Ball^n\}$ is the same packing, so
$\sum_i f(g\varphi_i^{-1}x)\le1$, and averaging over $g$ gives $\sum_i\bar f(\varphi_i^{-1}x)\le 1$.
Moreover $I(\bar f)=I(f)$, so $\bar f$ gives the same bound.
\end{addedblock}
So it suffices to consider functions of the distance $|x|$! Rankin does a complicated study to get a good
function \added{(R.~A.~Rankin, \emph{On the closest packing of spheres in $n$ dimensions}, Ann.\ of Math.\
48 (1947))}. We consider
\[
f_0(x)=\begin{cases}1-\tfrac12|x|^2, & |x|\le\sqrt2,\\[2pt] 0, & |x|>\sqrt2.\end{cases}
\]

We can show the following for a packing of unit radius balls.

\begin{lemma}\label{L2:lem:f0}
$f_0$ is a Blichfeldt gauge for $\Ball^n$.
\end{lemma}

\begin{proof}
Consider $m$ unit balls \added{in a packing} with centres $x_1,\dots,x_m\in\R^n$. Then clearly, for all
pairs,
\[
|x_i-x_j|^2\ge 4. \tag{$\dagger$}
\]
Summing $(\dagger)$ over $1\le i<j\le m$ (the cross terms cancel\added{, via the identity
$\sum_{i<j}|x_i-x_j|^2=m\sum_i|x_i|^2-\bigl|\sum_ix_i\bigr|^2$}), we get
\[
m\sum_{i=1}^m|x_i|^2-\Bigl|\sum_{i=1}^m x_i\Bigr|^2\;\ge\;2m(m-1),
\]
and therefore
\[
\sum_{i=1}^m|x_i|^2\;\ge\;2(m-1).
\]
\orig{Notes write the centres in coordinates $(a_i,b_i,\dots,n_i)$.}
The origin here is an arbitrary point. So, given $x\in\R^n$, consider the collection of the $m$ balls whose
centres lie within $\sqrt2$ of $x$ \added{(the only ones contributing to $\sum_if_0(x-x_i)$)}, with masses
defined by $f_0$ centred at the $x_i$. Writing $r_i=|x-x_i|$, we see that
\[
\sum_{i=1}^m\frac{2-r_i^2}{2}\;\le\;\frac{2m-2(m-1)}{2}\;=\;1,
\]
and so $f_0$ is a Blichfeldt gauge.
\end{proof}

% ---------------------------------------------------------------------------
\section{Blichfeldt's bound}\label{L2:sec:bound}

So we can integrate this, using spherical shells:
\[
I(f_0)=\int_0^{\sqrt2}\frac{2-r^2}{2}\,d\bigl(\vol(\Ball^n)\,r^n\bigr)
=\frac{2^{\frac n2+1}}{n+2}\,\vol(\Ball^n).
\]
\fixed{Notes first had $\frac{4\cdot2^{n/2}}{n+2}$; corrected in red to $\frac{2^{(n+2)/2}}{n+2}$.}

\begin{exercise}\label{L2:ex:integral}
Verify this integral using spherical shells. \added{($I(f_0)=n\vol(\Ball^n)\int_0^{\sqrt2}
(1-\tfrac{r^2}{2})r^{n-1}dr=\vol(\Ball^n)\,2^{n/2}\bigl(1-\tfrac{n}{n+2}\bigr)$.)}
\end{exercise}

\subsection*{Limit argument for a box}

Take a large box of side length $t$ with $K$ balls in it.
\orig{Notes write $|I|$ for $K$.}
\added{Their centres lie in the concentric box of side $t-2$, so the supports of the gauges
$f_0(x-x_i)$ lie in the concentric box of side} $t+2\sqrt2-2$. Since $f_0\ge 0$ and
$\sum_if_0(x-x_i)\le1$ everywhere, integrating over this larger box gives
\[
(t+2\sqrt2-2)^n\;\ge\;K\int_{\R^n}f_0(x)\,dx\;=\;\frac{K\,\vol(\Ball^n)\,2^{\frac n2+1}}{n+2}.
\]
Hence the density of these $K$ balls in the box $[0,t]^n$ satisfies
\[
\frac{K\vol(\Ball^n)}{t^n}\;\le\;\frac{n+2}{2^{\frac n2+1}}\Bigl(1+\frac{2\sqrt2-2}{t}\Bigr)^n,
\]
and letting $t\to\infty$:

\begin{theorem}[Blichfeldt's bound]\label{L2:thm:blichfeldt}
For any packing of $\R^n$ by unit balls,
\[
\dens\;\le\;\frac{n+2}{2^{\frac n2+1}} \;=\; \frac{n+2}{2}\,2^{-n/2}.
\]
\end{theorem}
\fixed{Notes had $\frac{n+2}{4\cdot2^{n/2}}$ (red correction in margin).}

One can do this explicitly for balls too; we are also still missing a boundary term.
\unclear{Probably: balls meeting the boundary of the box; they contribute $O(t^{n-1})$, negligible as $t\to\infty$.}

\begin{remark}
\added{For $n=2$ the bound is the trivial $\dens\le1$; for $n=3$ it gives $5/(4\sqrt2)\approx0.884$, compared
with $\pi/\sqrt{18}\approx0.7405$ for the FCC packing (Rogers' simplex bound gives $\approx0.7797$).
Blichfeldt's bound decays like $2^{-0.5n}$ up to polynomial factors; Rankin (1947) improved the constant,
and Kabatiansky--Levenshtein (1978) improved the exponent to $2^{-0.599n}$.}
\end{remark}

% MAT670 -- Lecture 3 (18.4.2016) -- transcribed from MAT670_3, pp. 32--48 of the handwritten notes
\chapter[Lecture 3: Blichfeldt gauges, Kabatjanskii--Levenshtein]{Lecture 3: Blichfeldt gauges and the Kabatjanskii--Levenshtein bound}
\label{L3:ch}

\noindent\emph{Lecture of 18 April 2016.}\orig{Notes pp.~32--48.}
Last time we proved the Fejes T\'oth inequality for point configurations on the sphere and
Blichfeldt's sphere packing bound. Today we make some remarks on both, look at the landscape of
known bounds for the packing density, prove the Kabatjanskii--Levenshtein-type bound via a Blichfeldt
gauge, and then extend the notion of a Blichfeldt gauge from balls to general convex bodies.

\added{Throughout, ``a packing of unit balls'' means a family of balls of radius $1$ in $\R^n$ with
pairwise disjoint interiors, i.e.\ with centers at mutual distance at least $2$.}

%======================================================================
\section{Some remarks}
\label{L3:sec:remarks}

\subsection*{Where is the Fejes T\'oth inequality sharp?}

Recall the Fejes T\'oth inequality: \added{if $n\ge 3$ points are placed on the unit sphere $\Sph^2\subset\R^3$,
then the minimal angular distance $d$ between two of them satisfies}
\begin{equation}\label{L3:eq:FT}
  d \;\le\; \arccos\!\left(\frac{\cot^2\omega-1}{2}\right),
  \qquad \omega=\frac{n}{n-2}\cdot\frac{\pi}{6}.
\end{equation}
\fixed{Notes have $\cot(\omega)$; the square is needed (check $n=12$).}%
The right-hand side is the angular edge length of the equilateral spherical triangle with area
$\frac{4\pi}{2n-4}$ on the unit sphere. \added{(A triangulation of the sphere with $n$ vertices has $2n-4$
triangular faces, so this is the area of a face when the $4\pi$ of total area is split evenly.)}

So the bound should be sharp exactly when the points are the vertices of an equilateral deltahedron
\added{(a polytope all of whose faces are congruent equilateral triangles) inscribed in the sphere}, i.e.\ for
\[
  n = 3,\ 4,\ 6,\ 12,
\]
and perhaps asymptotically as $n\to\infty$.
\begin{addedblock}
Check: $n=3$ gives $\omega=\pi/2$, $d\le\arccos(-\tfrac12)=120^\circ$ (an equilateral triangle on a great circle, the degenerate case);
$n=4$ gives $\cot^2\omega=\tfrac13$, $d\le\arccos(-\tfrac13)\approx109.47^\circ$ (regular tetrahedron);
$n=6$ gives $\omega=\pi/4$, $d\le 90^\circ$ (regular octahedron);
$n=12$ gives $\omega=\pi/5$, $d\le\arccos(1/\sqrt5)\approx 63.43^\circ$ (regular icosahedron).
Equality requires a triangulation by congruent equilateral triangles, so these are the only cases of equality;
as $n\to\infty$ the bound is asymptotically sharp.
\end{addedblock}
\orig{$n=3$ underlined in the notes.}

There are other proofs of the Fejes T\'oth inequality, but they are longer.

%======================================================================
\section{Aside: Blichfeldt's gauge revisited}
\label{L3:sec:aside}
\orig{Numbered ``3.2 aside''. Crossed out above: ``Again, there exists e.g.\ a trivial Blichfeldt gauge for $C$, given by $\chi_C$.''}

We saw last time that
\begin{equation}\label{L3:eq:f0}
  f_0(x)=\begin{cases} 1-\tfrac12\abs{x}^2, & \abs{x}\le\sqrt2,\\[2pt] 0, & \abs{x}>\sqrt2,\end{cases}
\end{equation}
is a Blichfeldt gauge for $\Ball^n$: \added{for every packing of unit balls with centers $c_1,c_2,\dots$,
$\sum_i f_0(x-c_i)\le 1$ for all $x\in\R^n$, and consequently the packing density satisfies
$\dens\le \vol(\Ball^n)/\int_{\R^n}f_0$.}
It is non-trivial for our purposes, that is,
\[
  \int_{\R^n} f_0(x)\,dV \;>\; \int_{\R^n}\chi_{\Ball^n}\,dV=\vol(\Ball^n)
  \qquad\text{(at least for } n>2).
\]
Indeed,
\begin{equation}\label{L3:eq:intf0}
  \int_{\R^n} f_0(x)\,dV=\frac{2^{1+\frac n2}}{n+2}\,\vol(\Ball^n).
\end{equation}
\begin{addedblock}
Proof: in polar coordinates, $dV = n\vol(\Ball^n)\,t^{n-1}\,dt\,d\sigma$ (with $\sigma$ the normalized surface measure), so
\[
  \int f_0 = n\vol(\Ball^n)\int_0^{\sqrt2}\Bigl(1-\frac{t^2}{2}\Bigr)t^{n-1}\,dt
  = \vol(\Ball^n)\Bigl(2^{n/2}-\frac{n}{2(n+2)}2^{(n+2)/2}\Bigr)
  = \frac{2\cdot 2^{n/2}}{n+2}\vol(\Ball^n).
\]
Hence $\int f_0>\vol(\Ball^n)$ iff $2^{1+n/2}>n+2$, which holds exactly for $n\ge3$; for $n=2$ there is equality and for
$n=1$ the gauge is worse than $\chi_{\Ball^1}$.
\end{addedblock}

\subsection*{From the gauge to the density bound}
Consider a packing $P$ of unit balls and a large cube $\mathrm{Box}(t)$ of side length $t$.
\begin{addedblock}
Split the balls meeting $\mathrm{Box}(t)$ into the \emph{inner} ones, whose centers lie in the concentric cube of side
$t-2\sqrt2$, and the \emph{boundary} ones. The support of $f_0(\cdot-c_i)$ for an inner center lies in $\mathrm{Box}(t)$, so
integrating $\sum_i f_0(x-c_i)\le 1$ over $\mathrm{Box}(t)$ gives
$N_{\mathrm{in}}\int f_0\le t^n$. The boundary centers lie in a shell of bounded width around $\partial\,\mathrm{Box}(t)$,
so there are at most $C\,t^{n-1}$ of them (by a volume argument, $C$ depending only on $n$).
\end{addedblock}
Therefore, for all $t$,
\[
  \frac{\vol\bigl(P\cap \mathrm{Box}(t)\bigr)}{\vol\bigl(\mathrm{Box}(t)\bigr)}
  \;\le\; \frac{n+2}{2^{1+\frac n2}} \;+\; \frac{C\,t^{n-1}}{t^n},
\]
\fixed{Notes: main term times $(1+\frac{\sqrt2-2}{t})^n$ (Lecture~2: $2\sqrt2-2$); not needed here, tends to $1$.}%
and letting $t\to\infty$,
\begin{equation}\label{L3:eq:blichfeldt}
  \dens(\Ball^n)\;\le\;\limsup_{t\to\infty}\frac{\vol(P\cap\mathrm{Box}(t))}{\vol(\mathrm{Box}(t))}\;\le\;\frac{n+2}{2^{1+\frac n2}}.
\end{equation}

\begin{remark}
We could equally well use
\[
  \dens=\limsup_{t\to\infty}\frac{\vol\bigcup\{\,\Ball\in P:\ \Ball\subseteq\mathrm{Box}(t)\,\}}{\vol\bigl(\mathrm{Box}(t)\bigr)},
\]
\added{counting only the balls entirely inside the box; since the balls meeting $\partial\,\mathrm{Box}(t)$ have total volume
$O(t^{n-1})$, this gives the same limit.}
\end{remark}

%======================================================================
\section{What does this look like?}
\label{L3:sec:landscape}
\orig{Section number overwritten: ``3.3'' / ``4''.}

The bound \eqref{L3:eq:blichfeldt} coming from $f_0$ gives the following density bounds
\added{(two-digit values as in the notes)}:
\begin{center}
\begin{tabular}{c|cccccccccc}
\toprule
$n$ & 1 & 2 & 3 & 4 & 5 & 6 & 7 & 8 & 9 & 10\\
\midrule
$\dfrac{n+2}{2^{1+n/2}}$ & 1.06 & 1 & .88 & .75 & .61 & .5 & .39 & .31 & .24 & .19\\
\bottomrule
\end{tabular}
\end{center}
\added{For comparison, the true optimal densities are $1$, $\pi/\sqrt{12}\approx .9069$, $\pi/\sqrt{18}\approx .7405$ in dimensions $1,2,3$,
and $\pi^4/384\approx .2537$ in dimension $8$.}

\subsection*{Asymptotic bounds}
\begin{center}
\begin{tabular}{p{0.44\textwidth}|p{0.44\textwidth}}
\textbf{Lower bounds} & \textbf{Upper bounds}\\
\midrule
$\dfrac{1}{2^n}$ \ (saturation) & $\dfrac{n+2}{2^{\frac n2+1}}$ \ (Blichfeldt)\\[10pt]
\quad$\vdots$ & \quad$\vdots$\\[4pt]
Ball (1992): $\dfrac{2n}{2^n}$ & $2^{-0.599n}$ \ (Kabatjanskii--Levenshtein)\\[10pt]
Vance, $n\equiv 0 \bmod 4$: $\dfrac{6}{e}\cdot\dfrac{n}{2^n}$ & \quad (Cohn, Elkies, Kumar, \dots)\\[10pt]
Venkatesh (sparse) \added{(2013)}: $\dfrac{1}{2}\cdot\dfrac{n\log\log n}{2^n}$ \ (!!!) & \\
\end{tabular}
\end{center}
\fixed{Notes: $e^{-\gamma}/2$; Venkatesh proves $\Delta_n\ge n\log\log n/2^{n+1}$, i.e.\ $1/2$.}%
\begin{addedblock}
Comments. (1) \emph{Saturation:} if a packing of unit balls is saturated (no further ball can be added), then the balls of
radius $2$ with the same centers cover $\R^n$, so $2^n\dens\ge 1$; saturated packings exist (e.g.\ by Zorn's lemma or a greedy
construction), hence $\dens(\Ball^n)\ge 2^{-n}$. (2) Ball (1992) proved $\dens\ge 2(n-1)\zeta(n)2^{-n}$ for lattice packings;
Vance (2011) proved $\dens\ge \frac{6n}{e\,2^n}$ for lattice packings when $4\mid n$; Venkatesh (\emph{A note on sphere packings in
high dimension}, IMRN 2013) proved $\dens\ge 65963\,n\,2^{-n}$ for all large $n$ and
$\dens\ge \frac12\,n\log\log n\,2^{-n}$ for infinitely many $n$ (the ``!!!'': the first improvement over linear growth of the
numerator). The $\log\log n$ comes from cyclotomic lattices: by Mertens' theorem $m/\varphi(m)$ can be as large as about
$e^{\gamma}\log\log m$. (3) \emph{Update after 2016:} Campos, Jenssen, Michelen and Sahasrabudhe (2023) proved
$\dens\ge(\frac12-o(1))\,n\log n\,2^{-n}$ for all $n$, and Klartag (2025) proved $\dens\ge c\,n^2 2^{-n}$ (even for lattice packings).
On the upper side, the exponent $0.599$ of Kabatjanskii--Levenshtein (1978) has only been improved by subexponential factors
(e.g.\ Cohn--Zhao 2014).
\end{addedblock}

\subsection*{Known good constructions}
There are known good constructions in low dimensions and special ones \added{(e.g.\ $D_4$, $E_8$, the Leech lattice $\Lambda_{24}$)}.
We need to talk about these at some point.\orig{Margin doodle (sad face) omitted.}

\subsection*{Exact results (upper bounds that are attained)}
\begin{itemize}
\item $d=1$: trivial.
\item $d=2$: we proved it. \added{Thue (1892, 1910); the first complete proof is due to L.~Fejes T\'oth (1940). The optimal lattice
packing was known to Lagrange and Gauss.}\unclear{Notes: ``(1892\dots ???) Thue''.}
\item $d=3$: Hales, twice, around 2000; a complicated geometric analysis. \added{Hales (with Ferguson), announced 1998, published in
Ann.\ of Math.\ 2005; the formal computer verification (Flyspeck project) was completed in 2014.}
\item $d=8$ and $d=24$: one month ago, via linear programming duality: one constructs a special (``magic'') auxiliary function.
\added{Viazovska (dimension 8, March 2016) and Cohn--Kumar--Miller--Radchenko--Viazovska (dimension 24, a week later);
both published in Ann.\ of Math.\ 2017. The LP bound itself is due to Cohn--Elkies (2003).}
\end{itemize}

\begin{remark}
Remember that the cubic lattice $\Z^n$ gives a packing of balls of radius $\frac12$ with density
\[
  \dens(\Z^n)=\vol\bigl(\tfrac12\Ball^n\bigr)=\frac{\pi^{n/2}}{2^n\,(\frac n2)!},
\]
\added{where $(\frac n2)!:=\Gamma(\frac n2+1)$. This decays superexponentially, much faster than the lower bound $2^{-n}$.}
\end{remark}

%======================================================================
\section{Kabatjanskii and Levenshtein}
\label{L3:sec:KL}

\begin{definition}\label{L3:def:M}
For $\varphi\in(0,\pi]$, let $M(d,\varphi)$ be the maximum number of points on $\Sph^{d-1}$ with pairwise
\added{angular} distance at least $\varphi$.
Equivalently, $M(d,\varphi)$ is the maximum number of points in $\R^d\setminus\{0\}$ such that the angle
\added{at the origin} spanned by any two of them is at least $\varphi$ \added{(project radially onto the sphere)}.
\end{definition}
\added{For example $M(d,\pi)=2$, and $M(d,\pi/3)=\kiss_d$ is the kissing number ($\kiss_2=6$, $\kiss_3=12$, $\kiss_4=24$, $\kiss_8=240$,
$\kiss_{24}=196560$).}

\begin{theorem}\label{L3:thm:KLgauge}
\added{In any dimension $d$}\orig{Notes: ``in high dimensions''.} and for $1<r<2$, the function
\[
  f_d(x)=\begin{cases} M\bigl(d,\arccos(1-\tfrac{2}{r^2})\bigr)^{-1}, & \abs{x}<r,\\[2pt] 0, & \abs{x}\ge r,\end{cases}
\]
is a Blichfeldt gauge for $\Ball^d$.
\end{theorem}

The key geometric fact is the following.

\begin{lemma}\label{L3:lem:angle}
Let $r<2$, let $x\in\R^d$ and let $y,z\in\R^d$ with $\abs{y-z}\ge 2$ and $\abs{x-y},\abs{x-z}\le r$. Then the angle $\theta$ at $x$
spanned by the segment $[y,z]$ satisfies $\cos\theta\le 1-\frac{2}{r^2}$, i.e.\ $\theta\ge\arccos(1-\frac2{r^2})$.
\end{lemma}

\begin{figure}[ht]
\centering
\begin{tikzpicture}[scale=1.5]
  \coordinate (X) at (0,0);
  \coordinate (Y) at (1.249,1.0);
  \coordinate (Z) at (1.249,-1.0);
  \draw[dashed,gray] (X) circle (1.6);
  \draw[thick] (X) -- (Y) node[midway,above left] {$r$};
  \draw[thick] (X) -- (Z) node[midway,below left] {$r$};
  \draw[thick] (Y) -- (Z) node[midway,right] {$2$};
  \draw (0.4,0.32) arc (38.66:-38.66:0.512);
  \node at (0.75,0) {$\theta$};
  \fill (X) circle (1.2pt) node[left] {$x$};
  \fill (Y) circle (1.2pt) node[above right] {$y$};
  \fill (Z) circle (1.2pt) node[below right] {$z$};
\end{tikzpicture}
\caption{The extremal case of Lemma~\ref{L3:lem:angle}: both centers at distance exactly $r$ from $x$ and at distance $2$ from each
other; then $\sin(\theta/2)=1/r$.}
\label{L3:fig:angle}
\end{figure}

\begin{proof}
First suppose both points are at distance exactly $r$ from $x$ (the worst case, see Figure~\ref{L3:fig:angle}).
Put $a=(y-x)/r$, $b=(z-x)/r$, unit vectors with $a\cdot b=\cos\theta$, and let $\ell=\abs{y-z}\ge2$. Then
\[
  \Bigl|\frac2r\Bigr|^2\le\Bigl|\frac{\ell}{r}\Bigr|^2=\abs{a-b}^2=2-2\,a\cdot b=2-2\cos\theta,
\]
so $\frac{4}{r^2}\le 2-2\cos\theta$, i.e.\ $\frac2{r^2}\le1-\cos\theta$, which is
\[
  \boxed{\cos\theta\le 1-\frac{2}{r^2}.}
\]
\begin{addedblock}
Why this is the worst case: in general, with $\alpha=\abs{x-y}$, $\beta=\abs{x-z}\in(0,r]$, the law of cosines gives
$\cos\theta=\frac{\alpha^2+\beta^2-\ell^2}{2\alpha\beta}\le h(\alpha,\beta):=\frac{\alpha^2+\beta^2-4}{2\alpha\beta}$, and
$\partial_\alpha h=\frac{\alpha^2-\beta^2+4}{2\alpha^2\beta}>0$ because $\beta\le r<2$. By symmetry $h$ is increasing in each variable,
so $h(\alpha,\beta)\le h(r,r)=1-\frac2{r^2}$. Equivalently, the apex angle of the isosceles triangle with legs $r$ and base $2$ is
$2\arcsin(1/r)$, and $\cos(2\arcsin\frac1r)=1-\frac2{r^2}$.
\end{addedblock}
\fixed{Side sketch says apex angle $2\arccos(1/r)$; it is $2\arcsin(1/r)$.}%
\end{proof}

\begin{proof}[Proof of Theorem~\ref{L3:thm:KLgauge}]
Take any packing of unit balls in $\R^d$ and $r<2$. For every $x\in\R^d$ there are at most $M\bigl(d,\arccos(1-2/r^2)\bigr)$ centers at
distance less than $r$ from $x$: by Lemma~\ref{L3:lem:angle}, the angle at $x$ spanned by any segment of length $\ge2$ whose ends are within
$r$ of $x$ is at least $\arccos(1-2/r^2)$, so the directions from $x$ to these centers form a configuration as in
Definition~\ref{L3:def:M}. \added{(If $x$ is itself a center, all other centers are at distance $\ge2>r$, and the count is $1$.)}
Hence $\sum_i f_d(x-c_i)\le 1$.
\end{proof}

\begin{corollary}\label{L3:cor:KL}
For $1\le r<2$,
\[
  \dens(\Ball^d)\le r^{-d}\,M\bigl(d,\arccos(1-\tfrac{2}{r^2})\bigr).
\]
Equivalently, substituting $r=1/\sin(\varphi/2)$,
\[
  \dens(\Ball^d)\le \sin\bigl(\tfrac\varphi2\bigr)^d\,M(d,\varphi),\qquad \frac\pi3<\varphi\le\pi.
\]
\end{corollary}
\begin{proof}
\added{$\int f_d=r^d\vol(\Ball^d)/M$, so $\dens\le\vol(\Ball^d)/\int f_d=r^{-d}M$. For the substitution: $1-\frac{2}{r^2}=1-2\sin^2\frac\varphi2=\cos\varphi$,
and $r\in[1,2)$ corresponds to $\varphi\in(\pi/3,\pi]$.}
\end{proof}
\added{Letting $r\to2$ recovers the simple bound $\dens(\Ball^d)\le \kiss_d\,2^{-d}$.}

For $d$ large, take $\varphi\approx63^\circ$; the Kabatjanskii--Levenshtein bound on spherical codes gives
\[
  \sin\bigl(\tfrac\varphi2\bigr)^d M(d,\varphi)\le 2^{-0.599d+o(d)},
\]
and hence
\[
  \dens(\Ball^d)\le 2^{-0.599d+o(d)}.
\]
\begin{addedblock}
More precisely, Kabatjanskii and Levenshtein (1978), using linear programming bounds for spherical codes, proved
$M(d,\varphi)\le 2^{d(c(\varphi)+o(1))}$ for $0<\varphi<\pi/2$, where with $s=\sin\varphi$,
\[
  c(\varphi)=\frac{1+s}{2s}\log_2\frac{1+s}{2s}-\frac{1-s}{2s}\log_2\frac{1-s}{2s}.
\]
The exponent $c(\varphi)+\log_2\sin(\varphi/2)$ is minimized at $\varphi^*\approx 62.9974^\circ$, where it equals $-0.5990\ldots$
(see Conway--Sloane, \emph{Sphere Packings, Lattices and Groups}, Ch.~1 and~9).
\end{addedblock}

%======================================================================
\section{Blichfeldt gauges for general bodies}
\label{L3:sec:general}
\orig{Section number overwritten, presumably 3.5.}

So we defined a Blichfeldt gauge for $\Ball^n$. Can we extend this? Abstractly:

\begin{definition}\label{L3:def:gauge}
Given a (convex) body $C$ in $\R^n$, a function $f:\R^n\to\R$ is a \emph{Blichfeldt gauge for $C$} if for any collection of
isometries $\{\varphi_i\}_{i=1}^\infty$ of $\R^n$ such that $\{\varphi_iC\}_{i=1}^\infty$ is a packing,
\[
  \sum_{i=1}^\infty f(\varphi_i^{-1}x)\le 1\qquad\forall x\in\R^n.
\]
\end{definition}
\added{As for balls, integrating over a large box shows that if $f$ is integrable and compactly supported, then every packing of
congruent copies of $C$ has density at most $\vol(C)/I(f)$, where $I(f):=\int_{\R^n}f$.}

\begin{remark}
One could add a non-negativity condition on $f$; what we really want is that the mass of $f$ is large: $I(f)\ge I(\chi_C)$.
\end{remark}
\begin{remark}
The index set could also be finite.
\end{remark}
\begin{remark}
The function $f$ could be much stranger than what we allow here.\unclear{Rest of the remark illegible: ``\dots accuracy first''?}
\end{remark}
\begin{exercise}[$\ast\ast\ast$]
Think about such functions distributionally.
\end{exercise}

There always exists a Blichfeldt gauge for $C$, e.g.\ $\chi_C$ \added{(the copies $\varphi_iC$ have disjoint interiors, so a point lies in
the interior of at most one of them; here one should take $\chi$ of the interior, or note that boundaries have measure zero)}.

\begin{lemma}\label{L3:lem:scaling}
If $f(x)$ is a Blichfeldt gauge for $C$, then $f(x/\lambda)$ is a Blichfeldt gauge for $\lambda C$, the $\lambda$-scaling of $C$.
\end{lemma}
\added{Indeed, a packing $\{\varphi_i(\lambda C)\}$ is the image under $x\mapsto\lambda x$ of the packing $\{\psi_iC\}$ with
$\psi_i(y)=\lambda^{-1}\varphi_i(\lambda y)$.}
For instance, $f_0$ (supported on the ball of radius $\sqrt2$) is a gauge for $\Ball^n$, and $f_0(x/\lambda)$ (supported on the ball of
radius $\lambda\sqrt2$) is a gauge for $\lambda\Ball^n$.\orig{Diagram of this scaling omitted.}

\begin{definition}\label{L3:def:inradius}
For a body $C$ define its \emph{inradius} $r(C)$ to be the radius of the largest sphere \added{(ball)} contained in $C$.
For $0\le\gamma\le r(C)$ define $C_\gamma$, the \emph{inner parallel body at depth $\gamma$}, to be the set of points
\[
  C_\gamma:=\{x\in C:\ x+\gamma\Ball^n\subseteq C\}.
\]
\end{definition}
\added{If $C$ is convex and compact, then $C_\gamma$ is convex, compact and non-empty for $\gamma\le r(C)$.}
We can define the straight-line distance from a point $x$ to a body $C$ by $d(x,C)$ \added{$=\min_{y\in C}\abs{x-y}$}.

\begin{figure}[ht]
\centering
\begin{tikzpicture}[scale=1.1]
  \draw[thick] (0,0) rectangle (4,2.4);
  \draw[dashed] (0.7,0.7) rectangle (3.3,1.7);
  \node at (2,1.2) {$C_\gamma$};
  \node at (0.35,2.1) {$C$};
  \draw[gray] (3.3,1.7) circle (0.7);
  \draw[->,gray] (3.3,1.7) -- ++(-150:0.7) node[midway,below,black] {\small$\gamma$};
  \fill (3.3,1.7) circle (1.2pt) node[below right] {$x_i$};
  \fill (6,3) circle (1.2pt) node[above right] {$x$};
  \draw[densely dotted] (6,3) -- (3.3,1.7) node[midway,above left] {$d(x,C_\gamma)$};
\end{tikzpicture}
\caption{The inner parallel body $C_\gamma$ and the nearest point $x_i\in C_\gamma$ to $x$; the ball $x_i+\gamma\Ball^n$ lies in $C$.}
\label{L3:fig:parallel}
\end{figure}

\begin{theorem}\label{L3:thm:parallel}
Let $C$ be a convex body, $0<\gamma\le r(C)$, and let $f$ be a radial Blichfeldt gauge for $\Ball^n$, i.e.\ $f(x)=F(\abs x)$.
Define
\[
  g(x)=F\Bigl(\frac{d(x,C_\gamma)}{\gamma}\Bigr).
\]
Then $g$ is a Blichfeldt gauge for $C$.
\end{theorem}

This is fairly easy to see from a picture (Figure~\ref{L3:fig:parallel}).

\begin{proof}
We would like to show that
\[
  \sum_{i=1}^\infty g(\varphi_i^{-1}x)\le1\qquad\forall x\in\R^n,\ \text{for all packings }\{\varphi_iC\}_{i=1}^\infty \text{ of } C.
\]
Consider a point $x\in\R^n$ and, for each $i$, a point $x_i\in\varphi_iC_\gamma$ at distance $d(x,\varphi_iC_\gamma)$ from $x$
\added{(it exists since $C_\gamma$ is compact)}. Then $x_i+\gamma\Ball^n$ is contained in $\varphi_iC$
\added{(note $\varphi_iC_\gamma=(\varphi_iC)_\gamma$)}, and so $\{x_i+\gamma\Ball^n\}_i$ is a packing.
\orig{Parenthetical: ``since $C$ is convex, things are well defined''.}%
Since $F(\abs{y}/\gamma)$ is a Blichfeldt gauge for $\gamma\Ball^n$ (Lemma~\ref{L3:lem:scaling}), and $x$ is at distance
$\abs{x-x_i}=d(x,\varphi_iC_\gamma)$ from $x_i$, we get
\[
  \sum_i g(\varphi_i^{-1}x)=\sum_i F\Bigl(\frac{d(x,\varphi_iC_\gamma)}{\gamma}\Bigr)=\sum_i F\Bigl(\frac{\abs{x-x_i}}{\gamma}\Bigr)\le 1,
\]
\added{using $d(\varphi_i^{-1}x,C_\gamma)=d(x,\varphi_iC_\gamma)$ because $\varphi_i$ is an isometry.}
\end{proof}

\subsection*{Is this useful?}
As with Blichfeldt gauges and other approaches: is this useful?\unclear{``is this useful?'' tentative; could be ``is there a catch?''} If we have a convex body like $A+\Ball^n$ \added{(a Minkowski sum
with a ball)}, it could be, provided there is enough high-dimensional mass.

\paragraph{A better gauge for balls.}
Consider
\[
  f(x)=\begin{cases} f_0(x), & \abs x>1,\\[2pt] 1-f_0\bigl(2-\abs{x}\bigr), & \abs x\le1,\end{cases}
\]
\added{with $f_0$ from \eqref{L3:eq:f0}, viewed as a function of the radius. Explicitly, $f=1$ for $\abs x\le2-\sqrt2$, $f=\frac12(2-\abs x)^2$ for
$2-\sqrt2\le\abs x\le1$, and $f=f_0$ beyond; $f$ is continuous and $f\ge f_0$ everywhere (Figure~\ref{L3:fig:profiles}).}
This leads to
\[
  \dens\le\Bigl(\frac{2}{d+2}\sqrt2^{\,d}\,(1+b_d)\Bigr)^{-1},\qquad
  b_d=\frac{1}{(\sqrt2)^d(d+1)}-(\sqrt2-1)^{d+1}\Bigl(1+\frac{\sqrt2}{d+1}\Bigr).
\]
\unclear{Notes do not verify that $f$ is a gauge; see the added paragraph below.}
\begin{addedblock}
The factor $\frac{2}{d+2}\sqrt2^{\,d}=\int f_0/\vol(\Ball^d)$ is Blichfeldt's, and $b_d=\int(f-f_0)\big/\int f_0$
(we checked this integral numerically against the formula for $d\le10$). The notes do not prove that $f$ is a Blichfeldt gauge.
The idea: a point $x$ lies in at most one ball of the packing, say at distance $s=\abs{x-c_0}\le1$ from its center, and all other centers
are at distance $\ge2-s$ from $x$; one has to show that their total $f_0$-contribution is at most $f_0(2-s)$, so that it is
compensated by replacing $f_0(s)$ with $1-f_0(2-s)$. Blichfeldt's inequality $\sum_i\abs{x-c_i}^2\ge 2(m-1)$ alone only gives the weaker
bound $s^2/2$ for this contribution, so a finer argument is needed. (A numerical search in dimensions $2,3,4$ found no violation.)
\end{addedblock}

\begin{figure}[ht]
\centering
\begin{tikzpicture}[xscale=3.4,yscale=2.2]
  \draw[->] (0,0) -- (1.7,0) node[right] {$\abs x$};
  \draw[->] (0,0) -- (0,1.2);
  \draw[thick,gray,dashed] plot[domain=0:1.41421,samples=60] (\x,{1-\x*\x/2});
  \draw[thick] (0,1) -- (0.58579,1);
  \draw[thick] plot[domain=0.58579:1,samples=40] (\x,{(2-\x)*(2-\x)/2});
  \draw[thick] plot[domain=1:1.41421,samples=40] (\x,{1-\x*\x/2});
  \draw[thick] (1.41421,0) -- (1.65,0);
  \foreach \t/\l in {0.58579/{2-\sqrt2},1/1,1.41421/{\sqrt2}} \draw (\t,0.03) -- (\t,-0.03) node[below] {\small$\l$};
  \node[left] at (0,1) {$1$};
  \node[gray] at (0.35,0.8) {\small$f_0$};
  \node at (0.85,0.95) {\small$f$};
\end{tikzpicture}
\caption{Radial profiles of Blichfeldt's gauge $f_0$ (dashed) and the improved function $f$ (solid); they agree for $\abs x\ge1$.}
\label{L3:fig:profiles}
\end{figure}

The resulting bounds compared with those from $f_0$ \added{(truncated to two digits)}:
\begin{center}
\begin{tabular}{c|cccccccccc}
\toprule
$d$ & 1 & 2 & 3 & 4 & 5 & 6 & 7 & 8 & 9 & 10\\
\midrule
bound from $f$ & 1 & .94 & .84 & .72 & .60 & .49 & .39 & .31 & .24 & .18\\
bound from $f_0$ & 1.06 & 1 & .88 & .75 & .61 & .5 & .39 & .31 & .24 & .19\\
\bottomrule
\end{tabular}
\end{center}
\orig{An arrow on p.~47 links $.94$ ($f$, $d=2$) with $1$ ($f_0$, $d=2$).}

\paragraph{Bodies with a lot of 2-dimensional mass.}
So this works for things with a lot of ``2-dimensional'' (or higher-dimensional) mass. For a long capsule in $\R^3$
(a segment plus a ball), the cross-sections are discs and Theorem~\ref{L3:thm:parallel} gives a bound of $.94+\text{error}$.
By contrast, for a long box \added{in the plane} the cross-sections are $1$-dimensional and we only get $1+\text{error}$: the method gives
nothing, since there is no nontrivial $1$-dimensional Blichfeldt gauge.
\unclear{Reading of the box sketch and of ``hence doesn't work as there is no 1d Blichfeldt gauge'' is tentative.}
\added{(In dimension $1$, intervals tile the line, so every gauge has $I(f)\le\vol(\Ball^1)$.)}

\begin{figure}[ht]
\centering
\begin{tikzpicture}[scale=0.9]
  \draw[thick] (0,1) -- (4,1) arc (90:-90:1) -- (0,-1) arc (270:90:1);
  \draw[very thick] (0,0) -- (4,0) node[midway,above] {$\ell$};
  \fill (0,0) circle (1.3pt) (4,0) circle (1.3pt);
  \draw[->,gray] (4,0) -- ++(-40:1) node[midway,below left,black] {\small$1$};
\end{tikzpicture}
\caption{An $\ell$-capsule: a segment of length $\ell$ plus a unit ball, drawn in a plane through the axis.}
\label{L3:fig:capsule}
\end{figure}

\begin{example}[$\ell$-capsules]\label{L3:ex:capsule}
Let $C=S_\ell+\Ball^d$, where $S_\ell$ is a segment of length $\ell$ \added{(the notes write ``$\ell\times\Ball^n$'')}.
\added{Then $r(C)=1$ and $C_1=S_\ell$. Take a radial gauge $F(\abs x)$ for $\Ball^d$ and $\gamma=1$, so $g(x)=F(d(x,S_\ell))$.
Write $I_k(F)=\int_{\R^k}F(\abs y)\,dy$. Points whose nearest point on $S_\ell$ is interior contribute a cylinder over $S_\ell$, and the two
ends together contribute a full ball's worth, so}
\[
  I(g)=\ell\cdot\frac{I_{d-1}(F)}{\vol(\Ball^{d-1})}\,\vol(\Ball^{d-1})+\frac{I_d(F)}{\vol(\Ball^d)}\,\vol(\Ball^d)
  =\ell\,I_{d-1}(F)+I_d(F),
\]
and
\[
  \vol(\ell\text{-capsule})=\ell\vol(\Ball^{d-1})+\vol(\Ball^d).
\]
Hence
\[
  \dens(\ell\text{-capsule})\le\frac{\ell\vol(\Ball^{d-1})+\vol(\Ball^d)}{\ell\,I_{d-1}(F)+I_d(F)}.
\]
\added{As $\ell\to\infty$ this tends to $\vol(\Ball^{d-1})/I_{d-1}(F)$, the $(d-1)$-dimensional ball bound for the same profile; for $d=3$ and
the improved $f$ above this is the value $.94$ from the table (assuming $f$ is a gauge in $\R^3$).}
\end{example}

\begin{remark}
\added{For a genuine cylinder $S_\ell\times\Ball^{d-1}$ (flat ends) the same computation works up to} end errors.
\end{remark}


\part{Lattices and root systems}
% MAT670 -- Lectures 4 and 5 (first part), transcribed from MAT670_46.pdf, pages 1--23
% Local drawing helpers for Dynkin diagrams (prefixed to avoid clashes).
% Vertices sit at (x,0); edges are drawn before vertices.
\newcommand{\LIVv}[1]{\filldraw[fill=white,draw=black] (#1) circle (0.09);}
\newcommand{\LIVe}[2]{\draw (#1) -- (#2);}
% double edge from (x,0) to (x+1,0), arrow pointing right (R) or left (L)
\newcommand{\LIVdR}[1]{\draw (#1,0.045) -- (#1+1,0.045); \draw (#1,-0.045) -- (#1+1,-0.045); \draw (#1+0.42,0.13) -- (#1+0.58,0) -- (#1+0.42,-0.13);}
\newcommand{\LIVdL}[1]{\draw (#1,0.045) -- (#1+1,0.045); \draw (#1,-0.045) -- (#1+1,-0.045); \draw (#1+0.58,0.13) -- (#1+0.42,0) -- (#1+0.58,-0.13);}
\newcommand{\LIVdN}[1]{\draw (#1,0.045) -- (#1+1,0.045); \draw (#1,-0.045) -- (#1+1,-0.045);}
% triple edge from (x,0) to (x+1,0), arrow pointing right / none
\newcommand{\LIVtR}[1]{\draw (#1,0.07) -- (#1+1,0.07); \draw (#1,0) -- (#1+1,0); \draw (#1,-0.07) -- (#1+1,-0.07); \draw (#1+0.4,0.16) -- (#1+0.6,0) -- (#1+0.4,-0.16);}
\newcommand{\LIVtN}[1]{\draw (#1,0.07) -- (#1+1,0.07); \draw (#1,0) -- (#1+1,0); \draw (#1,-0.07) -- (#1+1,-0.07);}

%======================================================================
\chapter[Lecture 4: Classification of root systems]{Lecture 4: The classification of (crystallographic) root systems}
\label{L04a:ch:L4}
\orig{Heading reads ``Lecture 4) (and 5?)''.}

%----------------------------------------------------------------------
\section{Motivation and lattice basics}
\label{L04a:sec:4.1}

In this lecture we examine \emph{root lattices}, coming from \emph{crystallographic} root systems.
One motivation is to find nice constructions of good packings in low dimensions. Another is to show
that such simple constructions do not really extend to high dimensions in a good way.
\orig{crossed out: ``there are not such simple constructions''}

The best known packings up through $\R^8$ are root lattices. In fact $A_1$, $A_2$, $A_3$ and $E_8$ are
optimal among \emph{all} packings in their dimensions.
\added{($A_1$: trivial; $A_2$: Thue, Fejes T\'oth 1940; $A_3$, the FCC packing: Hales 1998/2005, the Kepler
conjecture; $E_8$: Viazovska 2016. In dimension $24$ the Leech lattice, which is not a root lattice, is
optimal by Cohn--Kumar--Miller--Radchenko--Viazovska 2016.)}
\unclear{Box reads ``$A_1,A_2,A_3,E_8$ optimal as gen.\ packings''; ``$A_3$'' read from a blurred subscript.}

\begin{remark}
The classification we are about to give covers, in effect, half of the classification of complex
semisimple Lie algebras (equivalently, of complex semisimple Lie groups): the irreducible root systems
correspond to the simple ones.
\added{(See e.g.\ Humphreys, \emph{Introduction to Lie Algebras and Representation Theory}, Ch.~III, or
Fulton--Harris, \emph{Representation Theory}, Lecture~21.)}
\orig{Messy boxed aside, partly scribbled over; a book reference is unreadable.}
\end{remark}

\begin{definition}\label{L04a:def:lattice}
A \emph{lattice} $\Lat$ is a discrete additive subgroup of $\R^n$. For our purposes it will always be
co-compact, i.e.\ have full rank. Such a lattice has a $\Z$-basis $\{a_1,\dots,a_n\}$,
\added{$\Lat = \Z a_1\oplus\dots\oplus\Z a_n$,} that is at the same time an $\R$-basis of $\R^n$.
\end{definition}

\begin{example}
$\{(1,0)^T,(0,1)^T\}$ is a basis of the lattice $\Z^2\subset\R^2$.
\end{example}

A lattice can have many different bases. However, two bases $\{a_i\}$, $\{a_i'\}$ differ by an
\emph{integer} change of basis $K=(k_{ij})$ with integral inverse:
\[
  a_i' = \sum_{j=1}^n k_{ij} a_j, \qquad k_{ij}\in\Z
  \quad\Longrightarrow\quad \det K = \pm 1
\]
\added{(since $\det K$ and $\det K^{-1}$ are both integers whose product is $1$).}
So the volume of the torus $\R^n/\Lat$,
\[
  \vol(\R^n/\Lat) = \abs{\det\Lat} = \vol\Big(\Big\{\textstyle\sum_i t_i a_i : 0\le t_i\le 1\Big\}\Big),
\]
\added{i.e.\ the volume of the parallelepiped spanned by $a_1,\dots,a_n$, where $\det\Lat$ denotes the
determinant of the matrix with columns $a_1,\dots,a_n$,} is a clear invariant of $\Lat$.

The largest balls that can be placed at the lattice points without overlapping have radius
$\tfrac12\min_{\lambda\in\Lat\setminus\{0\}}\abs{\lambda}$, so each lattice gives a packing
\[
  \Big\{\, \varphi_i + \Ball\big(\tfrac12 \min_{\lambda\in\Lat\setminus\{0\}} \abs{\lambda}\big) \;:\; \varphi_i\in\Lat \Big\}.
\]
\begin{addedblock}
Its density is
$\dens(\Lat) = \vol \Ball^n\!\big(\tfrac12\lambda_1(\Lat)\big)\big/\abs{\det\Lat}$, where
$\lambda_1(\Lat)=\min_{\lambda\neq0}\abs{\lambda}$.
\end{addedblock}

Since scaling does not change the density, we can normalize to $\abs{\det\Lat}=1$ (i.e.\ consider
lattices $g\Z^n$ with $g\in\SL_n(\R)$).
\fixed{notes: ``(to $SL_n$, compact)''; the space is not compact}
\added{The space $\SL_n(\R)/\SL_n(\Z)$ of such lattices is \emph{not} compact (e.g.\ the lattices
$\operatorname{diag}(t,t^{-1},1,\dots,1)\Z^n$, $t\to0$, leave every compact set). What is true is
Mahler's compactness criterion: for each $\varepsilon>0$ the set of unimodular lattices with
$\lambda_1\ge\varepsilon$ is compact.}
Then there is clearly an upper bound on the length of the shortest vector among lattices with
$\abs{\det\Lat}=1$, and so a bound on the density.
\added{(Trivially, the balls are disjoint, so $\vol\Ball^n(\lambda_1/2)\le 1$; Minkowski's convex body
theorem gives the sharper $\vol\Ball^n(\lambda_1)\le 2^n$. Since lattices with very small $\lambda_1$
have small density, the supremum of the density can be taken over a set
$\{\lambda_1\ge\varepsilon\}$, which is compact by Mahler; so an optimal lattice exists in each
dimension.)}

So\dots\ what is the optimal lattice? The optimal lattice packings are known up to $n=8$
(Blichfeldt 1935 for $n=6,7,8$), and they are
\[
  A_1,\ A_2,\ A_3,\ D_4,\ D_5,\ E_6,\ E_7,\ E_8 .
\]
These are root lattices.
\orig{crossed out: ``still open''}
\added{(Earlier cases: Lagrange for $n=2$, Gauss for $n=3$, Korkine--Zolotarev for $n=4,5$. The only
other dimension where the optimal lattice is known is $n=24$: the Leech lattice, Cohn--Kumar 2009.)}

%----------------------------------------------------------------------
\section{Good candidates in low dimensions}
\label{L04a:sec:4.2}

\begin{definition}\label{L04a:def:integral}
A lattice $\Lat$ is \emph{integral} if $\ip{x}{y}\in\Z$ for all $x,y\in\Lat$.
A \emph{root lattice} is an integral lattice generated by roots, i.e.\ by elements of a root system
(Definition~\ref{L04a:def:rootsystem}).
\end{definition}

\begin{addedblock}
In lattice language one usually calls a vector $v$ of an integral lattice a root if
$\ip{v}{v}\in\{1,2\}$; then $2\ip{x}{v}/\ip{v}{v}\in\Z$ for all $x\in\Lat$, so the reflection $s_v$
(below) maps $\Lat$ to itself, and the roots of $\Lat$ form a root system. The irreducible lattices
generated by vectors of norm $2$ are exactly $A_n$ ($n\ge1$), $D_n$ ($n\ge4$), $E_6,E_7,E_8$ (Witt);
see Conway--Sloane, \emph{Sphere Packings, Lattices and Groups}, Ch.~4.
\end{addedblock}

\begin{definition}\label{L04a:def:rootsystem}
A \emph{root system} is a subset $\Delta\subset\R^n\setminus\{0\}$ satisfying:
\begin{enumerate}[label=\arabic*)]
\item $\Delta$ is finite and spans $\R^n$;
\item if $\alpha\in\Delta$, then $n\alpha\in\Delta \iff n=\pm1$ \quad(\emph{reduced});
\item $\Delta$ is invariant under the orthogonal reflection in the hyperplane $\alpha^\perp$, for any
  $\alpha\in\Delta$: writing
  \[
    \operatorname{proj}_\alpha\beta = \alpha\,\frac{\ip{\beta}{\alpha}}{\ip{\alpha}{\alpha}},\qquad
    s_\alpha\beta := \beta - 2\operatorname{proj}_\alpha\beta ,
  \]
  we require $s_\alpha\beta\in\Delta$ for all $\alpha,\beta\in\Delta$;
\item (\emph{crystallographic}) for all $\alpha,\beta\in\Delta$,
  \[
    n_{\beta\alpha} := \frac{2\ip{\beta}{\alpha}}{\ip{\alpha}{\alpha}}\in\Z .
  \]
\end{enumerate}
\added{The elements of $\Delta$ are called \emph{roots}.}
\end{definition}

\begin{center}
\begin{tikzpicture}[scale=1.1,>=Stealth]
  \draw[dashed,gray] (0,-0.4) -- (0,1.7) node[above] {\small$\alpha^\perp$};
  \draw[->,thick] (0,0) -- (2,0) node[right] {$\alpha$};
  \draw[->,thick] (0,0) -- (1.2,1.2) node[above right] {$\beta$};
  \draw[dotted] (1.2,1.2) -- (1.2,0);
  \draw[->,very thick,blue!60!black] (0,0) -- (1.2,0);
  \node[below] at (1.2,-0.05) {\small$\operatorname{proj}_\alpha\beta$};
  \draw[->,thick] (0,0) -- (-1.2,1.2) node[above left] {$s_\alpha\beta$};
  \draw[dotted] (-1.2,1.2) -- (1.2,1.2);
\end{tikzpicture}
\end{center}

Property 4) is the restriction that the projection of $\beta$ onto $\alpha$ is an integer or
half-integer multiple of $\alpha$; but it really forces $\Delta$ to give rise to a lattice
\added{(the $\Z$-span of $\Delta$ is discrete: it is the $\Z$-span of a fundamental system, which is a
basis of $\R^n$, see \S\ref{L04a:sec:4.3}).}

\begin{exercise}
Try building some other root systems, i.e.\ sets satisfying 1)--3), which are \emph{not}
crystallographic.
\added{(Hint: the $10$ unit vectors at angles $k\pi/5$ in $\R^2$; the $30$ vertices of an
icosidodecahedron in $\R^3$, i.e.\ type $H_3$.)}
\end{exercise}

\paragraph{Property 4) is strong.}
We have
\[
  n_{\beta\alpha} = \frac{2\ip{\beta}{\alpha}}{\ip{\alpha}{\alpha}} = \frac{2\abs{\beta}}{\abs{\alpha}}\cos\theta,
\]
where $\theta$ is the angle between $\alpha$ and $\beta$. Since both $n_{\beta\alpha}$ and $n_{\alpha\beta}$ are integers,
\[
  \boxed{\,n_{\beta\alpha}\,n_{\alpha\beta} = 4\cos^2\theta\in\Z\,}
  \qquad\Longrightarrow\qquad 4\cos^2\theta\in\{0,1,2,3,4\}.
\]
If $4\cos^2\theta=4$, then $\alpha$ and $\beta$ are parallel, so $\alpha=\beta$ or $\alpha=-\beta$
\added{by the reducedness axiom 2)}.

Assuming $\abs{\alpha}\ge\abs{\beta}$ (and $\alpha\neq\pm\beta$), this leaves the following table.
\begin{center}
\renewcommand{\arraystretch}{1.25}
\begin{tabular}{cccccc}
\toprule
$4\cos^2\theta$ & $n_{\beta\alpha}$ & $n_{\alpha\beta}$ & $\abs{\alpha}/\abs{\beta}$ & $\cos\theta$ & $\theta$\\
\midrule
$3$ & $+1$ & $+3$ & $\sqrt3$ & $+\sqrt3/2$ & $\pi/6$\\
$3$ & $-1$ & $-3$ & $\sqrt3$ & $-\sqrt3/2$ & $5\pi/6$\\
$2$ & $+1$ & $+2$ & $\sqrt2$ & $+\sqrt2/2$ & $\pi/4$\\
$2$ & $-1$ & $-2$ & $\sqrt2$ & $-\sqrt2/2$ & $3\pi/4$\\
$1$ & $+1$ & $+1$ & $1$ & $+1/2$ & $\pi/3$\\
$1$ & $-1$ & $-1$ & $1$ & $-1/2$ & $2\pi/3$\\
$0$ & $0$ & $0$ & (undetermined) & $0$ & $\pi/2$\\
\bottomrule
\end{tabular}
\end{center}
\added{(For $4\cos^2\theta\neq0$ the ratio follows from $n_{\alpha\beta}/n_{\beta\alpha}=\abs{\alpha}^2/\abs{\beta}^2$.)}

\begin{definition}\label{L04a:def:reducible}
A root system is \emph{reducible} (\emph{decomposable}) if $\Delta=\Delta_1\cup\Delta_2$
\added{with $\Delta_1,\Delta_2$ nonempty} such that $\ip{\alpha_1}{\alpha_2}=0$ for all
$\alpha_1\in\Delta_1$, $\alpha_2\in\Delta_2$; i.e.\ all the elements of one part are orthogonal to all
the elements of the other. Otherwise we call it \emph{irreducible} (\emph{indecomposable}).
\end{definition}

\begin{example}[Root systems in dimensions $1$ and $2$]\label{L04a:ex:rank2}
\leavevmode
\begin{center}
\begin{tikzpicture}[>=Stealth,scale=0.95]
  % A1
  \begin{scope}
    \draw[->,thick] (0,0) -- (1,0) node[right] {\small$\alpha$};
    \draw[->,thick] (0,0) -- (-1,0) node[left] {\small$-\alpha$};
    \fill (0,0) circle (1.2pt);
    \node at (0,-1.75) {$A_1$\quad($n=1$)};
  \end{scope}
  % A1 x A1
  \begin{scope}[xshift=4cm]
    \draw[->,thick] (0,0) -- (1,0) node[right] {\small$\alpha$};
    \draw[->,thick] (0,0) -- (-1,0) node[left] {\small$-\alpha$};
    \draw[->,thick] (0,0) -- (0,0.7) node[above] {\small$\beta$};
    \draw[->,thick] (0,0) -- (0,-0.7) node[below] {\small$-\beta$};
    \node at (0,-1.75) {$A_1\times A_1$\ (reducible), $\theta=\frac\pi2$};
  \end{scope}
  % A2
  \begin{scope}[xshift=9.4cm]
    \draw[gray!60] (0:1)--(60:1)--(120:1)--(180:1)--(240:1)--(300:1)--cycle;
    \foreach \a in {0,60,...,300} \draw[->,thick] (0,0) -- (\a:1);
    \node[right] at (0:1) {\small$\alpha$};
    \node[above left] at (120:1) {\small$\beta$};
    \node[above right] at (60:1) {\small$\alpha+\beta=s_\alpha\beta$};
    \node[left] at (180:1) {\small$-\alpha$};
    \node[below left] at (240:1) {\small$s_\alpha(-\beta)$};
    \node[below right] at (300:1) {\small$-\beta$};
    \node at (0,-1.75) {$A_2$, $\theta=\frac\pi3$};
  \end{scope}
\end{tikzpicture}

\medskip
\begin{tikzpicture}[>=Stealth,scale=0.95]
  % B2
  \begin{scope}
    \draw[gray!60] (-1,-1) rectangle (1,1);
    \foreach \v in {(1,0),(0,1),(-1,0),(0,-1),(1,1),(-1,1),(-1,-1),(1,-1)} \draw[->,thick] (0,0) -- \v;
    \node[right] at (1,0) {\small$\alpha$};
    \node[above left] at (-1,1) {\small$\beta$};
    \node[above] at (0,1) {\small$s_\beta\alpha$};
    \node at (0,-1.5) {$B_2$, $\theta=\frac\pi4$};
  \end{scope}
  % G2
  \begin{scope}[xshift=5cm]
    \foreach \a in {0,60,...,300} \draw[->,thick] (0,0) -- (\a:0.7);
    \foreach \a in {30,90,...,330} \draw[->,thick] (0,0) -- (\a:1.212);
    \draw[gray!60] (30:1.212)--(150:1.212)--(270:1.212)--cycle;
    \draw[gray!60] (90:1.212)--(210:1.212)--(330:1.212)--cycle;
    \node at (0,-1.5) {$G_2$, $\theta=\frac\pi6$};
  \end{scope}
\end{tikzpicture}
\end{center}
(Note: $\Delta$ is not a basis!)
\added{In $B_2$ the roots are $\pm e_1,\pm e_2,\pm e_1\pm e_2$; in $G_2$ there are $6$ short roots and
$6$ long roots of length ratio $\sqrt3$. In each case $\theta$ is the smallest angle between two roots.}
\end{example}

%----------------------------------------------------------------------
\section{Classification of root systems: preliminaries}
\label{L04a:sec:4.3}

From a root system $\Delta$ we can choose a (non-unique) subset, the \emph{simple roots}.

For each $\alpha\in\Delta$ consider $\alpha^\perp$, its (linear) orthogonal complement.
\orig{crossed out: ``there is a unique $\perp$ hyperplane''}
Since $\abs{\Delta}$ is finite, there exists $d\in\R^n$ such that $\ip{\alpha}{d}\neq0$ for all
$\alpha\in\Delta$ \added{(avoid the finitely many hyperplanes $\alpha^\perp$)}. Hence there is a
partition of $\Delta$ into
\[
  \text{positive roots } \Delta_d^+ = \{\alpha\in\Delta : \ip{\alpha}{d}>0\},\qquad
  \text{negative roots } \Delta_d^- = \{\alpha\in\Delta : \ip{\alpha}{d}<0\}
\]
\added{(note $\Delta_d^-=-\Delta_d^+$).}

\begin{definition}\label{L04a:def:simple}
A root $\alpha\in\Delta_d^+$ is \emph{simple} if it is not the sum of two other positive roots.
A set of simple roots is a \emph{fundamental system} of $\Delta$ \added{(we write $\Pi=\Pi_d$)}.
\end{definition}

\begin{center}
\begin{tikzpicture}[>=Stealth,scale=1.3]
  \fill[gray!12] (165:1.7) -- (-15:1.7) arc (-15:165:1.7) -- cycle;
  \draw[dashed,red!70!black] (165:1.8) -- (-15:1.8) node[right] {\small$d^\perp$};
  \foreach \a in {0,60,...,300} \draw[->,thick] (0,0) -- (\a:1.2);
  \node[right] at (0:1.2) {\small$\alpha$};
  \node[above left] at (120:1.2) {\small$\beta$};
  \node[above right] at (60:1.2) {\small$\alpha+\beta$};
  \draw[->,red!70!black,thick] (0,0) -- (75:0.9) node[above] {\small$d$};
  \node at (35:1.75) {\small$\Delta_d^+$};
  \node at (-110:1.4) {\small$\Delta_d^-$};
\end{tikzpicture}
\end{center}
\orig{Sketch on p.\,8, redrawn for $A_2$: simple roots $\alpha,\beta$.}

The choice of $d$ affects the fundamental system, but such systems are equivalent under the action of
the group of reflections $\langle s_\alpha\rangle_{\alpha\in\Delta}$
\added{(the \emph{Weyl group} $W$; see Proposition~\ref{L04a:prop:reconstruct}).}

\begin{theorem}\label{L04a:thm:basis}
A fundamental system of $\Delta$ is an $\R$-basis for $\R^n$.
\added{Moreover every positive root is a linear combination of simple roots with nonnegative integer
coefficients.}
\end{theorem}

\begin{proof}[Sketch]
\emph{Linear independence.} Suppose $\sum_i c_i\gamma_i = 0$ with $\gamma_i\in\Pi$. Partition the
indices into those with $c_i>0$ $(+)$, those with $c_i<0$ $(-)$, and those with $c_i=0$, and put
\added{$x := \sum_{c_i>0} c_i\gamma_i = \sum_{c_j<0}\abs{c_j}\gamma_j$.} Then
\[
  \ip{x}{x} = \sum_{c_i>0,\ c_j<0} c_i\,\abs{c_j}\,\ip{\gamma_i}{\gamma_j}\le 0
\]
by Lemma~\ref{L04a:lem:obtuse} (distinct simple roots make non-acute angles), so $x=0$.
\added{But then $0=\ip{x}{d}=\sum_{c_i>0}c_i\ip{\gamma_i}{d}$ with every $\ip{\gamma_i}{d}>0$, so no
$c_i$ is positive; likewise none is negative.}

\emph{Span.} Let $\gamma$ be orthogonal to all simple roots. Then by the orbit argument
(Proposition~\ref{L04a:prop:reconstruct}: every root is the image of a simple root under the group
generated by the simple reflections $s_\alpha$, $\alpha\in\Pi$, each of which fixes $\gamma$),
$\gamma$ is orthogonal to all roots, hence $\gamma=0$ since $\Delta$ spans.
\begin{addedblock}
A more elementary route: if $\alpha\in\Delta_d^+$ is not simple, then $\alpha=\beta+\beta'$ with
$\beta,\beta'\in\Delta_d^+$, and $\ip{\beta}{d},\ip{\beta'}{d}<\ip{\alpha}{d}$. Since $\Delta$ is
finite, induction on $\ip{\alpha}{d}$ shows every positive root is a sum of simple roots. Hence $\Pi$
spans $\operatorname{span}\Delta=\R^n$, and the integrality claim follows too.
\end{addedblock}
\end{proof}

\begin{proposition}\label{L04a:prop:reconstruct}
$\Delta$ can be reconstructed from a fundamental system via reflections.
\added{Precisely: if $W_\Pi$ denotes the group generated by $s_\gamma$, $\gamma\in\Pi$, then
$\Delta = W_\Pi\cdot\Pi$, and any two fundamental systems are related by an element of $W_\Pi$.}
\end{proposition}

\begin{proof}[Sketch]
Consider the \emph{Weyl chamber} of the fundamental system $\Pi$,
\[
  C(\Pi)=\{x\in\R^n : \ip{x}{\gamma}>0 \text{ for all }\gamma\in\Pi\}.
\]
\orig{crossed out: ``orbit $d$ to each\dots''}
Then $d\in C(\Pi)$, and moving $d$ within its chamber does not change the fundamental system.
Any other fundamental system is based on some $d'$ in another chamber, and there are finitely many
chambers \added{(the connected components of $\R^n\setminus\bigcup_{\alpha}\alpha^\perp$)}. Take a path
$d\to d'$ \added{in general position}; the walls it crosses give reflections $s_{\text{walls}}$, whose
composition takes $d$ to (the chamber of) $d'$, and hence takes the fundamental system of $d$ to that
of $d'$.
\begin{addedblock}
Details: if $d$ and $d'=s_\alpha d$ lie in adjacent chambers separated by the wall $\alpha^\perp$,
then $\Pi_{d'}=s_\alpha\Pi_d$, since $s_\alpha$ is an isometry permuting $\Delta$ with
$\ip{s_\alpha\beta}{s_\alpha d}=\ip{\beta}{d}$. The walls of $C(\Pi)$ are exactly the hyperplanes
$\gamma^\perp$, $\gamma\in\Pi$; writing each wall-crossing reflection in terms of the previous ones
($s_{w\gamma}=ws_\gamma w^{-1}$) shows the composite lies in $W_\Pi$. Finally every root $\alpha$ is
simple for a suitable $d'$ (take $d'$ very close to a generic point of $\alpha^\perp$, on the side where
$\ip{\alpha}{d'}>0$), so $\alpha=w\gamma$ for some $w\in W_\Pi$, $\gamma\in\Pi$. See Humphreys,
\emph{Introduction to Lie Algebras and Representation Theory}, \S10.3.
\end{addedblock}
\end{proof}

\begin{center}
\begin{tikzpicture}[>=Stealth,scale=1.6]
  % hyperplanes alpha^perp (90), beta^perp (30), (alpha+beta)^perp (150)
  \draw[magenta,thick] (90:1.5) -- (270:1.5) node[below] {\small$\alpha^\perp$};
  \draw[magenta,thick] (30:1.5) -- (210:1.5) node[below left] {\small$\beta^\perp$};
  \draw[magenta,thick] (150:1.5) -- (330:1.5) node[below right] {\small$(\alpha+\beta)^\perp$};
  \foreach \a in {0,60,...,300} \draw[->,thick] (0,0) -- (\a:1.15);
  \node[right] at (0:1.15) {\small$\alpha$};
  \node[left] at (180:1.15) {\small$s_\alpha(\alpha)$};
  \node[above left] at (120:1.15) {\small$\beta=s_\alpha(\alpha+\beta)$};
  \node[above right] at (60:1.15) {\small$\alpha+\beta=s_\alpha(\beta)$};
  \fill[red!70!black] (75:0.8) circle (1pt) node[right] {\small$d$};
  \fill[green!45!black] (105:0.8) circle (1pt) node[left] {\small$s_\alpha d$};
  \draw[dashed] (75:0.8) to[bend right=15] (105:0.8);
\end{tikzpicture}
\end{center}
\orig{Colour sketch on p.\,11, redrawn.}
For $A_2$: with $d$ as shown, $\Pi_d=\{\alpha,\beta\}$; reflecting across $\alpha^\perp$ gives
$s_\alpha d$, with fundamental system $\Pi_{s_\alpha d}=\{s_\alpha(\alpha),s_\alpha(\beta)\}=\{-\alpha,\alpha+\beta\}$.

\begin{lemma}\label{L04a:lem:diff}
If $\alpha,\beta\in\Delta$ are not parallel and $\ip{\alpha}{\beta}>0$, then $\alpha-\beta\in\Delta$.
\end{lemma}

\begin{proof}
We have $n_{\alpha\beta}=2\ip{\alpha}{\beta}/\ip{\beta}{\beta}>0$, and by the table
\added{($n_{\alpha\beta}n_{\beta\alpha}\in\{1,2,3\}$ with both factors positive integers)}
$n_{\alpha\beta}=1$ or $n_{\beta\alpha}=1$. If $n_{\alpha\beta}=1$, then
\[
  s_\beta(\alpha) = \alpha - 2\beta\,\frac{\ip{\alpha}{\beta}}{\ip{\beta}{\beta}}
  = \alpha-n_{\alpha\beta}\,\beta = \alpha-\beta\in\Delta.
\]
Similarly, if $n_{\beta\alpha}=1$, then $s_\alpha(-\beta)=-\beta+n_{\beta\alpha}\alpha=\alpha-\beta\in\Delta$.
\fixed{written $m_{\beta\alpha}$, $s_\beta(-\beta)$ (typos)}
\end{proof}

\begin{lemma}\label{L04a:lem:obtuse}
If $\alpha,\beta$ are distinct simple roots, then $\ip{\alpha}{\beta}\le0$.
\end{lemma}

\begin{proof}
If $\ip{\alpha}{\beta}>0$, then \added{(as $\alpha\neq\pm\beta$, both being positive)}
Lemma~\ref{L04a:lem:diff} gives $\gamma:=\alpha-\beta\in\Delta$. So $\gamma$ or $-\gamma$ lies in
$\Delta_d^+$, and then $\alpha=\beta+\gamma$ or $\beta=\alpha+(-\gamma)$ is a sum of two positive roots,
contradicting simplicity.
\end{proof}

\begin{proposition}\label{L04a:prop:decomp}
A fundamental system is decomposable \added{(into two nonempty mutually orthogonal subsets)} if and only
if the root system is decomposable.
\end{proposition}

\begin{proof}[Sketch]
Similar arguments.
\begin{addedblock}
If $\Delta=\Delta_1\cup\Delta_2$ orthogonally, then $\Pi=(\Pi\cap\Delta_1)\cup(\Pi\cap\Delta_2)$, and
both parts are nonempty because $\Pi$ spans $\R^n$. Conversely, if $\Pi=\Pi_1\cup\Pi_2$ orthogonally,
the reflections $s_\gamma$, $\gamma\in\Pi_1$, fix $\Pi_2$ pointwise and preserve
$\operatorname{span}\Pi_1$, and vice versa; so by
Proposition~\ref{L04a:prop:reconstruct} every root lies in $W_\Pi\Pi_1\subset\operatorname{span}\Pi_1$ or
in $W_\Pi\Pi_2\subset\operatorname{span}\Pi_2$, which are orthogonal.
\end{addedblock}
\end{proof}

%----------------------------------------------------------------------
\section{Classification via Coxeter graphs}
\label{L04a:sec:4.4}

Simple roots determine root systems (Proposition~\ref{L04a:prop:reconstruct}), so it is enough to
classify fundamental systems.

Elements $\alpha,\beta$ of a fundamental system are orthogonal or make an obtuse angle
(Lemma~\ref{L04a:lem:obtuse}). By the table, therefore,
\[
  \theta \in \Big\{\frac\pi2,\ \frac{2\pi}3,\ \frac{3\pi}4,\ \frac{5\pi}6\Big\},
  \qquad\text{that is,}\qquad n_{\alpha\beta}n_{\beta\alpha}=4\cos^2\theta\in\{0,1,2,3\}.
\]

\begin{definition}\label{L04a:def:coxeter}
The \emph{Coxeter graph} of $\Delta$ has a vertex for each simple root, and an edge of weight
$n_{\alpha\beta}\cdot n_{\beta\alpha}$ between $\alpha$ and $\beta$ \added{(drawn as that many edges;
no edge when $\alpha\perp\beta$)}.
\end{definition}

It is well defined \added{(independent of the choice of fundamental system)} by the Weyl chamber
argument, since fundamental systems are carried to fundamental systems by the reflections
$\langle s_\alpha\rangle$, \added{which are isometries and so preserve all the angles}.

\begin{definition}\label{L04a:def:dynkin}
The \emph{Dynkin diagram} is the Coxeter graph with arrows on the double and triple edges, pointing to
the shorter vector.
\end{definition}

\begin{remark}
In an irreducible system, the possible lengths are determined by the angles
\added{(by the table, the ratio of the lengths of two non-orthogonal roots is fixed by their angle;
connecting through the diagram fixes all lengths up to one overall scale)}. For a reducible system the
components can be rescaled independently.
\unclear{Margin note by $A_1\times A_1$, probably ``unrelated scale!''.}
\end{remark}

\begin{example}
\leavevmode
\begin{center}
\begin{tabular}{r@{\quad$\longrightarrow$\quad}l}
$A_1$ & \tikz[baseline=-3pt]{\LIVv{0,0}} \\[4pt]
$A_1\times A_1$ & \tikz[baseline=-3pt]{\LIVv{0,0}\LIVv{1,0}} \\[4pt]
$A_2$ & \tikz[baseline=-3pt]{\LIVe{0,0}{1,0}\LIVv{0,0}\LIVv{1,0}} \\[4pt]
$B_2$ & \tikz[baseline=-3pt]{\LIVdR{0}\LIVv{0,0}\LIVv{1,0}} \\[4pt]
$G_2$ & \tikz[baseline=-3pt]{\LIVtR{0}\LIVv{0,0}\LIVv{1,0}}
\end{tabular}
\end{center}
\end{example}

\orig{End of Lecture 4 (p.\,15).}

%======================================================================
\chapter[Lecture 5: Root systems (continued)]{Lecture 5: Classification of crystallographic root systems (continued)}
\label{L04a:ch:L5}

\paragraph{Recall.} A root system $\Delta\subset\R^n\setminus\{0\}$ satisfies
\begin{enumerate}[label=\arabic*)]
\item $\Delta$ is finite and spans $\R^n$;
\item $\alpha\in\Delta\implies(n\alpha\in\Delta\iff n=\pm1)$;
\item $\Delta$ is invariant under the reflections in the hyperplanes $\alpha^\perp$:
  $s_\alpha(\beta)\in\Delta$ for all $\alpha,\beta\in\Delta$;
\item $n_{\beta\alpha}:=2\ip{\beta}{\alpha}/\ip{\alpha}{\alpha}\in\Z$,
\end{enumerate}
and 4) implies $4\cos^2\theta_{\alpha\beta}\in\Z$, really $\in\{0,1,2,3,4\}$.

The plan: root lattices $\to$ root systems $\longleftrightarrow$ simple roots (last time)
$\longleftrightarrow$ Dynkin diagrams (today).

From \S\ref{L04a:sec:4.4}: the \emph{Coxeter graph} has a vertex for each simple root and an edge of
weight $n_{\alpha\beta}\cdot n_{\beta\alpha}$ between $\alpha$ and $\beta$. The \emph{Dynkin diagram} is
the Coxeter graph with arrows on the multi-edges pointing at the shorter root.
\orig{Page has ``(longer)'' with ``shorter'' written below.}

%----------------------------------------------------------------------
\section{Classification of Coxeter graphs}
\label{L04a:sec:5.1}

We ignore lengths.

\begin{definition}\label{L04a:def:admissible}
A linearly independent set of $n$ unit vectors $\{v_1,\dots,v_n\}$ spanning $\R^n$ is \emph{admissible}
if for all $i\neq j$
\[
  \ip{v_i}{v_j}\le0 \qquad\text{and}\qquad 4\ip{v_i}{v_j}^2 = 4\cos^2\theta_{ij}\in\{0,1,2,3\}.
\]
An \emph{admissible diagram} is the Coxeter graph of an admissible set
\added{(vertices $v_i$, with $4\ip{v_i}{v_j}^2$ edges between $v_i$ and $v_j$)}.
\end{definition}

\begin{center}
\fbox{\parbox{0.8\linewidth}{\centering A normalized set of simple roots is admissible.\\
(Note: the angles are the determining factor.)}}
\end{center}
\added{This follows from Theorem~\ref{L04a:thm:basis}, Lemma~\ref{L04a:lem:obtuse} and the table.}

The simple roots of an irreducible $\Delta$ are not all orthogonal
\added{--- more precisely, they cannot be split into two mutually orthogonal sets
(Proposition~\ref{L04a:prop:decomp})} --- so the Coxeter graph is connected. So we only need to classify
connected admissible diagrams.

\begin{theorem}\label{L04a:thm:classification}
The Dynkin diagram of an irreducible (crystallographic) root system is of one of the following types
($n$ = number of vertices).

\smallskip
\noindent\emph{Infinite families:}
\begin{center}
\begin{tabular}{lll}
$A_n$ & \tikz[baseline=-3pt]{\LIVe{0,0}{1,0}\draw[dotted] (1,0)--(2,0);\LIVe{2,0}{3,0}\foreach\x in{0,1,2,3}{\LIVv{\x,0}}} & $n\ge1$\\[6pt]
$B_n$ & \tikz[baseline=-3pt]{\LIVe{0,0}{1,0}\draw[dotted] (1,0)--(2,0);\LIVdR{2}\foreach\x in{0,1,2,3}{\LIVv{\x,0}}} & $n\ge2$\\[6pt]
$C_n$ & \tikz[baseline=-3pt]{\LIVe{0,0}{1,0}\draw[dotted] (1,0)--(2,0);\LIVdL{2}\foreach\x in{0,1,2,3}{\LIVv{\x,0}}} & $n\ge3$\\[6pt]
$D_n$ & \tikz[baseline=-3pt]{\LIVe{0,0}{1,0}\draw[dotted] (1,0)--(2,0);\LIVe{2,0}{2.8,0.45}\LIVe{2,0}{2.8,-0.45}\foreach\x in{0,1,2}{\LIVv{\x,0}}\LIVv{2.8,0.45}\LIVv{2.8,-0.45}} & $n\ge4$
\end{tabular}
\end{center}

\noindent\emph{Exceptional types:}
\begin{center}
\begin{tabular}{ll}
$E_6$ & \tikz[baseline=-3pt]{\foreach\x in{0,1,2,3}{\LIVe{\x,0}{\x+1,0}}\LIVe{2,0}{2,0.7}\foreach\x in{0,...,4}{\LIVv{\x,0}}\LIVv{2,0.7}}\\[12pt]
$E_7$ & \tikz[baseline=-3pt]{\foreach\x in{0,...,4}{\LIVe{\x,0}{\x+1,0}}\LIVe{2,0}{2,0.7}\foreach\x in{0,...,5}{\LIVv{\x,0}}\LIVv{2,0.7}}\\[12pt]
$E_8$ & \tikz[baseline=-3pt]{\foreach\x in{0,...,5}{\LIVe{\x,0}{\x+1,0}}\LIVe{2,0}{2,0.7}\foreach\x in{0,...,6}{\LIVv{\x,0}}\LIVv{2,0.7}}\\[6pt]
$F_4$ & \tikz[baseline=-3pt]{\LIVe{0,0}{1,0}\LIVdR{1}\LIVe{2,0}{3,0}\foreach\x in{0,...,3}{\LIVv{\x,0}}}\\[6pt]
$G_2$ & \tikz[baseline=-3pt]{\LIVtR{0}\LIVv{0,0}\LIVv{1,0}}
\end{tabular}
\end{center}
\added{Conversely, each of these diagrams is the Dynkin diagram of an irreducible root system, which is
unique up to isomorphism.}
\unclear{Parenthetical under $D_n$ reads roughly ``(each also unique)''.}
\added{(The ranges of $n$ avoid repetitions: $B_2=C_2$, $D_3=A_3$.)}
\end{theorem}

It will suffice to classify admissible diagrams. By definition, subsets of the vectors of an admissible
set satisfy the admissibility condition inside their span (the resulting diagram could be
disconnected); \added{that is, any subdiagram of an admissible diagram, obtained by deleting vertices,
is admissible.}

\begin{lemma}\label{L04a:lem:tree}
A connected admissible diagram is a tree.
\end{lemma}

\begin{proof}
Consider $v=\sum_{i=1}^n v_i$, where $\{v_i\}$ is admissible. The $v_i$ are linearly independent, so
$v\neq0$, and
\[
  0<\ip{v}{v} = \sum_i\ip{v_i}{v_i} + \sum_{i<j}2\ip{v_i}{v_j} = n+\sum_{i<j}2\ip{v_i}{v_j}.
\]
If $v_i$ is connected to $v_j$ in the Coxeter graph, then by the table
$2\ip{v_i}{v_j}\in\{-1,-\sqrt2,-\sqrt3\}$\added{, in particular $2\ip{v_i}{v_j}\le-1$}. Hence
$\sum_{i<j}2\ip{v_i}{v_j}$ has at most $n-1$ terms not equal to $0$, i.e.\ there are at most $n-1$
connected pairs of vertices. But by assumption the diagram is connected, so there are exactly $n-1$
pairs of connected vertices \added{(a connected graph on $n$ vertices has at least $n-1$ edges)}, and
it is a tree (maybe with multiple edges, but not too many).
\orig{a scribbled-out remark at the bottom of p.\,20}
\end{proof}

\begin{lemma}\label{L04a:lem:degree}
The degree of each vertex, counted with multiplicity, is at most $3$.
\end{lemma}

\begin{proof}
Fix a vertex $v$, connected to the vertices $v_1,\dots,v_k$.
\fixed{written ``fix vertex $v_1$'', then used as $v$}
Since the diagram is a tree \added{(Lemma~\ref{L04a:lem:tree}, applied to the subdiagram on
$v,v_1,\dots,v_k$)}, $\ip{v_i}{v_j}=0$ for $i\neq j$, so $\{v_1,\dots,v_k\}$ is orthonormal. Also,
$\{v_1,\dots,v_k,v\}$ is linearly independent (by admissibility; for simple roots this is Theorem~\ref{L04a:thm:basis}), so
the normalized projection
\[
  v_0 := \frac{\operatorname{proj}_{\langle v_1,\dots,v_k\rangle^\perp}(v)}{\norm{\operatorname{proj}_{\langle v_1,\dots,v_k\rangle^\perp}(v)}}
\]
is defined, and $\ip{v}{v_0}\neq0$. Now $\{v_0,\dots,v_k\}$ is an orthonormal set, and $v$ lies in its
span, so
\[
  v=\sum_{i=0}^k\ip{v}{v_i}v_i \qquad\text{and}\qquad \ip{v}{v}=\sum_{i=0}^k\ip{v}{v_i}^2=1 .
\]
Since $\ip{v}{v_0}^2\neq0$,
\[
  \sum_{i=1}^k 4\ip{v}{v_i}^2<4 .
\]
But $4\cos^2\theta = 4\ip{v}{v_i}^2$ is the number of edges from $v$ to $v_i$. So the degree of $v$,
counted with multiplicity, is at most $3$.
\end{proof}

\begin{corollary}\label{L04a:cor:G2}
\tikz[baseline=-3pt]{\LIVtN{0}\LIVv{0,0}\LIVv{1,0}} is the only connected admissible diagram with a
triple edge.
\end{corollary}

\begin{proof}[\added{Proof}]
\added{If $v$ and $w$ are joined by a triple edge, both already have degree $3$, so by
Lemma~\ref{L04a:lem:degree} neither is joined to any other vertex; by connectedness there are no
other vertices.}
\end{proof}

So from now on we consider only double and single edges.

\begin{definition}\label{L04a:def:chain}
A \emph{simple chain} is a collection of vertices connected by single edges:
\begin{center}
\begin{tikzpicture}
  \LIVe{0,0}{1,0}\draw[dotted] (1,0)--(2,0);\LIVe{2,0}{3,0}
  \foreach\x in{0,1,2,3}{\LIVv{\x,0}}
  \node[below] at (0,-0.1) {\small$v_1$};
  \node[below] at (1,-0.1) {\small$v_2$};
  \node[below] at (2,-0.1) {\small$v_{k-1}$};
  \node[below] at (3,-0.1) {\small$v_k$};
\end{tikzpicture}
\end{center}
\end{definition}

\begin{lemma}[Simple chain collapse]\label{L04a:lem:collapse}
In an admissible diagram, a simple chain represented by $\{v_1,\dots,v_k\}$ can be replaced by the single
vector $v=\sum_{i=1}^k v_i$, giving an admissible diagram.
\end{lemma}

\begin{proof}
We want to show that $v$ is a unit vector and that the collapsed diagram is admissible.
\added{(Linear independence of the new set is clear, since the $v_i$ and the remaining vectors are
linearly independent.)}
We have
\[
  \ip{v}{v} = k + \sum_{i<j}2\ip{v_i}{v_j},
\]
and $\ip{v_i}{v_j}=0$ for all $i<j$ except $j=i+1$ \added{(a tree has no other edges among the $v_i$)}.
For a single edge, recall $2\ip{v_i}{v_{i+1}}=-1$. Hence
\[
  \ip{v}{v} = k-(k-1) = 1 .
\]
Now consider a vertex $u$ not represented in the chain. It connects to at most one vertex in the chain
(by the tree property), say $v_j$ \added{(if none, $\ip{u}{v}=0$)}. Then
\[
  \ip{u}{v}=\sum_{i=1}^k\ip{u}{v_i}=\ip{u}{v_j},
\]
so the angles $\theta_{u,v}$ and $\theta_{u,v_j}$ are the same, and the collapsed set is admissible.
\orig{scribbled-out ``recall admissible\dots'' at the foot of p.\,23}
\end{proof}

% ============================================================
% MAT670_46, pages 24--46.
% Pages 24--31: continuation of Lecture 5 (classification of Coxeter graphs).
% Pages 32--46: Lecture 6 (root systems -> root lattices; codes).
% ============================================================
\tikzset{L04bv/.style={circle,draw,fill=white,inner sep=0pt,minimum size=6pt},
         L04bdbl/.style={double,double distance=2pt},
         L04blab/.style={font=\scriptsize}}

\section{Classification of Coxeter graphs (continued)}
\label{L04b:sec:coxeter-cont}
\orig{continues Lecture 5 from p.~23}

Recall from the previous pages: a connected admissible diagram is a tree; each vertex
meets at most three edges counted with multiplicity; the only admissible diagram with a
triple edge is \tikz[baseline=-3pt]{\node[L04bv] (a) at (0,0) {}; \node[L04bv] (b) at (0.6,0) {};
\draw (a.north east)--(b.north west); \draw (a)--(b); \draw (a.south east)--(b.south west);}
($G_2$); and a simple chain $\{v_1,\dots,v_k\}$ may be collapsed to the single unit vector
$v=\sum_i v_i$, leaving an admissible diagram. From now on only single and double edges occur.

\subsection{Forbidden configurations}

\begin{proposition}\label{L04b:prop:forbidden}
Coxeter graphs of admissible diagrams containing any of the following subgraphs are forbidden:
\begin{center}
\begin{tikzpicture}[scale=0.8]
  % two double edges
  \foreach \i in {0,1,2,3,4} \node[L04bv] (a\i) at (\i,0) {};
  \draw[L04bdbl] (a0)--(a1); \draw (a1)--(a2); \draw[dotted,thick] (a2)--(a3);
  \draw[L04bdbl] (a3)--(a4);
  % double edge + fork
  \begin{scope}[xshift=6cm]
  \foreach \i in {0,1,2,3} \node[L04bv] (b\i) at (\i,0) {};
  \node[L04bv] (b4) at (3.8,0.5) {}; \node[L04bv] (b5) at (3.8,-0.5) {};
  \draw[L04bdbl] (b0)--(b1); \draw (b1)--(b2); \draw[dotted,thick] (b2)--(b3);
  \draw (b3)--(b4); \draw (b3)--(b5);
  \end{scope}
  % two forks
  \begin{scope}[xshift=11.5cm]
  \node[L04bv] (c0) at (-0.8,0.5) {}; \node[L04bv] (c1) at (-0.8,-0.5) {};
  \foreach \i in {0,1,2} \node[L04bv] (d\i) at (\i,0) {};
  \node[L04bv] (c4) at (2.8,0.5) {}; \node[L04bv] (c5) at (2.8,-0.5) {};
  \draw (c0)--(d0); \draw (c1)--(d0); \draw (d0)--(d1); \draw[dotted,thick] (d1)--(d2);
  \draw (d2)--(c4); \draw (d2)--(c5);
  \end{scope}
\end{tikzpicture}
\end{center}
Hence a connected admissible diagram (with no triple edge) contains at most one branch
point \emph{or} at most one double edge, but not both.
\end{proposition}

\begin{proof}
\added{In each case collapse the simple chain joining the two special features (possibly of
length zero) to a single vertex. By the collapsing lemma the result is again admissible, but the
new vertex meets edges of total multiplicity $2+2$, $2+1+1$, or $1+1+1+1$, i.e.\ $4$, contradicting
the bound of $3$ edges (with multiplicity) at each vertex. The same argument rules out two double
edges, two branch points, or a double edge and a branch point anywhere in the tree, since they are
always joined by a simple chain.}
\end{proof}

So, besides $G_2$, the only possible connected Coxeter graphs are:
\begin{center}
\begin{tikzpicture}[scale=0.85]
  % chain with one double edge
  \foreach \x/\n/\l in {0/u1/u_1,1/u2/u_2,2.6/up/u_p,3.6/vq/v_q,5.2/v2/v_2,6.2/v1/v_1}
     {\node[L04bv] (\n) at (\x,0) {}; \node[L04blab,below=2pt] at (\x,0) {$\l$};}
  \draw (u1)--(u2); \draw[dotted,thick] (u2)--(up); \draw[L04bdbl] (up)--(vq);
  \draw[dotted,thick] (vq)--(v2); \draw (v2)--(v1);
  % branched
  \begin{scope}[xshift=8.3cm]
  \foreach \x/\n/\l in {0/a1/u_1,1/a2/u_2,2.4/ap/u_{p-1}}
     {\node[L04bv] (\n) at (\x,0) {}; \node[L04blab,below=2pt] at (\x,0) {$\l$};}
  \node[L04bv] (x) at (3.3,0) {}; \node[L04blab,below=2pt] at (3.3,0) {$x$};
  \node[L04bv] (bq) at (4.1,0.6) {}; \node[L04blab,above=2pt] at (4.1,0.6) {$v_{q-1}$};
  \node[L04bv] (b1) at (5.5,0.6) {}; \node[L04blab,above=2pt] at (5.5,0.6) {$v_1$};
  \node[L04bv] (cr) at (4.1,-0.6) {}; \node[L04blab,below=2pt] at (4.1,-0.6) {$w_{r-1}$};
  \node[L04bv] (c1) at (5.5,-0.6) {}; \node[L04blab,below=2pt] at (5.5,-0.6) {$w_1$};
  \draw (a1)--(a2); \draw[dotted,thick] (a2)--(ap); \draw (ap)--(x);
  \draw (x)--(bq); \draw[dotted,thick] (bq)--(b1);
  \draw (x)--(cr); \draw[dotted,thick] (cr)--(c1);
  \end{scope}
\end{tikzpicture}
\end{center}
together with the simple chains \tikz[baseline=-3pt]{\foreach \i/\x in {1/0,2/0.5,3/1.3,4/1.8} \node[L04bv] (s\i) at (\x,0) {};
\draw (s1)--(s2); \draw[dotted,thick] (s2)--(s3); \draw (s3)--(s4);}, which give $A_n$.

\subsection{One double edge: \texorpdfstring{$B_n$, $C_n$, $F_4$}{Bn, Cn, F4}}

Label the diagram with one double edge as above: $u_1,\dots,u_p$, then the double edge
$u_p = v_q$, then $v_q,\dots,v_1$. Put
\[
  u=\sum_{i=1}^p i\,u_i,\qquad v=\sum_{j=1}^q j\,v_j .
\]
Since the $u_i$ form a simple chain, $2\ip{u_i}{u_{i+1}}=-1$ for $1\le i\le p-1$ and all other pairs
are orthogonal. Hence
\begin{align*}
  \ip{u}{u} &= \sum_{i=1}^p i^2\ip{u_i}{u_i} + \sum_{i<j} 2ij\,\ip{u_i}{u_j}
   = \sum_{i=1}^p i^2 + \sum_{i=1}^{p-1} i(i+1)\cdot\underbrace{2\ip{u_i}{u_{i+1}}}_{=-1}\\
  &= \sum_{i=1}^p i^2 - \sum_{i=1}^{p-1} i(i+1) = p^2 - \sum_{i=1}^{p-1} i = \frac{p(p+1)}{2},
\end{align*}
and similarly $\ip{v}{v} = q(q+1)/2$.

Since the double edge $u_p = v_q$ is the only non-orthogonal relation between the two chains,
$\ip{u}{v} = pq\,\ip{u_p}{v_q}$, and $4\ip{u_p}{v_q}^2=2$. Thus
\[
  \ip{u}{v}^2 = \frac{p^2q^2}{2}.
\]
Since $u$ and $v$ are not parallel \added{(the $u_i, v_j$ are linearly independent)}, the
Cauchy--Schwarz inequality is strict: $\ip{u}{v}^2 < \ip{u}{u}\ip{v}{v}$, i.e.
\[
  \frac{p^2q^2}{2} < \frac{p(p+1)}{2}\cdot\frac{q(q+1)}{2},
  \qquad\text{so}\qquad 2pq < (p+1)(q+1),
\]
which, for integers $p,q\ge 1$, is
\[
  \boxed{(p-1)(q-1) < 2.}
\]
Hence either $p=q=2$, which is $F_4$, or $p=1$ (or $q=1$) with the other arbitrary, which is
$B_n$ and $C_n$ \added{($n=p+q$; the Coxeter graphs of $B_n$ and $C_n$ coincide, and they are
distinguished only by the direction of the arrow in the Dynkin diagram)}.

\subsection{One branch point: \texorpdfstring{$D_n$, $E_6$, $E_7$, $E_8$}{Dn, E6, E7, E8}}

Finally, take the branched diagram above, with branch vertex $x$ and three arms
$u_1,\dots,u_{p-1}$, $v_1,\dots,v_{q-1}$, $w_1,\dots,w_{r-1}$ ($p,q,r\ge 2$), the indices
increasing towards $x$. As before put
\[
  u=\sum_{i=1}^{p-1} i\,u_i,\qquad v=\sum_{i=1}^{q-1} i\,v_i,\qquad w=\sum_{i=1}^{r-1} i\,w_i .
\]
Note that $u,v,w$ are mutually orthogonal vectors (no edges between different arms), and
$x$ is not in their span $\langle u,v,w\rangle$ \added{(by linear independence of all the vertices)}.
By the previous computation with $p$ replaced by $p-1$,
$\ip{u}{u} = p(p-1)/2$, and similarly for $v,w$.

Since $x$ is a unit vector not in $\langle u,v,w\rangle$, and $u/\norm{u},v/\norm{v},w/\norm{w}$ is an
orthonormal basis of that space, the squared length of the projection of $x$ onto it is strictly
less than $1$:
\[
  1 = \ip{x}{x} > \Big\langle x,\frac{u}{\norm{u}}\Big\rangle^2 + \Big\langle x,\frac{v}{\norm{v}}\Big\rangle^2
  + \Big\langle x,\frac{w}{\norm{w}}\Big\rangle^2 .
\]
Also $\ip{x}{u_i}=0$ unless $i=p-1$, and $4\ip{x}{u_{p-1}}^2 = 1$ (single edge). Then
\[
  \ip{x}{u}^2 = \Big(\sum_{i=1}^{p-1} i\,\ip{x}{u_i}\Big)^2 = (p-1)^2\ip{x}{u_{p-1}}^2 = \frac{(p-1)^2}{4}
\]
\fixed{notes: $\sum i^2\langle x,u_i\rangle$, square misplaced}%
and therefore
\[
  \Big\langle x,\frac{u}{\norm{u}}\Big\rangle^2 = \frac{(p-1)^2}{4}\cdot\frac{2}{p(p-1)}
   = \frac12\Big(1-\frac1p\Big).
\]
The same holds for $v$ (with $q$) and $w$ (with $r$). Then
\[
  2 > \Big(1-\frac1p\Big) + \Big(1-\frac1q\Big) + \Big(1-\frac1r\Big)
  \quad\Longrightarrow\quad \frac1p + \frac1q + \frac1r > 1,
\]
where $p,q,r\ge 2$ are integers. Assume $p\ge q\ge r\ge 2$. Then $3/r > 1$ forces $r=2$, and:
\begin{itemize}
  \item $q=2$: any $p\ge 2$ works \added{--- this is $D_{p+2}$};
  \item $q=3$: $\frac1q+\frac1r = \frac56$, so $\frac1p>\frac16$, i.e.\ $3\le p<6$
        \added{--- $p=3,4,5$ give $E_6,E_7,E_8$};
  \item $q\ge4$: \added{$\frac1p+\frac1q+\frac12 \le \frac14+\frac14+\frac12=1$}, so there are no solutions.
\end{itemize}
\added{(The diagram has $p+q+r-2$ vertices, which gives the subscripts.)}

\begin{addedblock}
\textbf{Summary.} The connected Coxeter graphs of admissible diagrams (equivalently, of irreducible
crystallographic root systems) are
$A_n$ ($n\ge1$), $B_n=C_n$ as graphs ($n\ge2$), $D_n$ ($n\ge4$), $E_6$, $E_7$, $E_8$, $F_4$, $G_2$.
\begin{center}
\begin{tikzpicture}[scale=0.62, every node/.style={font=\small}]
  % A_n
  \node[anchor=east] at (-0.4,0) {$A_n$};
  \foreach \i in {0,1,3,4} \node[L04bv] (A\i) at (\i,0) {};
  \draw (A0)--(A1); \draw[dotted,thick] (A1)--(A3); \draw (A3)--(A4);
  % B_n / C_n
  \node[anchor=east] at (-0.4,-1.2) {$B_n,C_n$};
  \foreach \i in {0,1,3,4} \node[L04bv] (B\i) at (\i,-1.2) {};
  \draw (B0)--(B1); \draw[dotted,thick] (B1)--(B3); \draw[L04bdbl] (B3)--(B4);
  % D_n
  \node[anchor=east] at (-0.4,-2.6) {$D_n$};
  \foreach \i in {0,1,3} \node[L04bv] (D\i) at (\i,-2.6) {};
  \node[L04bv] (Du) at (3.8,-2.1) {}; \node[L04bv] (Dd) at (3.8,-3.1) {};
  \draw (D0)--(D1); \draw[dotted,thick] (D1)--(D3); \draw (D3)--(Du); \draw (D3)--(Dd);
  % F4, G2
  \node[anchor=east] at (-0.4,-4.0) {$F_4$};
  \foreach \i in {0,1,2,3} \node[L04bv] (F\i) at (\i,-4.0) {};
  \draw (F0)--(F1); \draw[L04bdbl] (F1)--(F2); \draw (F2)--(F3);
  \node[anchor=east] at (-0.4,-5.2) {$G_2$};
  \node[L04bv] (G0) at (0,-5.2) {}; \node[L04bv] (G1) at (1,-5.2) {};
  \draw (G0)--(G1); \draw ([yshift=2pt]G0.east)--([yshift=2pt]G1.west);
  \draw ([yshift=-2pt]G0.east)--([yshift=-2pt]G1.west);
  % E6, E7, E8
  \begin{scope}[xshift=9cm]
  \foreach \n/\k/\y in {E_6/5/0, E_7/6/-2, E_8/7/-4} {
    \node[anchor=east] at (-0.4,\y) {$\n$};
    \foreach \i in {1,...,\k} \node[L04bv] (e\i) at (\i-1,\y) {};
    \foreach \i [evaluate=\i as \j using int(\i-1)] in {2,...,\k} \draw (e\i)--(e\j);
    \node[L04bv] (top) at (2,\y+0.9) {}; \draw (e3)--(top);
  }
  \end{scope}
\end{tikzpicture}
\end{center}
\end{addedblock}

\begin{remark}
It remains to show that each of
$A_n$, $B_n$, $C_n$, $D_n$, $E_6$, $E_7$, $E_8$, $F_4$, $G_2$ actually exists, i.e.\ to construct
the corresponding root systems. This is done in the next lecture; $G_2$ is left as an exercise
\added{(its $12$ roots were drawn in Lecture~4: check axioms 1)--4) directly)}.
\orig{p.~31: ``Existence of $A_n,\dots,F_4$: construct next time; $G_2$ -- ex.''}
\end{remark}

% ============================================================
\chapter{Lecture 6: Root systems and root lattices}
\label{L04b:ch:lecture6}

\section{Construction of crystallographic root systems}
\label{L04b:sec:construction}

\subsection*{\texorpdfstring{$A_n$}{An}}
In $A_n$ the simple roots are pairwise orthogonal or at angle $2\pi/3$.

Consider $x\in\Z^{n+1}$ such that $\sum_{i=1}^{n+1} x_i = 0$. These form an $n$-dimensional
\added{sublattice} $S$ of $\Z^{n+1}$. Let
\[
  \Delta = \{x\in S : \norm{x}=\sqrt2\}.
\]
That is, all but two coordinates are $0$, one coordinate is $+1$ and one is $-1$: the vectors
are the permutations of $(1,-1,0,\dots,0)$, so $|\Delta| = n(n+1)$. Choose the simple roots
\[
  \alpha_i = e_i - e_{i+1},\qquad i=1,\dots,n,
\]
\orig{diagram labelled $e_{i+1}-e_i$}
with matrix of coordinates and Dynkin diagram
\[
  \begin{bmatrix} 1 & -1 & & & \\ & 1 & -1 & & \\ & & \ddots & \ddots & \\ & & & 1 & -1\end{bmatrix}
  \qquad
  \begin{tikzpicture}[baseline=-3pt,scale=0.9]
    \foreach \x/\n/\l in {0/a/e_1-e_2,1.4/b/e_2-e_3,3.4/c/e_n-e_{n+1}}
      {\node[L04bv] (\n) at (\x,0) {}; \node[L04blab,below=2pt] at (\x,0) {$\l$};}
    \draw (a)--(b); \draw[dotted,thick] (b)--(c);
  \end{tikzpicture}
\]
To check the Dynkin diagram: the number of edges joining $\alpha_i$ and $\alpha_j$ is
$4\ip{\alpha_i}{\alpha_j}^2/(\norm{\alpha_i}^2\norm{\alpha_j}^2) = \ip{\alpha_i}{\alpha_j}^2\in\{1,0\}$,
and it is $1$ exactly when $j=i\pm1$.

Note that these roots generate the full sublattice $\{x\in\Z^{n+1}: \sum x_i=0\}$
\added{(any such $x$ equals $\sum_{k=1}^{n} (x_1+\dots+x_k)\,\alpha_k$)}. So the $A_n$ root lattice
is all of $S$. This description is maybe not so useful for visualizing, since $S$ sits inside
$\R^{n+1}$.

\begin{example}[$A_2$]
Simple roots $(-1,1,0)$ and $(0,-1,1)$, whose sum is $(-1,0,1)$; together with the negatives
$(1,-1,0)$, $(0,1,-1)$, $(1,0,-1)$ these are the six roots, lying in the plane
$x_1+x_2+x_3=0$. In that plane they form the familiar hexagonal configuration:
\begin{center}
\begin{tikzpicture}[scale=1.1]
  \foreach \a in {0,60,...,300} \draw[-{Stealth}] (0,0)--(\a:1);
  \draw[-{Stealth},very thick,addcol] (0,0)--(0:1) node[right,font=\scriptsize] {$\alpha_1$};
  \draw[-{Stealth},very thick,addcol] (0,0)--(120:1) node[above left,font=\scriptsize] {$\alpha_2$};
\end{tikzpicture}
\end{center}
\orig{p.~33: pink star in the plane $\sum x_i=0$}
\end{example}

\begin{example}[$A_3$]
The $12$ roots of $A_3$ are the $12$ nearest neighbours of the origin in the FCC lattice
(``the FCC system'') \added{--- the $A_3$ root lattice is the face-centred cubic lattice, with
kissing number $12$}.
\end{example}

\subsection*{\texorpdfstring{$B_n$}{Bn}}
$\Delta$ is the set of all integer vectors $x\in\Z^n$ of norm $1$ or norm $\sqrt2$, i.e.\
$\pm e_i$ and $\pm e_i\pm e_j$. So $|\Delta| = 2n + 4\binom n2 = 2n^2$. Simple roots:
\[
  \alpha_i = e_i - e_{i+1}\ (1\le i\le n-1),\qquad \text{short root } \alpha_n = e_n,
\]
\[
  \begin{bmatrix} 1 & -1 & & & \\ & 1 & -1 & & \\ & & \ddots & \ddots & \\ & & & 1 & -1\\ & & & & 1\end{bmatrix}
  \qquad
  \begin{tikzpicture}[baseline=-3pt,scale=0.9]
    \foreach \x/\n/\l in {0/a/e_1-e_2,1.4/b/e_2-e_3,3.2/c/e_{n-1}-e_n,4.6/d/e_n}
      {\node[L04bv] (\n) at (\x,0) {}; \node[L04blab,below=2pt] at (\x,0) {$\l$};}
    \draw (a)--(b); \draw[dotted,thick] (b)--(c); \draw[L04bdbl] (c)--(d);
    \draw[thick] (3.8,0.12)--(3.98,0)--(3.8,-0.12);
  \end{tikzpicture}
\]
Sadly, from this construction the associated root lattice is all of $\Z^n$ \added{(it contains
every $e_i$)}, so we get nothing new.

\subsection*{\texorpdfstring{$C_n$}{Cn}}
$\Delta$ is the set of all integer vectors of the form $2x$ with $x\in\Z^n$ of length $1$, or
$x\in\Z^n$ of length $\sqrt2$; i.e.\ $\pm 2e_i$ and $\pm e_i\pm e_j$. Again
$|\Delta| = 2n + 4\binom n2 = 2n^2$. Simple roots:
\[
  \alpha_i = e_i - e_{i+1}\ (1\le i\le n-1),\qquad \text{long root } \alpha_n = 2e_n,
\]
\[
  \begin{bmatrix} 1 & -1 & & & \\ & 1 & -1 & & \\ & & \ddots & \ddots & \\ & & & 1 & -1\\ & & & & 2\end{bmatrix}
  \qquad
  \begin{tikzpicture}[baseline=-3pt,scale=0.9]
    \foreach \x/\n/\l in {0/a/e_1-e_2,1.4/b/e_2-e_3,3.2/c/e_{n-1}-e_n,4.6/d/2e_n}
      {\node[L04bv] (\n) at (\x,0) {}; \node[L04blab,below=2pt] at (\x,0) {$\l$};}
    \draw (a)--(b); \draw[dotted,thick] (b)--(c); \draw[L04bdbl] (c)--(d);
    \draw[thick] (4.0,0.12)--(3.82,0)--(4.0,-0.12);
  \end{tikzpicture}
\]
\fixed{diagram drawn with a single last edge}
So the $C_n$ root lattice is the set of all integer vectors with \emph{even} coordinate sum.

\begin{example}[$B_2$ and $C_2$]
The two root systems, with the simple roots highlighted:
\begin{center}
\begin{tikzpicture}[scale=0.9]
  \foreach \v in {(1,0),(0,1),(-1,0),(0,-1),(1,1),(1,-1),(-1,1),(-1,-1)} \draw[-{Stealth},gray] (0,0)--\v;
  \draw[-{Stealth},very thick,addcol] (0,0)--(1,-1) node[below right,font=\scriptsize] {$\alpha_1=e_1-e_2$};
  \draw[-{Stealth},very thick,addcol] (0,0)--(0,1) node[above,font=\scriptsize] {$\alpha_2=e_2$};
  \node at (0,-1.7) {$B_2$};
  \begin{scope}[xshift=5cm]
  \foreach \v in {(2,0),(0,2),(-2,0),(0,-2),(1,1),(1,-1),(-1,1),(-1,-1)} \draw[-{Stealth},gray] (0,0)--\v;
  \draw[-{Stealth},very thick,addcol] (0,0)--(1,-1) node[below right,font=\scriptsize] {$\alpha_1=e_1-e_2$};
  \draw[-{Stealth},very thick,addcol] (0,0)--(0,2) node[above,font=\scriptsize] {$\alpha_2=2e_2$};
  \node at (0,-2.4) {$C_2$};
  \end{scope}
\end{tikzpicture}
\end{center}
\orig{redrawn; sketch showed only $\alpha_1,\alpha_2$}
\added{In both, the simple roots meet at angle $3\pi/4$; $C_2$ is $B_2$ rotated by $\pi/4$ and rescaled.}
\end{example}

\subsection*{\texorpdfstring{$D_n$}{Dn}}
$\Delta$ is the set of all integer vectors $x\in\Z^n$ of length $\sqrt2$, i.e.\ $\pm e_i\pm e_j$, so
\[
  |\Delta| = 4\binom n2 = 2n(n-1).
\]
\fixed{notes: $2\binom n2$}
Simple roots:
\[
  \alpha_i = e_i - e_{i+1}\ (1\le i\le n-1)\qquad\text{and}\qquad \alpha_n = e_{n-1}+e_n,
\]
\[
  \begin{bmatrix} 1 & -1 & & & \\ & \ddots & \ddots & & \\ & & 1 & -1 & \\ & & & 1 & -1\\ & & & 1 & 1\end{bmatrix}
  \qquad
  \begin{tikzpicture}[baseline=-3pt,scale=0.9]
    \foreach \x/\n/\l in {0/a/e_1-e_2,1.4/b/e_2-e_3,3.2/c/e_{n-2}-e_{n-1}}
      {\node[L04bv] (\n) at (\x,0) {}; \node[L04blab,below=2pt] at (\x,0) {$\l$};}
    \node[L04bv] (u) at (4.3,0.55) {}; \node[L04blab,right=2pt] at (4.3,0.55) {$e_{n-1}+e_n$};
    \node[L04bv] (d) at (4.3,-0.55) {}; \node[L04blab,right=2pt] at (4.3,-0.55) {$e_{n-1}-e_n$};
    \draw (a)--(b); \draw[dotted,thick] (b)--(c); \draw (c)--(u); \draw (c)--(d);
  \end{tikzpicture}
\]
This also generates the lattice $\{x\in\Z^n : \text{the coordinate sum is even}\}$, the same as
for $C_n$. (I can't draw $D_4$, sadly.)

\subsection*{\texorpdfstring{$F_4$}{F4}}
$\Delta$ is the set of vectors $x\in\R^4$ of length $1$ or $\sqrt2$ such that the coordinates of
$2x$ are all even or all odd integers. \added{Explicitly: $\pm e_i$ ($8$), $\pm e_i\pm e_j$ ($24$) and
$\tfrac12(\pm1,\pm1,\pm1,\pm1)$ ($16$), so $|\Delta|=48$.} Simple roots:
\[
  \alpha_1 = (1,-1,0,0),\quad \alpha_2=(0,1,-1,0),\quad \alpha_3=(0,0,1,0),\quad
  \alpha_4 = -\tfrac12(1,1,1,1),
\]
\begin{center}
  \begin{tikzpicture}[baseline=-3pt,scale=0.9]
    \foreach \x/\n in {0/1,1/2,2/3,3/4}
      {\node[L04bv] (f\n) at (\x,0) {}; \node[L04blab,below=2pt] at (\x,0) {$\alpha_\n$};}
    \draw (f1)--(f2); \draw[L04bdbl] (f2)--(f3); \draw (f3)--(f4);
    \draw[thick] (1.42,0.12)--(1.6,0)--(1.42,-0.12);
  \end{tikzpicture}
\end{center}
(check by hand). \added{Indeed $\ip{\alpha_1}{\alpha_2}=-1$ with both of length $\sqrt2$ (angle $2\pi/3$);
$\ip{\alpha_2}{\alpha_3}=-1$ with lengths $\sqrt2,1$ (angle $3\pi/4$, double edge);
$\ip{\alpha_3}{\alpha_4}=-\frac12$ with both of length $1$ (angle $2\pi/3$); all other pairs are orthogonal.}

The $F_4$ root lattice: all coordinates are integers, or all coordinates are halves of odd integers
--- no mixing. \added{That is, $\Z^4\cup(\Z^4+\tfrac12(1,1,1,1))$, which is $D_4^*$, a scaled copy of $D_4$.}

\subsection*{\texorpdfstring{$G_2$}{G2}}
We already drew this one (the $12$ roots form a ``star of David'', two hexagons of lengths $1$ and
$\sqrt3$). \orig{star sketch, see Lecture~4}
The $G_2$ root lattice coincides with the $A_2$ root lattice \added{(the hexagonal lattice)}.

% ------------------------------------------------------------
\section{The \texorpdfstring{$E_8$}{E8} root system}
\label{L04b:sec:E8}

\begin{definition}\label{L04b:def:E8}
The $E_8$ root system consists of the $|\Delta| = 240$ vectors in $\R^8$ of the forms
\[
  \pm e_i \pm e_j\ (i\neq j) \qquad\text{and}\qquad \tfrac12(\pm e_1 \pm \dots \pm e_8)
  \ \added{\text{with an even number of minus signs}}.
\]
\fixed{added the even-sign condition}
\added{There are $4\binom82=112$ of the first kind and $2^7=128$ of the second.}
\end{definition}

Simple roots:
\[
  \alpha_i = e_i-e_{i+1}\ (1\le i\le 6),\qquad \alpha_7 = e_6+e_7,\qquad \alpha_8 = -\tfrac12(1,\dots,1),
\]
\[
  \begin{bmatrix}
   1&-1&0&0&0&0&0&0\\ &1&-1&0&0&0&0&0\\ &&1&-1&0&0&0&0\\ &&&1&-1&0&0&0\\ &&&&1&-1&0&0\\
   &&&&&1&-1&0\\ &&&&&1&1&0\\ -\frac12&-\frac12&-\frac12&-\frac12&-\frac12&-\frac12&-\frac12&-\frac12
  \end{bmatrix}
  \qquad
  \begin{tikzpicture}[baseline=-3pt,scale=0.8]
    \foreach \i in {1,...,6} {\node[L04bv] (g\i) at (\i-1,0) {}; \node[L04blab,below=2pt] at (\i-1,0) {$\alpha_{\i}$};}
    \node[L04bv] (g7) at (4,1) {}; \node[L04blab,above=2pt] at (4,1) {$\alpha_7$};
    \node[L04bv] (g8) at (5,1) {}; \node[L04blab,above=2pt] at (5,1) {$\alpha_8$};
    \foreach \i/\j in {1/2,2/3,3/4,4/5,5/6,5/7,7/8} \draw (g\i)--(g\j);
  \end{tikzpicture}
\]
\added{(Check: $\ip{\alpha_7}{\alpha_5}=-1$, $\ip{\alpha_7}{\alpha_6}=0$, $\ip{\alpha_8}{\alpha_7}=-1$ and
$\ip{\alpha_8}{\alpha_i}=0$ for $i\le 6$; the arms at $\alpha_5$ give $(p,q,r)=(5,3,2)$.)}

This generates the lattice $E_8$ of all $x\in\R^8$ such that
\begin{enumerate}[label=\arabic*)]
  \item the coordinates are all integers or all half-integers \added{(halves of odd integers)}, and
  \item the sum of the coordinates is even.
\end{enumerate}

This is clearly related to $D_8$. Consider the ``deep holes'' of the $D_n$ lattice, i.e.\ the
half-integer points such as $h=\frac12(1,\dots,1)$. In dimension $8$,
\[
  \norm{\tfrac12(1,\dots,1)} = \sqrt2,
\]
and a typical shortest lattice vector has $\norm{(1,1,0,\dots,0)}=\sqrt2$. Similarly,
\[
  \norm{(1,1,0,\dots,0)-\tfrac12(1,\dots,1)} = \sqrt2 .
\]
So if we have $D_8$, there is another copy $D_8+h$ that fits into the holes of $D_8$ without
decreasing the minimal distance: $E_8 = D_8\cup(D_8+h)$.
\added{(For general $n$, $h$ is at distance $\sqrt n/2$ from $D_n$; this equals the minimal
distance $\sqrt2$ exactly when $n=8$.)}

\subsection*{\texorpdfstring{$E_7$, $E_6$}{E7, E6}}
Let $v=(1,1,\dots,1)$ and $w=-e_7-e_8$. Then
\[
  E_7 = E_8\cap v^\perp,\qquad E_6 = E_8\cap \langle w,v\rangle^\perp .
\]
\added{Here $\frac12 v=-\alpha_8$ and $w$ are roots with $\ip{\tfrac12 v}{w}=-1$, so they span an
$A_2$; thus $E_7$ and $E_6$ are the orthogonal complements of an $A_1$ and an $A_2$ in $E_8$.}

These descriptions are not so satisfying, but perhaps there is an easy-to-see construction given
by collapsing simple chains: replace $\alpha_1,\alpha_2$ by $\alpha_1+\alpha_2$, then
$\alpha_1+\alpha_2,\alpha_3$ by $\alpha_1+\alpha_2+\alpha_3$, and so on.
\begin{center}
\begin{tikzpicture}[scale=0.85]
  \foreach \x/\n in {0/1,1/2,2/3,3/4,4/5,5/7,6/8}
    {\node[L04bv] (h\n) at (\x,0) {}; \node[L04blab,below=2pt] at (\x,0) {$\alpha_\n$};}
  \node[L04bv] (h6) at (4,1) {}; \node[L04blab,above=2pt] at (4,1) {$\alpha_6$};
  \foreach \i/\j in {1/2,2/3,3/4,4/5,5/7,7/8,5/6} \draw (h\i)--(h\j);
  \draw (0,0.35)--(0,0.55)--(1,0.55)--(1,0.35);
  \draw (0,0.65)--(0,0.85)--(2,0.85)--(2,0.65);
\end{tikzpicture}
\end{center}
By the collapsing lemma, $\alpha_1+\alpha_2,\alpha_3,\dots,\alpha_8$ is a simple system of type $E_7$,
and $\alpha_1+\alpha_2+\alpha_3,\alpha_4,\dots,\alpha_8$ one of type $E_6$, \added{so the lattices they
generate are copies of the $E_7$, $E_6$ root lattices inside $E_8$.} This continues
\added{(collapsing suitable chains)} to
\[
  D_5,\quad D_4,\quad A_3,\quad A_2,\quad A_1 .
\]

\begin{exercise}
Find other constructions of $E_8$. \added{(For instance via the extended Hamming code and
Construction~A, \S\ref{L04b:sec:codes}.)}
\end{exercise}

% ------------------------------------------------------------
\section{Laminated lattices}
\label{L04b:sec:laminated}

Before (Lecture~4) we had the lattices
\[
  A_1,\ A_2,\ A_3,\ D_4,\ D_5,\ E_6,\ E_7,\ E_8
\]
as the best lattice packings, and best known packings, in dimensions $1$ to $8$. Somehow these
are all related to $E_8$ (as we just saw, they all sit inside $E_8$).

\begin{addedblock}
Optimality among lattices is proven in dimensions $1$--$8$ (Lagrange $n=2$, Gauss $n=3$,
Korkine--Zolotarev $n=4,5$, Blichfeldt $n=6,7,8$); among \emph{all} packings it is proven for
$n=2$ (Thue, Fejes T\'oth), $n=3$ (Hales) and $n=8$ (Viazovska, posted March 2016, during this course).
Similarly the Leech lattice is optimal in dimension $24$ (Cohn--Kumar among lattices;
Cohn--Kumar--Miller--Radchenko--Viazovska 2016/17 among all packings).
\end{addedblock}

Dimension $9$ is, as far as I know, still unknown.
\added{(This is still the case, to the editor's knowledge.)}

The best constructions for lattices come from an ``inductive'' process. Start with the even
integers $\Lambda_1 = 2\Z\subset\R^1$.

\begin{definition}\label{L04b:def:laminated}
The \emph{$n$-dimensional laminated lattice} $\Lambda_n$ is a lattice of maximal density among
all lattices in $\R^n$ with shortest vector of length $\norm{x}=2$ that contain $\Lambda_{n-1}$
as a sublattice.
\end{definition}
So $\Lambda_n$ is obtained by gluing: $\Lambda_n=\langle \Lambda_{n-1}, g\rangle$\added{, i.e.\ by stacking
layers $\Lambda_{n-1}+kg$, $k\in\Z$, with each new layer placed over the deepest holes of the previous one}.
\begin{center}
\begin{tikzpicture}[scale=0.45]
  % A1
  \foreach \i in {0,...,4} \fill (\i,0) circle (2.5pt);
  \node at (2,-1.6) {$A_1$};
  \draw[-{Stealth}] (5,0)--(6.5,0);
  % A2
  \begin{scope}[xshift=7.5cm]
    \foreach \j in {-2,...,2} \foreach \i in {0,...,4} {
      \pgfmathsetmacro{\xx}{\i+0.5*mod(abs(\j),2)} \fill (\xx,{0.866*\j}) circle (2.5pt); }
    \node at (2.2,-3.0) {$A_2$};
  \end{scope}
  \draw[-{Stealth}] (13.3,0)--(14.8,0);
  % A3: hexagonal layer + next layer over holes
  \begin{scope}[xshift=16cm]
    \foreach \j in {-2,...,2} \foreach \i in {0,...,4} {
      \pgfmathsetmacro{\xx}{\i+0.5*mod(abs(\j),2)} \fill (\xx,{0.866*\j}) circle (2.5pt); }
    \foreach \j in {-2,...,1} \foreach \i in {0,...,3} {
      \pgfmathsetmacro{\xx}{\i+0.5+0.5*mod(abs(\j),2)} \draw (\xx,{0.866*\j+0.2887}) circle (3.5pt); }
    \node at (2.2,-3.0) {$A_3$};
  \end{scope}
  \draw[-{Stealth}] (21.8,0)--(23.3,0);
  \node at (24.6,0) {$\cdots\ E_8$};
\end{tikzpicture}
\end{center}
\orig{p.~40 sketch; circles = next layer}
The laminated lattices $\Lambda_1,\dots,\Lambda_8$ coincide with the root lattices above
\added{(scaled to minimal length $2$)}.

\begin{remark}
Dimension $9$ is unusual: there is a continuous family of packings built from layers of $\Lambda_8$,
\fixed{notes: ``cts family of lattices''}
all of the same density as $\Lambda_9$. \added{These are the ``fluid'' packings; in general they are
\emph{not} lattice packings, so ``family of lattices'' should be read as ``family of packings'' (see
Conway--Sloane, \emph{Sphere Packings, Lattices and Groups}, Ch.~1 and Ch.~6).}
\end{remark}

% ------------------------------------------------------------
\section{Dimension 10: codes and Construction A}
\label{L04b:sec:codes}

In dimension $10$, a \emph{non-lattice} packing beats the best known lattices.

\begin{definition}\label{L04b:def:code}
A \emph{binary code} of length $n$ is a subset $C\subseteq\{0,1\}^n$ of the binary words of length $n$.
The elements of $C$ are the \emph{code words}. \added{The \emph{Hamming distance} between two words is the
number of positions in which they differ; $d(C)$ denotes the minimal distance between distinct code words.}
\end{definition}

\begin{definition}[Construction A]\label{L04b:def:constrA}
Given a binary code $C$ of length $n$, there is a packing
\fixed{notes: $|C|=n$}
\[
  P(C) = \{x\in\Z^n : x \bmod 2 \in C\}.
\]
\end{definition}
\begin{addedblock}
Two distinct points of $P(C)$ either differ by an element of $2\Z^n$ (distance $\ge2$) or reduce to
different code words, hence differ by $\pm1$ in at least $d(C)$ coordinates (distance $\ge\sqrt{d(C)}$).
So $P(C)$ is a packing of balls of radius $\frac12\min(2,\sqrt{d(C)})$ with $|C|$ centres per cube
$[0,2)^n$, and its centre density is $|C|\,2^{-n}\big(\tfrac12\min(2,\sqrt{d(C)})\big)^n$.
\end{addedblock}

This gives some relation between packing problems in high dimensions and error detecting /
error correcting codes. Think of a code word sent through a continuous channel and received
with an error of size up to $\varepsilon$: the centres of the $\varepsilon$-spheres of a packing correspond
to code words that are at distance $\ge 2\varepsilon$ apart. Dense packings correspond to large collections of
code words that are well separated.
\orig{sketch: $\varepsilon$-ball around a code word}

\begin{example}[Block codes]
The $(6,3)$ repetition code repeats the pattern: $101\mapsto101101$. The $(3,2)$ parity check code
(mod $2$) appends a parity bit: $11\mapsto 110$. We can see the latter on the Hamming cube: its code
words $000,011,101,110$ all have Hamming distance $2$, i.e.\ are $2$ bits apart.
\begin{center}
\begin{tikzpicture}[scale=1.3]
  \coordinate (000) at (0,0); \coordinate (100) at (1,0); \coordinate (010) at (0.45,0.35);
  \coordinate (001) at (0,1); \coordinate (110) at (1.45,0.35); \coordinate (101) at (1,1);
  \coordinate (011) at (0.45,1.35); \coordinate (111) at (1.45,1.35);
  \draw (000)--(100)--(101)--(001)--cycle; \draw (001)--(011)--(111)--(101); \draw (100)--(110)--(111);
  \draw[dashed] (000)--(010)--(110); \draw[dashed] (010)--(011);
  \foreach \p in {000,011,101,110} {\fill[addcol] (\p) circle (2.2pt);}
  \node[below left,font=\scriptsize] at (000) {$000$}; \node[above,font=\scriptsize] at (011) {$011$};
  \node[right,font=\scriptsize] at (101) {$101$}; \node[right,font=\scriptsize] at (110) {$110$};
\end{tikzpicture}
\end{center}
In general, parity check codes (all even-weight words) correspond to $D_n$, the lattice with even
coordinate sums.
\end{example}

Construction A thus links binary codes and lattices: $P(C)$ is a lattice if and only if $C$ is a
\emph{linear} code, that is, sums of code words (mod $2$) are code words.
The Hamming codes give $E_7$ and $E_8$. \added{More precisely (after scaling by $1/\sqrt2$): the extended
Hamming code $[8,4,4]$ gives $E_8$ ($16+14\cdot16=240$ vectors of norm $4$ in $P(C)$), and the $[7,3,4]$
simplex code, the dual of the Hamming code $[7,4,3]$, gives $E_7$ ($14+7\cdot16=126$).}

\subsection*{The code \texorpdfstring{$C_{10}$}{C10}}
So in dimension $10$ there is a special binary code $C_{10}$ \added{(Best's code)}, which is unique
\added{(Litsyn--Vardy)} and quite dense. It is based on a Gray code, or reflected binary code.

\begin{definition}
A \emph{Gray code} is a map $(\Z_2)^n\to(\Z_2)^n$, \added{a re-ordering of all binary words of length $n$},
where the successor of an element differs by a single bit flip.
\end{definition}
\[
  \begin{array}{ccc}
    \text{binary} & & \text{Gray}\\ \hline
    00 & 0 & 00\\ 01 & 1 & 01\\ 10 & 2 & 11\\ 11 & 3 & 10
  \end{array}
\]
There is a replacement procedure to convert between binary and Gray codes.

\begin{exercise}
Find the conversion binary $\leftrightarrow$ Gray, and iterate it to get the Gray code of length $n$.
\added{(Hint: $g = b \oplus (b\gg1)$; conversely $b_k = g_1\oplus\dots\oplus g_k$ from the leading bit.)}
\end{exercise}

In dimension $10$ \added{we look for} binary codes with minimum distance at least $4$. We construct
$C_{10}$ as follows. Take $(a,b,c,d,e)\in(\Z_4)^5$ satisfying
\[
  b,c,d\in\{\pm1\},\qquad a = c-d,\qquad e = b+c,
\]
\unclear{$\{\pm1\}$: a crossed-out entry; $\{\pm1\}$ fits the count $2^3$}
together with all cyclic permutations thereof. Then apply the Gray map
\[
  0\mapsto 00,\quad 1\mapsto 01,\quad 2\mapsto 11,\quad 3\mapsto 10.
\]
This gives $8\cdot5=40$ code words ($2^3$ choices of $(b,c,d)$ times $5$ cyclic shifts) of length $10$.

\begin{claim}
The Hamming distance between distinct code words of $C_{10}$ is at least $4$.
\end{claim}
\begin{addedblock}
The Gray map is an isometry from $\Z_4$ with the Lee metric ($d(0,2)=d(1,3)=2$, other distinct pairs
distance $1$) to $\Z_2^2$ with the Hamming metric, so it suffices to check Lee distance $\ge4$ in $(\Z_4)^5$.
(Editor: confirmed by direct computer check of all $\binom{40}{2}$ pairs; the $40$ words are distinct.)
\end{addedblock}

Moreover $A(10,4)=40$, i.e.\ $40$ is the maximal size of a binary code of length $10$ and minimal
distance $4$, so $C_{10}$ is optimal for dimension $10$.
\unclear{notes: ``LP bound $\Rightarrow<40.\ldots$''; $A(10,4)=40$ is due to Best (c.~1980)}
\begin{addedblock}
Construction A applied to $C_{10}$ gives the non-lattice packing $P_{10c}$: radius $1$ and $40$ centres
per $2^{10}$ volume, so centre density $40/2^{10}=5/128\approx0.0391$, larger than that of the best known
lattice $\Lambda_{10}$ (centre density $1/(16\sqrt3)\approx0.0361$). See Conway--Sloane, Ch.~5.
\end{addedblock}

% ------------------------------------------------------------
\section{An aside: lattice vs.\ non-lattice packings of other objects}
\label{L04b:sec:aside}

For other objects than balls, the question of lattice versus non-lattice packings looks different:
\begin{itemize}
  \item ellipsoid packings;
  \item polygon packings;
  \item honeycomb (HCB) and bubble problems. \unclear{``HCB'', ``Bubble prob.''?}
\end{itemize}
\added{For example, in the plane a lattice packing is optimal for any centrally symmetric convex body
(Fejes T\'oth, Rogers), but not for some non-symmetric ones; and congruent ellipsoids in $\R^3$ can be packed
more densely by non-lattice arrangements than by any lattice (Bezdek--Kuperberg). Hales proved the honeycomb
conjecture in 1999.} Density issues are subtle here, e.g.\ in $2$d versus $1$d.

On the same page there is a computation: with $a^2b=1$,
\[
  \begin{bmatrix} 1&&\\ &a&\\ &&a^2\end{bmatrix}
  \begin{pmatrix}1\\1\\1\end{pmatrix} = \begin{pmatrix}1\\a\\a^2\end{pmatrix},\qquad
  \begin{bmatrix} 1&&\\ &a&\\ &&a^2\end{bmatrix}
  \begin{pmatrix}a\\a\\b\end{pmatrix} = \begin{pmatrix}a\\a^2\\1\end{pmatrix}.
\]
\unclear{purpose unclear; maybe an ellipsoid-packing example}

% ------------------------------------------------------------
\section*{Projects and quiz}
\orig{admin (date, room) omitted}
Ideas for projects:
\begin{itemize}
  \item error correcting codes and information theory;
  \item other sphere packing bounds;
  \item entropy of jammed systems;
  \item translation of an older work on packings. \unclear{last item barely legible}
\end{itemize}
Format of the quiz: definitions, short answers, sketches.


% Interlude (new material): numbered "I" so that lecture numbers stay aligned with chapter numbers
\begingroup
\renewcommand{\thechapter}{I}\renewcommand{\theHchapter}{interlude}
% MAT670 -- Interlude after Lecture 6: from Dynkin diagrams to Lie algebras and Lie groups.
% Entirely new material (not in the handwritten notes). Labels prefixed LIE:.

\chapter[Interlude: Lie algebras and Lie groups]{Interlude: From Dynkin diagrams to Lie algebras and Lie groups}
\label{LIE:ch}

\begin{addedblock}
This interlude is new: it is not in the handwritten notes. It picks up the remark at the start of
Lecture~4, that the classification of root systems is ``half of the classification of complex
semisimple Lie algebras'', and follows it. To keep it readable it is set in black; all of it should be
read as added material.
\end{addedblock}

In Lectures~4--6 we took the root system axioms as given, classified the possible Dynkin diagrams, and
then built each root system by hand from lattices. Two questions were left hanging.
\begin{itemize}
\item Where do the axioms come from? Why should anyone care about a finite set of vectors closed under
  reflections, with the integrality condition $n_{\beta\alpha}\in\Z$?
\item The same root system $A_{n-1}$ appears twice in these notes, once as the vectors
  $e_i-e_j$ in the hyperplane $\sum x_i=0$ (Lecture~6) and, later, once as the walls $x_i=x_j$ of
  $\Conf(N,\R)$ (Lecture~8). Is there a single object behind both?
\end{itemize}
The answer to both is \emph{matrices}. In this interlude we find root systems inside matrix
algebras, see that the crystallographic axioms are forced by a small piece of linear algebra, and
then use the exponential map to pass from the algebra to groups of matrices. At the end the root
lattice comes back, this time sitting inside a group.

Throughout, $M_n$ denotes the $n\times n$ complex matrices and $E_{ij}$ the matrix with a $1$ in
position $(i,j)$ and zeros elsewhere, so that
\[
  E_{ij}E_{kl}=\delta_{jk}E_{il}.
\]

%======================================================================
\section{The roots of \texorpdfstring{$A_{n-1}$}{A(n-1)} inside matrices}
\label{LIE:sec:sln}

Let $h=\operatorname{diag}(h_1,\dots,h_n)$ be a diagonal matrix. Multiplying a matrix unit by $h$ on
the left scales its row and on the right scales its column, so
\[
  hE_{ij}=h_iE_{ij},\qquad E_{ij}h=h_jE_{ij},
  \qquad\text{hence}\qquad
  hE_{ij}-E_{ij}h=(h_i-h_j)\,E_{ij}.
\]
In words: the linear map
\[
  \operatorname{ad}_h\colon M_n\to M_n,\qquad X\mapsto [h,X]:=hX-Xh
\]
is \emph{diagonal} in the basis $\{E_{ij}\}$, with eigenvalue $h_i-h_j$ on $E_{ij}$.

Now regard the eigenvalue as a function of $h$. On $E_{ij}$ it is the linear functional
$h\mapsto h_i-h_j$. If we identify functionals with vectors using the standard inner product, so that
$h\mapsto h_i$ corresponds to $e_i$, then the eigenvalue on $E_{ij}$ is the vector $e_i-e_j$. So:
\begin{center}
\emph{the nonzero eigenvalues of $\operatorname{ad}_h$, as $h$ runs over diagonal matrices, are exactly
the roots $e_i-e_j$ of $A_{n-1}$.}
\end{center}
The diagonal matrices themselves have eigenvalue $0$.

Where is the hyperplane $\sum x_i=0$? Scalar matrices $h=cI$ commute with everything, so they tell us
nothing, and it is natural to discard them. That leaves the \emph{traceless} matrices
\[
  \mathfrak{sl}_n=\{X\in M_n:\operatorname{tr}X=0\},
\]
whose diagonal part
\[
  \mathfrak h=\{\operatorname{diag}(h_1,\dots,h_n): h_1+\dots+h_n=0\}
\]
is exactly the hyperplane $\sum x_i=0$ of Lecture~6, of dimension $n-1$. We get a decomposition into
eigenspaces
\begin{equation}\label{LIE:eq:rootdecomp}
  \mathfrak{sl}_n=\mathfrak h\ \oplus\ \bigoplus_{\alpha\in\Delta}\mathfrak g_\alpha,
  \qquad
  \mathfrak g_{e_i-e_j}=\C E_{ij},
\end{equation}
where $\Delta=\{e_i-e_j: i\neq j\}$ is the root system $A_{n-1}$. As a check on dimensions:
$n^2-1=(n-1)+n(n-1)$, which is (rank) $+$ (number of roots), using the count
$\abs{\Delta}=n(n-1)$ from Lecture~6.

\begin{example}[$\mathfrak{sl}_3$]\label{LIE:ex:sl3}
Here $\mathfrak h$ is two-dimensional and the six off-diagonal matrix units sit at the six roots of
$A_2$. Placing each $E_{ij}$ at its root gives the familiar hexagon:
\begin{center}
\begin{tikzpicture}[scale=1.55]
  \foreach \a/\lab/\pos in {0/E_{12}/right, 60/E_{13}/above right, 120/E_{23}/above left,
                            180/E_{21}/left, 240/E_{31}/below left, 300/E_{32}/below right}
    {\draw[-{Stealth}] (0,0)--(\a:1); \fill (\a:1) circle (0.03);
     \node[\pos,font=\small] at (\a:1) {$\lab$};}
  \fill (0,0) circle (0.045); \fill (0.09,0) circle (0.045);
  \node[below,font=\small] at (0.05,-0.08) {$\mathfrak h$};
  \draw[dashed,-{Stealth},addcol,thick] (120:1.05) to[bend left=35] node[above,font=\scriptsize]{$\operatorname{ad}_{E_{12}}$} (60:1.05);
  \node[font=\scriptsize,gray] at (2.6,0.55) {$\alpha_1=e_1-e_2$};
  \node[font=\scriptsize,gray] at (2.6,0.25) {$\alpha_2=e_2-e_3$};
  \node[font=\scriptsize,gray] at (2.6,-0.05) {$\alpha_1+\alpha_2=e_1-e_3$};
\end{tikzpicture}
\end{center}
The two dots at the centre stand for the two-dimensional zero eigenspace $\mathfrak h$. The dashed
arrow anticipates Section~\ref{LIE:sec:bracket}: bracketing with $E_{12}$ moves $E_{23}$ to $E_{13}$,
that is, it \emph{adds the root} $\alpha_1$.
\end{example}

The choices made in Lecture~4 also have matrix meanings. Take the direction $d=(n,n-1,\dots,1)$, so
that $\ip{e_i-e_j}{d}=j-i$. Then:
\begin{itemize}
\item the positive roots $e_i-e_j$ with $i<j$ are the matrix units \emph{above} the diagonal, and
  the negative roots are those below it;
\item the simple roots $\alpha_i=e_i-e_{i+1}$ are the entries just above the diagonal,
  $E_{12},E_{23},\dots,E_{n-1,n}$;
\item every positive root is a sum of simple roots: $e_i-e_j=\alpha_i+\alpha_{i+1}+\dots+\alpha_{j-1}$.
  On the matrix side this is the statement that $E_{ij}$ is an iterated product (indeed an iterated
  commutator) of superdiagonal units.
\end{itemize}
The Weyl group, generated by the reflections $s_{e_i-e_j}$, acts on $\mathfrak h$ by swapping $h_i$ and
$h_j$. It is the symmetric group $\Sigma_n$ permuting the diagonal entries, which is conjugation by
permutation matrices. The Weyl chamber containing $d$ is $\{h_1>h_2>\dots>h_n\}$, one of the $n!$
chambers of $\Conf(n,\R)$ from Lecture~8.

%======================================================================
\section{The bracket and the definition of a Lie algebra}
\label{LIE:sec:bracket}

What made the computation work was the operation $[X,Y]=XY-YX$, the \emph{commutator}. It has three
properties that are easy to check:
\begin{enumerate}[label=(\roman*)]
\item it is bilinear;
\item it is antisymmetric, $[X,Y]=-[Y,X]$;
\item it satisfies the \emph{Jacobi identity}
  \[
    [X,[Y,Z]]=[[X,Y],Z]+[Y,[X,Z]],
  \]
  which says that $\operatorname{ad}_X=[X,\cdot\,]$ obeys the product rule (it is a derivation of the
  bracket). Expanding both sides into products of three matrices proves it.
\end{enumerate}

\begin{definition}\label{LIE:def:liealg}
A \emph{Lie algebra} over a field $k$ is a vector space $\mathfrak g$ with a bilinear map
$[\cdot,\cdot]\colon\mathfrak g\times\mathfrak g\to\mathfrak g$ that is antisymmetric and satisfies the
Jacobi identity. A \emph{subalgebra} is a subspace closed under the bracket; an \emph{ideal} is a
subspace $I$ with $[\mathfrak g,I]\subseteq I$. A Lie algebra is \emph{simple} if it is not abelian and
has no ideals other than $0$ and $\mathfrak g$. A \emph{representation} of $\mathfrak g$ on a vector space
$V$ is a linear map $\rho\colon\mathfrak g\to\operatorname{End}(V)$ with
$\rho([X,Y])=\rho(X)\rho(Y)-\rho(Y)\rho(X)$.
\end{definition}

\begin{example}\leavevmode
\begin{enumerate}[nosep]
\item $\mathfrak{gl}_n=M_n$ with the commutator.
\item $\mathfrak{sl}_n$ is a subalgebra, since $\operatorname{tr}(XY-YX)=0$ for all $X,Y$. It is simple for
  $n\ge2$.
\item $\mathfrak{so}_n=\{X\in M_n: X^T=-X\}$ is a subalgebra: if $X,Y$ are antisymmetric, then
  $[X,Y]^T=Y^TX^T-X^TY^T=YX-XY=-[X,Y]$.
\item $\R^3$ with the cross product is a Lie algebra (Jacobi is the identity
  $a\times(b\times c)=(a\times b)\times c+b\times(a\times c)$). Sending $\omega\in\R^3$ to the matrix
  $X_\omega$ with $X_\omega v=\omega\times v$ gives an isomorphism $(\R^3,\times)\cong\mathfrak{so}_3(\R)$.
\item By the Jacobi identity, $X\mapsto\operatorname{ad}_X$ is a representation of $\mathfrak g$ on itself,
  the \emph{adjoint representation}. Section~\ref{LIE:sec:sln} was a study of the adjoint
  representation of $\mathfrak{sl}_n$, restricted to $\mathfrak h$.
\end{enumerate}
\end{example}

The Jacobi identity also explains the dashed arrow in Example~\ref{LIE:ex:sl3}. Suppose
$[h,x]=\alpha(h)x$ and $[h,y]=\beta(h)y$ for all $h\in\mathfrak h$. Then
\[
  [h,[x,y]]=[[h,x],y]+[x,[h,y]]=\big(\alpha(h)+\beta(h)\big)[x,y].
\]

\begin{lemma}\label{LIE:lem:add}
$[\mathfrak g_\alpha,\mathfrak g_\beta]\subseteq\mathfrak g_{\alpha+\beta}$, where
$\mathfrak g_{\alpha+\beta}=0$ if $\alpha+\beta$ is neither a root nor $0$, and $\mathfrak g_0=\mathfrak h$.
\end{lemma}

So the bracket moves around the root diagram by \emph{adding roots}. In $\mathfrak{sl}_3$,
$[E_{12},E_{23}]=E_{13}$ is the statement $\alpha_1+\alpha_2=e_1-e_3$, while $[E_{12},E_{13}]=0$
because $2e_1-e_2-e_3$ is not a root.

%======================================================================
\section{The atom: \texorpdfstring{$\mathfrak{sl}_2$}{sl(2)}}
\label{LIE:sec:sl2}

Fix a root $\alpha=e_i-e_j$ of $A_{n-1}$ and put
\[
  e_\alpha=E_{ij},\qquad f_\alpha=E_{ji},\qquad h_\alpha=[e_\alpha,f_\alpha]=E_{ii}-E_{jj}\in\mathfrak h.
\]
A short computation gives
\begin{equation}\label{LIE:eq:sl2rel}
  [h_\alpha,e_\alpha]=2e_\alpha,\qquad [h_\alpha,f_\alpha]=-2f_\alpha,\qquad [e_\alpha,f_\alpha]=h_\alpha .
\end{equation}
For $n=2$ these three matrices
\[
  e=\begin{pmatrix}0&1\\0&0\end{pmatrix},\quad
  f=\begin{pmatrix}0&0\\1&0\end{pmatrix},\quad
  h=\begin{pmatrix}1&0\\0&-1\end{pmatrix}
\]
are a basis of $\mathfrak{sl}_2$. So every root $\alpha$ comes with a copy
$\mathfrak s_\alpha=\operatorname{span}(e_\alpha,f_\alpha,h_\alpha)\cong\mathfrak{sl}_2$ sitting inside $\mathfrak{sl}_n$.

Now look at how $h_\alpha$ acts on the other root spaces. For a root $\beta=e_k-e_l$,
\[
  \beta(h_\alpha)=(\delta_{ki}-\delta_{kj})-(\delta_{li}-\delta_{lj})=\ip{\beta}{\alpha}
  =\frac{2\ip{\beta}{\alpha}}{\ip{\alpha}{\alpha}}=n_{\beta\alpha},
\]
since $\ip{\alpha}{\alpha}=2$. So \emph{the Cartan integer $n_{\beta\alpha}$ is the eigenvalue of
$h_\alpha$ on $\mathfrak g_\beta$}. ($h_\alpha$ is called the \emph{coroot} of $\alpha$.) To see why these
eigenvalues have to be integers, we need one theorem about $\mathfrak{sl}_2$.

\begin{theorem}[Representations of $\mathfrak{sl}_2$]\label{LIE:thm:sl2}
Let $V$ be a finite-dimensional complex representation of $\mathfrak{sl}_2$. Then $h$ acts
diagonalizably on $V$, its eigenvalues are integers, and for every $k\in\Z$ the eigenspaces satisfy
$\dim V_k=\dim V_{-k}$. If $V$ is irreducible of dimension $m+1$, then the eigenvalues of $h$ are
$m,m-2,\dots,-m$, each with multiplicity one.
\end{theorem}

\begin{proof}[Proof for irreducible $V$]
The key observation is that $e$ raises $h$-eigenvalues by $2$ and $f$ lowers them by $2$. If
$hv=\lambda v$, then by \eqref{LIE:eq:sl2rel}
\[
  h(ev)=e(hv)+[h,e]v=(\lambda+2)\,ev,\qquad h(fv)=(\lambda-2)\,fv .
\]
Take any eigenvector of $h$ and apply $e$ repeatedly. The vectors obtained are eigenvectors for
distinct eigenvalues while they are nonzero, and $V$ is finite-dimensional, so eventually we reach a
\emph{highest weight vector}: $v\neq0$ with $hv=\lambda v$ and $ev=0$. Put $v_k=f^kv$. Then
$hv_k=(\lambda-2k)v_k$, and induction on $k$ using $ef=fe+h$ gives
\[
  e\,v_k=k(\lambda-k+1)\,v_{k-1}.
\]
Let $m$ be the largest index with $v_m\neq0$. Then $0=e\,v_{m+1}=(m+1)(\lambda-m)\,v_m$, so
$\lambda=m$. In particular $\lambda$ is a non-negative integer. The span of $v_0,\dots,v_m$ is closed
under $e$, $f$ and $h$, so by irreducibility it is all of $V$. The eigenvalues are therefore
$m,m-2,\dots,-m$, as claimed.
\end{proof}

\begin{figure}[ht]
\centering
\begin{tikzpicture}[xscale=1.7]
  \foreach \x/\l in {0/{-m},1/{-m+2},2/{\cdots},3/{m-2},4/{m}} {\fill (\x,0) circle (0.05); \node[font=\small] at (\x,-0.62) {$\l$};}
  \foreach \x in {0,1,2,3} {\draw[-{Stealth},addcol] (\x+0.1,0.1) to[bend left=40] (\x+0.9,0.1);
                            \draw[-{Stealth},flagcol] (\x+0.9,-0.1) to[bend left=40] (\x+0.1,-0.1);}
  \node[addcol,font=\small] at (2,0.7) {$e$ raises by $2$};
  \node[flagcol,font=\small] at (2,-1.05) {$f$ lowers by $2$};
\end{tikzpicture}
\caption{The $(m+1)$-dimensional irreducible representation of $\mathfrak{sl}_2$: a string of
$h$-eigenvalues, symmetric about $0$.}
\end{figure}

For a general finite-dimensional $V$ one uses \emph{Weyl's complete reducibility theorem}: $V$ is a
direct sum of irreducible representations (see Humphreys, \emph{Introduction to Lie Algebras and
Representation Theory}, \S6.3, or Hall, \emph{Lie Groups, Lie Algebras, and Representations}, Chs.~4 and~10).
The statements about eigenvalues then follow from the irreducible case.

%======================================================================
\section{Why the root system axioms hold}
\label{LIE:sec:axioms}

Now apply Theorem~\ref{LIE:thm:sl2} to the copy $\mathfrak s_\alpha\cong\mathfrak{sl}_2$, acting by
brackets. Fix roots $\beta\neq\pm\alpha$ and form the \emph{$\alpha$-string through $\beta$},
\[
  V=\bigoplus_{k\in\Z}\mathfrak g_{\beta+k\alpha}.
\]
By Lemma~\ref{LIE:lem:add}, bracketing with $e_\alpha$ or $f_\alpha$ shifts $k$ by $\pm1$, so $V$ is a
representation of $\mathfrak s_\alpha$. On the summand $\mathfrak g_{\beta+k\alpha}$, $h_\alpha$ acts by the
scalar
\[
  (\beta+k\alpha)(h_\alpha)=n_{\beta\alpha}+2k .
\]

\begin{itemize}
\item \textbf{Axiom 4 (crystallographic).} $n_{\beta\alpha}$ is an eigenvalue of $h_\alpha$ on $V$
  (take $k=0$), so by Theorem~\ref{LIE:thm:sl2} it is an integer.
\item \textbf{Axiom 3 (reflections).} The eigenvalues are symmetric about $0$, so $-n_{\beta\alpha}$ is
  also an eigenvalue: $n_{\beta\alpha}+2k=-n_{\beta\alpha}$ for some $k$ with
  $\mathfrak g_{\beta+k\alpha}\neq0$. That is, $k=-n_{\beta\alpha}$, and
  \[
    \beta-n_{\beta\alpha}\,\alpha=s_\alpha(\beta)
  \]
  is a root. The reflection $s_\alpha$ maps $\Delta$ to itself.
\item \textbf{Strings have no gaps.} Each irreducible piece of $V$ has eigenvalues forming an unbroken
  string with step $2$. In fact $V$ is irreducible here, since the root spaces turn out to be
  one-dimensional. So the roots $\beta+k\alpha$ form an unbroken string $\beta-p\alpha,\dots,\beta+q\alpha$
  with $p-q=n_{\beta\alpha}$. This is the ``$\alpha-\beta$'' phenomenon behind Lemma~\ref{L04a:lem:diff} of Lecture~4.
\end{itemize}

For $\mathfrak{sl}_n$ all of this can of course be checked directly. The point is that the argument used
only three ingredients: a commutative subalgebra $\mathfrak h$ acting diagonalizably, the copies of
$\mathfrak{sl}_2$ attached to the roots, and finite-dimensionality. These are exactly what a
\emph{complex semisimple Lie algebra} provides. In such an algebra a maximal commutative subalgebra of
diagonalizable elements (a \emph{Cartan subalgebra}) always exists, and its nonzero eigenvalues form a
reduced, crystallographic root system. (Reducedness, axiom~2, takes one more argument; see Humphreys
\S8.4.) \emph{The root system axioms are the shadow of $\mathfrak{sl}_2$ representation theory.}

%======================================================================
\section{The other classical families}
\label{LIE:sec:classical}

The families $B_n$, $C_n$, $D_n$ from Lecture~6 come from the Lie algebras that preserve a bilinear
form. Let $J$ be an invertible $N\times N$ matrix and put
\[
  \mathfrak g_J=\{X\in M_N: X^TJ+JX=0\}.
\]
This is a subalgebra (the same computation as for $\mathfrak{so}_n$), and in Section~\ref{LIE:sec:exp}
we will see that it consists of the \emph{infinitesimal symmetries} of the form $\ip{v}{w}_J=v^TJw$.
For $J=I$ this is $\mathfrak{so}_N$.

The trick is to choose $J$ so that a large family of \emph{diagonal} matrices lies in $\mathfrak g_J$.
Write $i'=N+1-i$ and let $J$ be the antidiagonal matrix with $J_{i,i'}=1$. For $X\in M_N$ one computes
$(JX)_{ij}=X_{i'j}$ and $(X^TJ)_{ij}=X_{j'i}$, so
\[
  X\in\mathfrak g_J\iff X_{ij}=-X_{j'i'}\ \text{for all } i,j .
\]
(Over $\C$ any two nondegenerate symmetric forms are equivalent, so $\mathfrak g_J\cong\mathfrak{so}_N(\C)$.)
\begin{itemize}
\item The diagonal matrices in $\mathfrak g_J$ are
  $h=\operatorname{diag}(h_1,\dots,h_n,[0],-h_n,\dots,-h_1)$, where the middle $0$ is present only when
  $N=2n+1$ is odd. These form a commutative subalgebra $\mathfrak h$ of dimension $n$.
\item The matrices $E_{ij}-E_{j'i'}$, for pairs with $i\neq j$ and $j\neq i'$, span the rest of
  $\mathfrak g_J$ (for $j=i'$ the expression is $0$: the antidiagonal of $X$ vanishes). Each is
  an eigenvector of $\operatorname{ad}_h$ with eigenvalue $\tilde h_i-\tilde h_j$, where $\tilde h$ is the
  diagonal of $h$. (The two matrix units give the same eigenvalue, because $\tilde h_{k'}=-\tilde h_k$.)
\item Reading off the eigenvalues: $i,j\le n$ gives $\pm(e_i-e_j)$; $i\le n<j$ gives
  $\pm(e_i+e_k)$ with $j=k'$ and $k\neq i$ (the case $k=i$ is excluded since then $j=i'$); and when $N$
  is odd the middle index gives $\pm e_i$.
\end{itemize}
So $\mathfrak{so}_{2n}$ has root system $\{\pm e_i\pm e_j\}=D_n$ and $\mathfrak{so}_{2n+1}$ has
$\{\pm e_i\pm e_j\}\cup\{\pm e_i\}=B_n$. If instead
$J=\left(\begin{smallmatrix}0&K\\-K&0\end{smallmatrix}\right)$, with $K$ the $n\times n$ antidiagonal
matrix of ones, the form is \emph{symplectic}. Now signs enter: with $\varepsilon_i=1$ for $i\le n$
and $\varepsilon_i=-1$ for $i>n$, the condition becomes $X_{ij}=-\varepsilon_i\varepsilon_jX_{j'i'}$, the
diagonal part is the same $\mathfrak h$ as before, and the off-diagonal basis is
$E_{ij}-\varepsilon_i\varepsilon_jE_{j'i'}$. The same bookkeeping gives the symplectic algebra
$\mathfrak{sp}_{2n}$, whose root system is $\{\pm e_i\pm e_j\}\cup\{\pm 2e_i\}=C_n$: now the case $k=i$ is
allowed, since for $j=i'$ the basis element is $2E_{ii'}\neq0$, and it produces the long roots $\pm2e_i$.

\begin{center}
\renewcommand{\arraystretch}{1.25}
\begin{tabular}{llllll}
\toprule
type & Lie algebra & condition & roots & rank $+\abs{\Delta}$ & $\dim\mathfrak g$\\
\midrule
$A_n$ & $\mathfrak{sl}_{n+1}$ & $\operatorname{tr}X=0$ & $e_i-e_j$ & $n+n(n+1)$ & $(n+1)^2-1$\\
$B_n$ & $\mathfrak{so}_{2n+1}$ & $X^TJ+JX=0$ & $\pm e_i\pm e_j,\ \pm e_i$ & $n+2n^2$ & $n(2n+1)$\\
$C_n$ & $\mathfrak{sp}_{2n}$ & $X^TJ+JX=0$ ($J$ skew) & $\pm e_i\pm e_j,\ \pm2e_i$ & $n+2n^2$ & $n(2n+1)$\\
$D_n$ & $\mathfrak{so}_{2n}$ & $X^TJ+JX=0$ & $\pm e_i\pm e_j$ & $n+2n(n-1)$ & $n(2n-1)$\\
\bottomrule
\end{tabular}
\end{center}
In every row, $\dim\mathfrak g=\operatorname{rank}+\abs{\Delta}$, as it must be by the root decomposition
\eqref{LIE:eq:rootdecomp}. The root counts are the ones computed in Lecture~6.

\begin{remark}[The ``repetitions'' in the classification]
Theorem~\ref{L04a:thm:classification} restricted the ranges of $n$ ($B_n$ for $n\ge2$, $C_n$ for $n\ge3$, $D_n$ for $n\ge4$) to
avoid repeating diagrams. Each coincidence of small diagrams is an isomorphism of Lie algebras:
$B_1=C_1=A_1$ gives $\mathfrak{so}_3\cong\mathfrak{sp}_2=\mathfrak{sl}_2$; $B_2=C_2$ gives
$\mathfrak{so}_5\cong\mathfrak{sp}_4$; $D_2=A_1\times A_1$ gives $\mathfrak{so}_4\cong\mathfrak{sl}_2\times\mathfrak{sl}_2$; and
$D_3=A_3$ gives $\mathfrak{so}_6\cong\mathfrak{sl}_4$.
\end{remark}

%======================================================================
\section{The classification of simple Lie algebras}
\label{LIE:sec:classification}

Everything so far points to a dictionary between Dynkin diagrams and Lie algebras. The result is
due to Killing (1888--90) and \'E.~Cartan (thesis, 1894); the diagrams themselves are due to
Dynkin (1946).

\begin{theorem}[Killing--Cartan classification]\label{LIE:thm:KC}
Let $\mathfrak g$ be a complex semisimple Lie algebra.
\begin{enumerate}[nosep]
\item $\mathfrak g$ has a Cartan subalgebra $\mathfrak h$, unique up to automorphisms of $\mathfrak g$, and the
  nonzero eigenvalues of $\operatorname{ad}\mathfrak h$ form a root system $\Delta(\mathfrak g)$ in (the real
  span of) $\mathfrak h^*$, with one-dimensional root spaces.
\item $\mathfrak g$ is simple if and only if $\Delta(\mathfrak g)$ is irreducible.
\item $\mathfrak g\cong\mathfrak g'$ if and only if $\Delta(\mathfrak g)\cong\Delta(\mathfrak g')$.
\item (Existence; Serre's theorem) Every root system arises this way.
\end{enumerate}
Consequently the complex simple Lie algebras correspond exactly to the connected Dynkin diagrams
$A_n$ ($n\ge1$), $B_n$ ($n\ge2$), $C_n$ ($n\ge3$), $D_n$ ($n\ge4$), $E_6$, $E_7$, $E_8$, $F_4$, $G_2$.
\end{theorem}

Serre's theorem even writes the algebra down from the diagram. Let $\alpha_1,\dots,\alpha_r$ be the
simple roots and $a_{ij}=n_{\alpha_j\alpha_i}$ the \emph{Cartan matrix}. Then $\mathfrak g$ is generated by
$e_i,f_i,h_i$ ($1\le i\le r$), which are exactly the $\mathfrak{sl}_2$-triples attached to the simple
roots, subject to the relations
\[
  [h_i,h_j]=0,\quad [h_i,e_j]=a_{ij}e_j,\quad [h_i,f_j]=-a_{ij}f_j,\quad [e_i,f_j]=\delta_{ij}h_i,
\]
\[
  (\operatorname{ad}e_i)^{1-a_{ij}}(e_j)=0,\qquad (\operatorname{ad}f_i)^{1-a_{ij}}(f_j)=0\qquad(i\neq j).
\]
The last two relations say that an $\alpha_i$-string through $\alpha_j$ has exactly the length the
Cartan integer predicts. For $\mathfrak{sl}_n$ this recovers $e_i=E_{i,i+1}$, $f_i=E_{i+1,i}$ and
$h_i=E_{ii}-E_{i+1,i+1}$.

The dimension formula $\dim\mathfrak g=\operatorname{rank}+\abs{\Delta}$ and the root counts ($12$ for
$G_2$ from Lecture~4; $48$, $126$, $240$ from Lecture~6; and $72$ for $E_6$, which one checks by counting
the roots of $E_8$ orthogonal to an $A_2$) give the dimensions of the exceptional algebras:
\begin{center}
\renewcommand{\arraystretch}{1.2}
\begin{tabular}{lccccc}
\toprule
 & $G_2$ & $F_4$ & $E_6$ & $E_7$ & $E_8$\\
\midrule
rank & 2 & 4 & 6 & 7 & 8\\
$\abs{\Delta}$ & 12 & 48 & 72 & 126 & 240\\
$\dim\mathfrak g$ & 14 & 52 & 78 & 133 & 248\\
\bottomrule
\end{tabular}
\end{center}
In particular, the $240$ minimal vectors of the $E_8$ lattice, the kissing configuration in
dimension $8$, \emph{literally} label the root spaces of the largest exceptional Lie algebra,
$\mathfrak e_8=\mathfrak h\oplus\bigoplus_{\alpha\in E_8}\mathfrak g_\alpha$, $248=8+240$. The smallest one,
$\mathfrak g_2$, is the algebra of derivations of the octonions (\'E.~Cartan).

%======================================================================
\section{From the algebra to the group: the exponential map}
\label{LIE:sec:exp}

The name ``Lie algebra'' comes from Sophus Lie's idea that a group of transformations is determined,
near the identity, by its \emph{infinitesimal} transformations. For matrices this can be made
completely explicit.

\begin{definition}
For $X\in M_n$ (real or complex) the \emph{matrix exponential} is
\[
  \exp X=e^X=\sum_{k=0}^\infty\frac{X^k}{k!}.
\]
The series converges absolutely, since $\norm{X^k}\le\norm{X}^k$ in the operator norm.
\end{definition}

\begin{proposition}\label{LIE:prop:exp}
For $X,Y\in M_n$ and $g\in\GL_n$:
\begin{enumerate}[nosep]
\item if $XY=YX$, then $e^{X+Y}=e^Xe^Y$; in particular $e^{(s+t)X}=e^{sX}e^{tX}$ and $(e^X)^{-1}=e^{-X}$;
\item $\frac{d}{dt}e^{tX}=Xe^{tX}=e^{tX}X$;
\item $e^{gXg^{-1}}=g\,e^Xg^{-1}$ and $(e^X)^T=e^{X^T}$;
\item $\det e^X=e^{\operatorname{tr}X}$;
\item $\exp$ maps a neighbourhood of $0$ diffeomorphically onto a neighbourhood of $I$, with inverse
  $\log(I+A)=\sum_{k\ge1}(-1)^{k+1}A^k/k$ for $\norm{A}<1$.
\end{enumerate}
\end{proposition}

\begin{proof}[Sketch]
(1) For commuting $X$ and $Y$ the binomial theorem applies to $(X+Y)^k$, and the Cauchy product of the
two series gives the result. (2) Differentiate the series term by term. (3) Apply the identities
$(gXg^{-1})^k=gX^kg^{-1}$ and $(X^k)^T=(X^T)^k$ to each term. (4) Over $\C$, write $X=gTg^{-1}$ with
$T$ upper triangular (Schur). Then $e^T$ is upper triangular with diagonal entries $e^{t_{ii}}$, so
$\det e^X=\det e^T=e^{\sum t_{ii}}=e^{\operatorname{tr}X}$. (5) The derivative of $\exp$ at $0$ is the
identity, so the inverse function theorem applies; the series for $\log$ inverts $\exp$ near $I$ as a
formal identity of power series.
\end{proof}

By (1), $t\mapsto e^{tX}$ is a homomorphism $(\R,+)\to\GL_n$: a \emph{one-parameter subgroup}. This is
the matrix version of the one-parameter groups of diffeomorphisms of Lecture~9. The linear vector field $v\mapsto Xv$ on $\R^n$ generates the
one-parameter group of diffeomorphisms $\varphi_t(v)=e^{tX}v$, because $\frac{d}{dt}e^{tX}v=X(e^{tX}v)$
by (2). So $X$ is the ``velocity'' of the subgroup at the identity.

\begin{figure}[ht]
\centering
\begin{tikzpicture}[scale=1.4]
  \draw[thick] (0,0) circle (1);
  \draw[addcol,thick] (1,-1.35)--(1,1.35) node[above,font=\small]{$T_I\SO(2)=\mathfrak{so}_2\cong\R$};
  \fill (1,0) circle (0.03) node[right,font=\small]{$I$};
  \foreach \t in {0.6,1.2} {\fill[addcol] (1,\t) circle (0.03); \fill ({cos(\t r)},{sin(\t r)}) circle (0.03);
     \draw[-{Stealth},gray,dashed] (0.97,\t) to[bend right=20] ({0.97*cos(\t r)},{0.97*sin(\t r)});}
  \node[font=\small] at (0,0) {$\SO(2)\cong\Sph^1$};
  \node[font=\scriptsize,gray] at (2.45,0.45) {$\theta\mapsto\begin{pmatrix}\cos\theta&-\sin\theta\\ \sin\theta&\cos\theta\end{pmatrix}$};
\end{tikzpicture}
\caption{The exponential map wraps the tangent line $\mathfrak{so}_2$ around the circle $\SO(2)$.
Its kernel is $2\pi\Z$, a lattice.}
\label{LIE:fig:so2}
\end{figure}

\subsection*{Matrix groups and their Lie algebras}

\begin{definition}
A \emph{matrix group} is a closed subgroup $G\subseteq\GL_n(\R)$ or $\GL_n(\C)$. Its \emph{Lie algebra} is
\[
  \operatorname{Lie}(G)=\{X\in M_n: e^{tX}\in G\text{ for all }t\in\R\}.
\]
\end{definition}

\begin{theorem}[von Neumann 1929; Cartan 1930]\label{LIE:thm:closed}
Let $G$ be a matrix group. Then $\operatorname{Lie}(G)$ is a real vector space closed under the commutator,
so it is a Lie algebra. Moreover $G$ is a smooth submanifold of $M_n$ (a \emph{Lie group}) with tangent
space $T_IG=\operatorname{Lie}(G)$, and $\exp$ maps a neighbourhood of $0$ in $\operatorname{Lie}(G)$
diffeomorphically onto a neighbourhood of $I$ in $G$.
\end{theorem}

We will not prove that $G$ is a manifold (see Hall, Ch.~3, or Stillwell, \emph{Naive Lie Theory}, Ch.~7), but it is worth seeing where the bracket comes from. Closure under addition follows from the
\emph{Lie product formula} $e^{X+Y}=\lim_{k\to\infty}\big(e^{X/k}e^{Y/k}\big)^k$. For the bracket, take
$X,Y\in\operatorname{Lie}(G)$. For each $t$ the matrix $e^{tX}Ye^{-tX}$ lies in $\operatorname{Lie}(G)$, because
$\exp\!\big(s\,e^{tX}Ye^{-tX}\big)=e^{tX}e^{sY}e^{-tX}\in G$. Since $\operatorname{Lie}(G)$ is a closed
subspace, it also contains the derivative
\[
  \frac{d}{dt}\Big|_{t=0}e^{tX}Ye^{-tX}=XY-YX=[X,Y].
\]
Equivalently, expanding the group commutator gives
\[
  e^{sX}e^{tY}e^{-sX}e^{-tY}=I+st\,[X,Y]+(\text{terms of total degree}\ge3\text{ in }s,t).
\]
\emph{The bracket is the infinitesimal shadow of the failure of the group to commute.}

\begin{example}[The classical groups recover the classical algebras]\leavevmode
\begin{itemize}[nosep]
\item $\operatorname{Lie}(\GL_n)=M_n=\mathfrak{gl}_n$.
\item $\operatorname{Lie}(\SL_n)=\mathfrak{sl}_n$. If $\operatorname{tr}X=0$ then $\det e^{tX}=1$ by
  Proposition~\ref{LIE:prop:exp}(4); conversely, $\det e^{tX}=e^{t\operatorname{tr}X}=1$ for all $t$ forces
  $\operatorname{tr}X=0$.
\item $\operatorname{Lie}(G_J)=\mathfrak g_J$ for the isometry group $G_J=\{g: g^TJg=J\}$ of the form $J$. If
  $X^TJ+JX=0$, then $X^T=-JXJ^{-1}$, so $(e^{X})^TJe^{X}=Je^{-X}J^{-1}Je^{X}=J$; applying this to $tX$ shows $\mathfrak g_J\subseteq
  \operatorname{Lie}(G_J)$. Conversely, differentiating $(e^{tX})^TJe^{tX}=J$ at $t=0$ gives $X^TJ+JX=0$. Thus $\mathfrak{so}_N$
  belongs to $\Orth(N)$ and $\SO(N)$, and $\mathfrak{sp}_{2n}$ to $\mathrm{Sp}_{2n}$. This is why $\mathfrak g_J$
  deserves the name ``infinitesimal isometries''.
\item The unitary group $\mathrm U(n)=\{g\in\GL_n(\C): g^*g=I\}$ has Lie algebra
  $\mathfrak u(n)=\{X: X^*=-X\}$ (skew-Hermitian matrices), and $\mathrm{SU}(n)$ has
  $\mathfrak{su}(n)=\mathfrak u(n)\cap\mathfrak{sl}_n$. Note that $\mathfrak{su}(n)\otimes\C=\mathfrak{sl}_n(\C)$: the
  compact group $\mathrm{SU}(n)$ is a \emph{real form} of $\mathfrak{sl}_n(\C)$.
\end{itemize}
\end{example}

\begin{example}[$\SO(2)$ and $\SO(3)$]\label{LIE:ex:so3}
For $X=\theta\left(\begin{smallmatrix}0&-1\\1&0\end{smallmatrix}\right)$ we have $X^2=-\theta^2I$, and
splitting the series into even and odd terms gives
\[
  e^X=\cos\theta\,I+\sin\theta\begin{pmatrix}0&-1\\1&0\end{pmatrix},
\]
which is rotation by $\theta$ (Figure~\ref{LIE:fig:so2}). In $\SO(3)$, write $X_\omega v=\omega\times v$
with $\theta=\abs{\omega}$. Using $X_\omega^3=-\theta^2X_\omega$, the same splitting gives
\emph{Rodrigues' formula}
\[
  e^{X_\omega}=I+\frac{\sin\theta}{\theta}X_\omega+\frac{1-\cos\theta}{\theta^2}X_\omega^2,
\]
which is rotation by angle $\theta$ about the axis $\omega$. Every rotation has an axis, so $\exp$ is
onto $\SO(3)$. It maps the closed ball $\abs{\omega}\le\pi$ onto $\SO(3)$, identifying only antipodal
points of the boundary sphere, so $\SO(3)\cong\R P^3$.
\end{example}

\begin{remark}[Same algebra, different groups]
The group $\mathrm{SU}(2)$ of unit quaternions is a $3$-sphere. It acts on the purely imaginary
quaternions $\cong\R^3$ by conjugation, and this gives a two-to-one homomorphism
$\mathrm{SU}(2)\to\SO(3)$ with kernel $\{\pm I\}$. The two groups have isomorphic Lie algebras,
$\mathfrak{su}(2)\cong\mathfrak{so}(3)\cong(\R^3,\times)$, but they are not isomorphic: one is simply connected,
the other has $\pi_1=\Z/2$. The Lie algebra determines a connected group only up to coverings. The
missing global information turns out to be a \emph{lattice}, as the next section shows.
\end{remark}

%======================================================================
\section{Tori, Weyl groups, and the lattice inside the group}
\label{LIE:sec:torus}

Return to $A_{n-1}$, now in the compact group $\mathrm{SU}(n)$. The diagonal matrices in $\mathrm{SU}(n)$ form
the \emph{maximal torus}
\[
  T=\big\{\operatorname{diag}(e^{i\theta_1},\dots,e^{i\theta_n}):\ \theta_1+\dots+\theta_n\equiv0\pmod{2\pi}\big\},
\]
with Lie algebra
\[
  \mathfrak t=\{\operatorname{diag}(i\theta_1,\dots,i\theta_n):\ \theta_1+\dots+\theta_n=0\}
  \cong\{\theta\in\R^n:\textstyle\sum\theta_j=0\},
\]
which is again the hyperplane of Lecture~6. Since $\mathfrak t$ is commutative, $\exp\colon\mathfrak t\to T$ is a
surjective homomorphism, and its kernel consists of the $\theta$ with every $\theta_j\in2\pi\Z$ and
$\sum\theta_j=0$:
\[
  \ker\big(\exp|_{\mathfrak t}\big)=2\pi\cdot\{x\in\Z^n:\textstyle\sum x_j=0\}=2\pi\cdot A_{n-1},
\]
the $A_{n-1}$ root lattice from Lecture~6. So
\[
  T\;\cong\;\mathfrak t\,/\,2\pi A_{n-1}:
\]
\emph{the maximal torus of $\mathrm{SU}(n)$ is Euclidean space modulo the $A_{n-1}$ root lattice.} For
$\mathrm{SU}(3)$ it is the plane modulo the hexagonal lattice, the lattice of the densest disk packing
from Lecture~1.

\paragraph{The Weyl group.} Permutation matrices (with one sign changed if necessary, to make the
determinant $1$) normalize $T$ and permute the $\theta_j$. The resulting group
\[
  W=N(T)/T\cong\Sigma_n
\]
is the Weyl group of $A_{n-1}$, generated by the reflections $s_{e_i-e_j}$, which swap $\theta_i$ and
$\theta_j$. By the spectral theorem every element of $\mathrm{SU}(n)$ is conjugate to an element of $T$,
uniquely up to the action of $W$, so
\[
  \{\text{conjugacy classes of }\mathrm{SU}(n)\}\;\longleftrightarrow\;T/W
  \;\longleftrightarrow\;\{\text{unordered $n$-tuples of eigenvalues on }\Sph^1\text{ with product }1\}.
\]

\begin{remark}[Back to configuration spaces]
An element of $\mathrm{SU}(n)$ is \emph{regular} when its eigenvalues are distinct, that is, when its
$n$ eigenvalues form a configuration of distinct points on the circle. Choose real eigenvalue angles with
$\sum\theta_j=0$, ordered so that $\theta_1\ge\theta_2\ge\dots\ge\theta_n\ge\theta_1-2\pi$ (this is
always possible, in exactly one way), and record the gaps
\[
  g_k=\theta_k-\theta_{k+1}\ (1\le k<n),\qquad g_n=2\pi-(\theta_1-\theta_n).
\]
The gaps are non-negative and sum to $2\pi$, and together with $\sum\theta_j=0$ they determine the
$\theta_j$. So $T/W$ is a closed simplex of dimension $n-1$, the \emph{Weyl alcove}. Its interior, the
regular classes, is identified through the gaps with the reduced configuration space of $n$ points on a
circle in a fixed cyclic order (one of the open simplices of gaps in Lecture~7). For $n=3$ it is the same open triangle of gap coordinates $(x,y,1-x-y)$, up to the scale
factor $2\pi$. Its boundary walls, where two eigenvalues collide, are the walls $\theta_i=\theta_j$ of
$\Conf(n,\R)$ from Lecture~8, wrapped around the circle.
\end{remark}

All of this generalizes.

\begin{theorem}[Maximal tori; Weyl, Cartan]\label{LIE:thm:torus}
Let $G$ be a compact connected Lie group. Then $G$ has a maximal torus $T$, unique up to conjugacy,
and every element of $G$ is conjugate into $T$. The Weyl group $W=N(T)/T$ is finite and is the Weyl
group of the root system of $\mathfrak g\otimes\C$. The kernel of $\exp\colon\mathfrak t\to T$ is a lattice
containing $2\pi$ times the coroot lattice (the coroots $h_\alpha$ placed in $\mathfrak t$ as $ih_\alpha$,
as in the $\mathrm{SU}(n)$ example), with equality exactly when $G$ is simply connected; in general the
quotient is $\pi_1(G)$. The
compact simply connected simple Lie groups correspond to the connected Dynkin diagrams:
\[
  \mathrm{SU}(n+1)\ (A_n),\quad \mathrm{Spin}(2n+1)\ (B_n),\quad \mathrm{Sp}(n)\ (C_n),\quad
  \mathrm{Spin}(2n)\ (D_n),\quad G_2,\ F_4,\ E_6,\ E_7,\ E_8 .
\]
\end{theorem}

For example, $\mathrm{SU}(2)$ has $T=\{\operatorname{diag}(e^{i\theta},e^{-i\theta})\}$ and kernel
$\theta\in2\pi\Z$, while $\SO(3)=\mathrm{SU}(2)/\{\pm I\}$ has the finer kernel $\theta\in\pi\Z$. The covering
$\mathrm{SU}(2)\to\SO(3)$ of the previous section is recorded by an index-two inclusion of lattices.

%======================================================================
\section{Summary: a dictionary}
\label{LIE:sec:dictionary}

\begin{center}
\renewcommand{\arraystretch}{1.25}
\begin{tabular}{p{0.42\textwidth}p{0.50\textwidth}}
\toprule
Packings and lattices (Lectures 4--8) & Lie theory\\
\midrule
root system $\Delta$ & nonzero eigenvalues of $\operatorname{ad}\mathfrak h$ on $\mathfrak g$\\
Cartan integer $n_{\beta\alpha}\in\Z$ & eigenvalue of the coroot $h_\alpha$; integral by $\mathfrak{sl}_2$ theory\\
reflection $s_\alpha$, Weyl group $W$ & symmetry of $\alpha$-strings; $W=N(T)/T$\\
choice of $d$, positive roots, Weyl chamber & upper triangular matrices (a Borel subalgebra)\\
simple roots, Dynkin diagram & generators $e_i,f_i,h_i$ and Serre relations\\
rank $+\abs{\Delta}$ & $\dim\mathfrak g$\\
root lattice $\Z\Delta$ (coroot lattice, in general) & $\tfrac1{2\pi}\,$kernel of $\exp$ on the maximal torus (simply connected, compact)\\
$E_8$ kissing configuration ($240$ vectors) & root spaces of $\mathfrak e_8$, $248=8+240$\\
walls $x_i=x_j$ of $\Conf(N,\R)$ & eigenvalue collisions; non-regular elements of $\mathrm{SU}(N)$\\
one-parameter groups of diffeomorphisms (Lecture 9) & $t\mapsto e^{tX}$\\
\bottomrule
\end{tabular}
\end{center}

%======================================================================
\section*{Exercises}
\addcontentsline{toc}{section}{Exercises}

\begin{exercise}
Write out the root decomposition of $\mathfrak{sl}_3$ explicitly. Check that $[E_{12},E_{23}]=E_{13}$,
$[E_{13},E_{32}]=E_{12}$ and $[E_{12},E_{13}]=0$, and interpret each identity on the hexagon of
Example~\ref{LIE:ex:sl3}.
\end{exercise}

\begin{exercise}
Prove the Jacobi identity for the commutator, and show that $\omega\mapsto X_\omega$ is an isomorphism
$(\R^3,\times)\to\mathfrak{so}_3(\R)$.
\end{exercise}

\begin{exercise}
Expand $e^{sX}e^{tY}e^{-sX}e^{-tY}$ to second order and confirm that the result is $I+st[X,Y]$ plus
higher-order terms.
\end{exercise}

\begin{exercise}
Show that $\exp\colon\mathfrak{sl}_2(\R)\to\SL_2(\R)$ is not surjective: the matrix
$\left(\begin{smallmatrix}-1&1\\0&-1\end{smallmatrix}\right)$ is not in the image. \emph{Hint:} a real
traceless $2\times2$ matrix has eigenvalues $\pm\lambda$ or $\pm i\mu$. Compute the possible traces of
$e^X$, and decide when the trace can equal $-2$.
\end{exercise}

\begin{exercise}
Using the antidiagonal form $J$, find the roots of $\mathfrak{so}_6$ and check that they form a copy of
$A_3$ (twelve roots, simply laced, rank $3$). Do the same for $\mathfrak{so}_4$ and $A_1\times A_1$, and for
$\mathfrak{sp}_4$ and $\mathfrak{so}_5$ (show that $C_2$ and $B_2$ differ by a rotation of $45^\circ$ and a scaling).
\end{exercise}

\begin{exercise}
Prove Rodrigues' formula, and show that $\exp$ maps the closed ball of radius $\pi$ in $\mathfrak{so}_3$
onto $\SO(3)$, identifying exactly the antipodal points of the boundary.
\end{exercise}

\begin{exercise}
Draw the kernel of $\exp$ on the maximal torus of $\mathrm{SU}(3)$ inside the plane
$\theta_1+\theta_2+\theta_3=0$, together with the Weyl alcove. Check that the kernel is a hexagonal
lattice and that the alcove is an equilateral triangle. How many alcoves fit into a fundamental domain
of the lattice?
\end{exercise}

\section*{Further reading}
\addcontentsline{toc}{section}{Further reading}
\begin{itemize}[nosep]
\item J.~Stillwell, \emph{Naive Lie Theory}, Springer UTM, 2008. Matrix groups first, with minimal prerequisites.
\item B.~C.~Hall, \emph{Lie Groups, Lie Algebras, and Representations}, 2nd ed., Springer GTM 222, 2015.
  Matrix groups, the exponential map, and a full treatment of the classification.
\item J.~E.~Humphreys, \emph{Introduction to Lie Algebras and Representation Theory}, Springer GTM 9, 1972.
  The algebraic route from $\mathfrak{sl}_2$ to root systems and Serre's theorem.
\item W.~Fulton and J.~Harris, \emph{Representation Theory: A First Course}, Springer GTM 129, 1991.
  Lectures 11--13 work through $\mathfrak{sl}_2$ and $\mathfrak{sl}_3$ by hand, and Lecture 21 does the classification.
\item J.~H.~Conway and N.~J.~A.~Sloane, \emph{Sphere Packings, Lattices and Groups}, Ch.~4 and 8, for the
  lattice side of $E_8$.
\end{itemize}

\endgroup
\setcounter{chapter}{6}

\part{Configuration spaces and Morse theory}
% ======================================================================
% MAT670 -- Lectures 7 and 8 (source: MAT670_78, pp. 1--31)
% First-pass transcription: "clean + fill gaps".
% ======================================================================

\chapter{Lecture 7: A short introduction to configuration spaces}
\label{L07:ch:lecture7}

\section{Configuration spaces}
\label{L07:sec:confspaces}

A \emph{configuration space} is a topological space that records the possible
configurations of an object we wish to study. For example:
\begin{itemize}
  \item the $2^N$ possible results of $N$ coin flips
    \added{(a discrete space $\{0,1\}^N$ with $2^N$ points)};
  \item the tilings of a checkerboard by dominoes \added{(again a finite set)};
  \item the positions and orientations of a polytope in space
    \added{(for a polytope with no symmetries this is $\R^3\times\SO(3)$; in general
    one divides $\SO(3)$ by the rotational symmetry group of the polytope)};
  \item the configurations of a mechanical linkage;
  \item the positions (and velocities) of particles in a box.
\end{itemize}

This last example is the one most closely related to what we want to study, and to the
first few lectures: positions of points in space.

\begin{definition}[Configuration space of $N$ points]
\label{L07:def:conf}
Let $X$ be a topological space. The \emph{(classical) configuration space of $N$
distinct labelled points in $X$} is
\[
  C^N(X) = \Conf(N,X) = \prod_{i=1}^N X \setminus \Delta,
  \qquad
  \Delta = \{(x_i)_{i=1}^N : x_i = x_j \text{ for some } i\neq j\},
\]
\added{with the subspace topology from the product $X^N$.}
\orig{crossed out: ``statistical mechanics'', ``also study''}
\end{definition}

\begin{exercise}
\label{L07:ex:weyl}
Think about how this might be related to the Weyl group associated to the root system
$A_n$.
\added{(Hint: take $X=\R$ and $N=n+1$. The Weyl group of $A_n$ is the symmetric group
$\Sigma_{n+1}$, acting on $\R^{n+1}$ by permuting coordinates, and its reflecting
hyperplanes are exactly the hyperplanes $x_i=x_j$ that make up $\Delta$. This is taken up
again in Section~\ref{L07:sec:weyl}.)}
\end{exercise}

There are also \emph{unlabelled} versions, $\Conf(N,X)/\Sigma_N$, where the symmetric
group $\Sigma_N$ acts by permuting the indices. When $X$ is endowed with a metric we can
also consider \emph{reduced} versions, $\Conf(N,X)/\operatorname{Isom}(X)$.
\added{(One can of course do both: $\Conf(N,X)/(\Sigma_N\times\operatorname{Isom}(X))$.)}
So the unlabelled version picks a single representative class for each permutation of the
indices, and the reduced version divides out by the symmetries of the ambient space.

\begin{example}[$\Conf(N,\R)$]
\label{L07:ex:confR}
We have $\Conf(N,\R)=\R^N\setminus\Delta$.
\begin{itemize}
  \item \emph{Unlabelled:} each point of $\Conf(N,\R)/\Sigma_N$ has exactly one
    representative with $x_1<x_2<x_3<\dots<x_N$, so we can choose to work in this region
    of the space.
    \added{$\R^N\setminus\Delta$ has $N!$ connected components, one for each ordering of
    the coordinates; each is an open convex cone, hence contractible, and $\Sigma_N$
    permutes them simply transitively.}
  \item \emph{Reduced:} we can shift one of the $x_i$ to $0$, for example the smallest.
    Together with the ordering this gives $\{0=x_1<x_2<\dots<x_N\}$.
    \added{(Modulo translations, the ordered region becomes the open cone
    $\{0<x_2<\dots<x_N\}\subset\R^{N-1}$; if reflections $x\mapsto -x$ are divided out as well, one may in addition
    fix the order of the two ``outer'' gaps.)}
\end{itemize}
\end{example}

\begin{center}
\begin{tikzpicture}[scale=0.9]
  % N = 2
  \fill[blue!8] (0,0) -- (2.4,2.4) -- (0,2.4) -- cycle;
  \draw[->] (-1.2,0) -- (2.6,0) node[right] {$x_1$};
  \draw[->] (0,-1.2) -- (0,2.6) node[above] {$x_2$};
  \draw[dashed] (-1.2,-1.2) -- (2.4,2.4) node[above right] {$\Delta$};
  \draw[very thick,addcol] (0,0.05) -- (0,2.4);
  \fill[white] (0,0) circle (2pt); \draw (0,0) circle (2pt);
  \node[align=center] at (0.7,1.75) {\small $x_1<x_2$};
  \node[addcol,left] at (0,1.2) {\small $x_1=0$};
\end{tikzpicture}
\end{center}
\noindent\textit{Figure:} $\Conf(2,\R)=\R^2\setminus\Delta$. The shaded half-plane
$\{x_1<x_2\}$ is a copy of the unlabelled space; the thick ray $\{0=x_1<x_2\}$ is the
reduced unlabelled space.
\orig{p.\,3 also sketches $N=3$, labelled ``$x_1=0$'', ``$x_1<x_2$'', ``$0<x_2<x_3$''}

\begin{example}[$\Conf(N,\Sph^2)$]
\label{L07:ex:confS2}
Consider labelled distinct points on a sphere. The reduced space is the space of $N$
labelled points on $\Sph^2$ modulo \emph{ambient} rotations.
For $N=4$: two configurations that differ by a rotation of the sphere are the same point of
the reduced space; a configuration and its mirror image are in general \emph{not} related by
a rotation, but they may become identified in the unlabelled space after relabelling
(e.g.\ swapping the labels $2$ and $3$).
\orig{p.\,4: sketches of 4 labelled points on spheres}
\added{(If one divides by all of $\operatorname{Isom}(\Sph^2)=\Orth(3)$ rather than by
$\SO(3)$, mirror images are identified as well.)}
\end{example}

\section{Braid groups}
\label{L07:sec:braid}

These spaces are related to the Artin braid groups, which we may define as follows.

\begin{definition}[Braid groups]
\label{L07:def:braid}
The \emph{braid group} and the \emph{pure braid group} of $X$ (on $N$ strands) are
\[
  B_N(X) = \pi_1\big(\Conf(N,X)/\Sigma_N\big),
  \qquad
  P_N(X) = \pi_1\big(\Conf(N,X)\big).
\]
\added{For $X=\R^2$ these are Artin's braid group $B_N$ and pure braid group $P_N$.
Since $\Conf(N,X)\to\Conf(N,X)/\Sigma_N$ is a covering map with deck group $\Sigma_N$, for
connected $\Conf(N,X)$ there is a short exact sequence
$1\to P_N(X)\to B_N(X)\to\Sigma_N\to 1$.}
\end{definition}

To see why these are called braid groups, consider what a path is in $\Conf(N,\R^2)$.
Fix a base point $x_0=(x_0^1,\dots,x_0^N)$. A loop $\gamma\colon[0,1]\to\Conf(N,\R^2)$
consists of $N$ paths $\gamma_1,\dots,\gamma_N$ in the plane which never collide at any
time $t$. Drawing the graphs $\{(\gamma_i(t),t)\}\subset\R^2\times[0,1]$ gives $N$
disjoint strands running from the plane $t=0$ to the plane $t=1$: a braid.
\added{In the unlabelled space a loop may end at a permutation of the starting labels,
which is why $B_N$ surjects onto $\Sigma_N$; homotopic loops give braids that are isotopic
with endpoints fixed.}

\begin{center}
\begin{tikzpicture}[scale=0.8,yscale=1.1]
  % top and bottom planes
  \draw (-0.6,0.15) -- (2.6,0.15) -- (3.0,0.65) -- (-0.2,0.65) -- cycle;
  \draw (-0.6,-3.15) -- (2.6,-3.15) -- (3.0,-2.65) -- (-0.2,-2.65) -- cycle;
  \node[left] at (-0.6,0.4) {\small $x_0$};
  \node[left] at (-0.6,-2.9) {\small $x_0$};
  % level 0: sigma_1 (strand from pos 0 over)
  \draw[thick] (1,0) .. controls (1,-0.5) and (0,-0.5) .. (0,-1);
  \draw[white,line width=5pt] (0,0) .. controls (0,-0.5) and (1,-0.5) .. (1,-1);
  \draw[thick] (0,0) .. controls (0,-0.5) and (1,-0.5) .. (1,-1);
  \draw[thick] (2,0) -- (2,-1);
  % level 1: sigma_2^{-1}
  \draw[thick] (1,-1) .. controls (1,-1.5) and (2,-1.5) .. (2,-2);
  \draw[white,line width=5pt] (2,-1) .. controls (2,-1.5) and (1,-1.5) .. (1,-2);
  \draw[thick] (2,-1) .. controls (2,-1.5) and (1,-1.5) .. (1,-2);
  \draw[thick] (0,-1) -- (0,-2);
  % level 2: sigma_1
  \draw[thick] (1,-2) .. controls (1,-2.5) and (0,-2.5) .. (0,-3);
  \draw[white,line width=5pt] (0,-2) .. controls (0,-2.5) and (1,-2.5) .. (1,-3);
  \draw[thick] (0,-2) .. controls (0,-2.5) and (1,-2.5) .. (1,-3);
  \draw[thick] (2,-2) -- (2,-3);
  \foreach \x/\l in {0/1,1/2,2/3} {
    \fill (\x,0) circle (1.6pt); \node[above] at (\x,0.05) {\scriptsize $\l$};
    \fill (\x,-3) circle (1.6pt); \node[below] at (\x,-3.05) {\scriptsize $\l$};
  }
  \draw[->] (3.6,0) -- (3.6,-3) node[midway,right] {\small $t$};
\end{tikzpicture}
\end{center}
\noindent\textit{Figure:} a loop in $\Conf(3,\R^2)/\Sigma_3$ drawn as a braid on three
strands. \orig{p.\,5: second sketch (strand graphs) illegible, omitted}

\begin{remark}
\added{For $X=\R$ the components of $\Conf(N,\R)$ are contractible
(Example~\ref{L07:ex:confR}), so all braid groups of $\R$ are trivial. For $X=\R^n$ with
$n\geq 3$ the diagonal has codimension $n\geq3$, so $\Conf(N,\R^n)$ is simply connected:
$P_N(\R^n)=1$ and $B_N(\R^n)\cong\Sigma_N$. The plane is the interesting case.}
\end{remark}

\section{Functions on configuration spaces}
\label{L07:sec:functions}

To configurations we would like to associate real-valued functions, based on what they
model. For example:
\begin{itemize}
  \item an electrical potential, or more generally a potential given as a sum of pair
    interactions \added{$E(x_1,\dots,x_N)=\sum_{i<j}\phi(x_i,x_j)$};
  \item a hard shell (hard sphere) potential
    \added{($\phi=+\infty$ if two spheres overlap, $0$ otherwise)};
\end{itemize}
which can describe different energy levels and can also constrain the system.

We may ask about
\begin{itemize}
  \item minima and maxima (local and global);
  \item critical points;
  \item generic configurations.
\end{itemize}
This breaks down into a couple of categories:
\begin{itemize}
  \item small $N$: by hand;
  \item large $N$: by statistical mechanics, or by averaging.
\end{itemize}
The classic example for large $N$ is the expected value of $N$ coin flips with values in
$\{0,1\}$, which leads to the central limit theorem.

\subsection{The Boltzmann distribution}
\label{L07:subsec:boltzmann}

So, for example, consider $N$ particles, each in one of $K$ states, with energy $E_i$
associated to the state $i\in\{1,\dots,K\}$. Let $N_i$ be the number of particles in state
$i$. The system is closed, so the total energy is constant, and
\[
  E=\sum_{i=1}^K E_i N_i, \qquad N=\sum_{i=1}^K N_i .
\]
The number of ways to place the $N$ particles according to the partition
$\{N_1,\dots,N_K\}$ is the multinomial coefficient
\[
  W=\frac{N!}{N_1!\cdots N_K!}.
\]
For large $N$ the optimal partition dominates, so we want to determine the partition that
maximises $W$ \added{subject to the two constraints above}. One can derive the
\emph{Boltzmann measure}
\[
  \frac{N_i}{N}=\frac{e^{-\beta E_i}}{\sum_{j=1}^K e^{-\beta E_j}},
  \qquad\text{i.e.}\qquad
  \Pr(\text{state } i)=\frac{e^{-\beta E_i}}{Z},
\]
where $Z=\sum_{j=1}^K e^{-\beta E_j}$ is the \emph{partition function} and
$\beta=1/(k_B T)$, with $k_B$ \added{Boltzmann's} constant and $T$ the temperature.
\orig{derivation deferred to ``next time'' in the notes}

\begin{addedblock}
Sketch of the derivation. By Stirling's formula,
$\ln W\approx N\ln N-\sum_i N_i\ln N_i$. Maximising this subject to $\sum_i N_i=N$ and
$\sum_i E_iN_i=E$ with Lagrange multipliers $\alpha,\beta$ gives
$-\ln N_i-1-\alpha-\beta E_i=0$, i.e.\ $N_i\propto e^{-\beta E_i}$; normalising gives the
formula above. The multiplier $\beta$ is fixed by the energy constraint, and physics
identifies it with $1/(k_BT)$.
\end{addedblock}

\begin{remark}
Note:
\begin{itemize}
  \item as $T\to\infty$ (so $\beta\to0$), $\Pr(i)\to 1/K$: all states become equally likely;
  \item as $T\to 0$ (so $\beta\to\infty$), $\Pr$ concentrates on the states of minimal
    energy \added{(uniformly on the ground states if there are several)}.
\end{itemize}
\end{remark}

\subsection{A gas in a box}
\label{L07:subsec:gas}

But let us consider another (sketchy) physics analysis: the distribution of a gas in a box,
in the simplest non-interacting setting. Take $2N$ points, i.e.\ a configuration in
\[
  \Conf\big(2N,[-1,1]^3\big),
\]
\added{chosen uniformly at random.} Let $P_m$ be the probability that exactly $N+m$
particles have positive first coordinate.

There are $2^{2N}$ choices for the signs of the first coordinates, all equally likely, and
$\binom{2N}{N+m}$ of them have exactly $N+m$ particles with first coordinate $>0$. (This is
really the coin-flip analysis.) So
\[
  P_m=2^{-2N}\,\frac{(2N)!}{(N+m)!\,(N-m)!}.
\]

We use a very bad version of Stirling's approximation,
\[
  n!\sim\Big(\frac ne\Big)^n\sqrt{2\pi n}\;\leadsto\;n!\approx\Big(\frac ne\Big)^n .
\]
Then
\begin{align*}
  P_m&\approx 2^{-2N}\Big(\frac{2N}{e}\Big)^{2N}\Big/
     \Big[\Big(\frac{N+m}{e}\Big)^{N+m}\Big(\frac{N-m}{e}\Big)^{N-m}\Big]
   =\frac{N^{2N}}{(N+m)^{N+m}(N-m)^{N-m}}\\
  &=\Big(1+\frac mN\Big)^{-(N+m)}\Big(1-\frac mN\Big)^{-(N-m)}
   =\Big(1-\frac{m^2}{N^2}\Big)^{-N}\Big(1+\frac mN\Big)^{-m}\Big(1-\frac mN\Big)^{m}\\
  &\overset{m\ll N}{\approx}\big(e^{-m^2/N^2}\big)^{-N}\big(e^{m/N}\big)^{-m}\big(e^{-m/N}\big)^{m}
   = e^{-m^2/N}.
\end{align*}
Since the crude Stirling formula loses the factors $\sqrt{2\pi n}$, we only keep the shape
and write
\[
  P_m\approx P_0\,e^{-m^2/N}
\]
(a crude central limit theorem). Normalising,
\[
  1=\sum_{m=-N}^{N}P_m\approx\int_{-\infty}^{\infty}P_0\,e^{-m^2/N}\,dm=P_0\sqrt{\pi N}
  \quad\Longrightarrow\quad
  P_0=\frac{1}{\sqrt{\pi N}},\qquad P_m\approx\frac{1}{\sqrt{\pi N}}\,e^{-m^2/N}.
\]
Comparing with $e^{-m^2/(2\sigma^2)}$ we get $\sigma^2=N/2$, so $\sigma=\sqrt{N/2}$.
\added{This agrees with the exact variance $2N\cdot\frac12\cdot\frac12=N/2$ of the
binomial distribution, and $P_0=1/\sqrt{\pi N}$ is exactly what the full Stirling formula
gives for $\binom{2N}{N}4^{-N}$.}

So the number of points with positive first coordinate is approximately Gaussian, centred
at the expected value $N$ with standard deviation $\sqrt{N/2}$. For example, for large $N$,
about $95\%$ of configurations are within $2\sigma=2\sqrt{N/2}$ of the expected value.
\added{The relative fluctuation $\sigma/N=1/\sqrt{2N}$ tends to $0$.}

\begin{remark}
\added{The approximations $\ln(1\pm x)\approx\pm x-x^2/2$ used above each drop cubic
terms of size $m^3/N^2$, but these cancel between the factors $(1+\frac mN)^{-m}$ and
$(1-\frac mN)^{m}$ (the exact exponent is even in $m$). The first neglected term is
$-m^4/(6N^3)$, so the Gaussian approximation is accurate for $m\ll N^{3/4}$; this covers
the typical range $m=O(\sqrt N)$.}\fixed{error order was stated as $m^3/N^2$}
\end{remark}

So if we pick a point in high dimensions, i.e.\ a configuration in
$[-1,1]^{3\cdot 2N}$\fixed{notes: $[-1,1]^{3N}$; $2N$ points in $\R^3$ give $6N$ coordinates}
(and the same argument works for the other variables), we find roughly half of the
coordinates positive and half negative, just like the coin flips. This makes sense.
We can also talk about energy, etc., and generic configurations in this way.

But we may be interested in smaller configurations, like nearest neighbours, where
Stirling's approximation is bad.
\unclear{reading of ``smaller configurations, like nearest neighbours'' uncertain}
How do we analyse such things?

\subsection{Points on the circle}
\label{L07:subsec:circle}

Let us look at points on $\Sph^1$, \added{which we take to have length $1$}. Consider
$\operatorname{Reduced}(3,\Sph^1)$, three points on the circle modulo rotation. Use the
rotation to fix one point at $0$; going around the circle, the three gaps between
consecutive points are $x$, $y$ and $1-x-y$.
\fixed{notes write $1-x+y$}
So the reduced space is described by
\[
  \{(x,y): x>0,\ y>0,\ x+y<1\},
\]
an open triangle, whose boundary edges $x=0$, $y=0$, $x+y=1$ are the collisions.
\added{(For labelled points there are two such triangles, one for each cyclic order of the
labels $1,2,3$; compare Example~\ref{L07:ex:confS1} below.)}
In general, for $N$ points we have a simplex
\added{$\{(g_1,\dots,g_N): g_i>0,\ \sum g_i=1\}$ of gaps, an open $(N-1)$-simplex
(one for each of the $(N-1)!$ cyclic orders of the labels)}.

We can attach the \emph{injectivity radius} function
\[
  \rho=\min_{i\neq j}\tfrac12\,\operatorname{geodesic}(x_i,x_j).
\]
\added{Since the smallest gap is at most $1/N\le 1/2$, $\rho$ is half the smallest gap:
$\rho=\frac12\min(x,y,1-x-y)$ for $N=3$. It has a unique maximum $\rho=1/6$ at the
equilateral configuration $x=y=1/3$; its level sets are triangles, and it fails to be
smooth along the three segments from the centre to the vertices. In general the maximum is
$1/(2N)$, attained at the regular $N$-gon.}

\begin{center}
\begin{tikzpicture}[scale=3.6]
  \draw[dashed] (0,0) -- (1,0) -- (0,1) -- cycle;
  \draw[thin,gray] (1/3,1/3) -- (0,0) (1/3,1/3) -- (1,0) (1/3,1/3) -- (0,1);
  \draw[addcol] (1/6,1/6) -- (2/3,1/6) -- (1/6,2/3) -- cycle;
  \fill (1/3,1/3) circle (0.5pt);
  \draw[gray,thin] (1/3,1/3) -- (0.62,0.62); \node[right] at (0.62,0.64) {\scriptsize $x=y=\tfrac13$};
  \node[below] at (0.5,0) {\scriptsize $y=0$};
  \node[left] at (0,0.5) {\scriptsize $x=0$};
  \node[above right] at (0.72,0.3) {\scriptsize $x+y=1$};
  \node[addcol] at (0.47,0.21) {\tiny $\rho=\frac1{12}$};
\end{tikzpicture}
\end{center}
\noindent\textit{Figure:} $\operatorname{Reduced}(3,\Sph^1)$ as an open triangle
(dashed boundary = collisions), with a level set of $\rho$ and the lines where $\rho$ is
not smooth.

\begin{remark}
\orig{p.\,13: closing remark largely illegible}
\unclear{remark ends ``\dots it's a (box--?)\dots''}
The notes end Lecture~7 with a remark about computing with the injectivity radius.
\added{The point to keep in mind for the next lecture is that $\rho$ is a minimum of smooth
functions, hence only piecewise smooth, so smooth Morse theory does not apply to it
directly.}
\end{remark}

% ======================================================================

\chapter{Lecture 8: (Smooth) Morse theory}
\label{L07:ch:lecture8}

\section{Configuration spaces on a line and Weyl chambers}
\label{L07:sec:weyl}

We started to talk about the spaces
\[
  \Conf(N,X)=\prod_{i=1}^N X\setminus\Delta,
  \qquad
  \Delta=\{(x_i)_{i=1}^N: x_i=x_j \text{ for some } i\neq j\}.
\]

\begin{example}[$\Conf(3,\Sph^1)$]
\label{L07:ex:confS1}
Fixing the first point at $0$ by a rotation, the other two points lie in
$\Sph^1\setminus\{0\}\cong(0,1)$ and must be distinct, so
\[
  \Conf(3,\Sph^1)\cong \Sph^1\times\big((0,1)^2\setminus\{x_2=x_3\}\big),
\]
\added{and the reduced space is the open square minus its diagonal: two open triangles, one
for each cyclic order (cf.\ Section~\ref{L07:subsec:circle}).}
\fixed{notes: ``$\Conf(3,\Sph^1)\approx$ square''; this is the reduced space}
\end{example}

\begin{center}
\begin{tikzpicture}[scale=1.6]
  \draw[dashed] (0,0) rectangle (1,1);
  \draw[dashed] (0,0) -- (1,1);
  \node[below] at (0.5,0) {\scriptsize $x_2\in(0,1)$};
  \node[left] at (0,0.5) {\scriptsize $x_3$};
\end{tikzpicture}
\end{center}

\begin{example}[$\Conf(N,\R)$ and the $A_{N-1}$ Weyl chambers]
Write $n=N$. Consider the affine subspaces
\[
  V=V_c=\Big\{(x_i)_{i=1}^n:\ \sum_i x_i=c\Big\}\subset\R^n .
\]
\added{Translation along $(1,\dots,1)$ splits off, $\Conf(n,\R)\cong\R\times(V_c\setminus\Delta)$.}
In these subspaces the hyperplanes $\{x_i=x_j\}\cap V$ are the boundaries of the Weyl
chambers in our construction of the $A_{n-1}$ root system in $V\subset\R^n$.
\added{Recall $A_{n-1}=\{e_i-e_j: i\neq j\}$, with $n(n-1)$ roots, simple roots
$e_i-e_{i+1}$, and Weyl group $\Sigma_n$ generated by the reflections in the hyperplanes
$x_i=x_j$ (i.e.\ swapping two coordinates). The connected components of
$V\setminus\Delta$ are exactly the $n!$ open Weyl chambers.}

For $n=3$\fixed{notes: ``$n=2$''; the example is $A_2$, $n=3$} we had the simple roots
\[
  (1,-1,0),\qquad (0,1,-1).
\]
So the root system is given by the six vectors
\[
  \pm(1,-1,0),\qquad \pm(0,1,-1),\qquad \pm(1,0,-1),
\]
generated by swapping coordinates, i.e.\ by reflecting through the hyperplanes $x_i=x_j$,
$i\neq j$. To draw them in the plane $V_3=\{x+y+z=3\}$ we translate by $(1,1,1)$.
\end{example}

\begin{center}
\begin{tikzpicture}[scale=1.25]
  \coordinate (A) at (-1.732,-1);
  \coordinate (B) at (1.732,-1);
  \coordinate (C) at (0,2);
  \draw (A) -- (B) -- (C) -- cycle;
  \node[below left] at (A) {\scriptsize $(3,0,0)$};
  \node[below right] at (B) {\scriptsize $(0,3,0)$};
  \node[above] at (C) {\scriptsize $(0,0,3)$};
  \coordinate (O) at (barycentric cs:A=1,B=1,C=1);
  % walls (medians)
  \draw[dashed] (C) -- (barycentric cs:A=1,B=1);
  \draw[dashed] (A) -- (barycentric cs:B=1,C=1);
  \draw[dashed] (B) -- (barycentric cs:A=1,C=1);
  \node[below] at (barycentric cs:A=1,B=1) {\scriptsize $x=y$};
  \node[above right] at (barycentric cs:B=1,C=1) {\scriptsize $y=z$};
  \node[above left] at (barycentric cs:A=1,C=1) {\scriptsize $x=z$};
  % roots + (1,1,1)
  \foreach \a/\b/\c in {2/0/1, 0/2/1, 1/2/0, 2/1/0, 1/0/2, 0/1/2} {
    \draw[->,addcol,thick] (O) -- (barycentric cs:A=\a,B=\b,C=\c);
    \fill[addcol] (barycentric cs:A=\a,B=\b,C=\c) circle (1.2pt);
  }
  \fill (O) circle (1.2pt);
  \node[below right] at ($(O)+(0.05,-0.02)$) {\tiny $(1,1,1)$};
  \node[addcol,left] at ($(barycentric cs:A=2,B=0,C=1)+(-0.05,0)$) {\tiny $(2,0,1)$};
\end{tikzpicture}
\end{center}
\noindent\textit{Figure:} the plane $x+y+z=3$. The dashed lines are the walls
$x_i=x_j$; they cut the plane into the six Weyl chambers (the components of
$V_3\setminus\Delta$). The arrows are the six roots of $A_2$, translated by $(1,1,1)$; they
end at the trisection points of the edges, e.g.\ $(1,1,1)+(1,-1,0)=(2,0,1)$, and each
root is perpendicular to one wall.

\subsection{Manifolds}
\label{L07:subsec:manifolds}

We will consider spaces $M$ that are Hausdorff, second countable topological spaces: the
topology separates points, and there is a countable basis for the topology. (Basically
this is to eliminate pathological examples.)

\begin{definition}[Homeomorphism]
A \emph{homeomorphism} $\varphi\colon M\to N$ is a continuous bijection with continuous
inverse. Equivalently: $\varphi$ is a bijection, and $U\subset M$ is open if and only if
$\varphi(U)\subset N$ is open.
\end{definition}

\begin{definition}[Coordinate chart]
A \emph{coordinate chart} for a space $M$ is a homeomorphism from a coordinate
neighbourhood $U$ \added{(an open subset of $M$)} onto an open subset of $\R^n$,
\[
  \varphi_U\colon U\subset M\longrightarrow V\subset\R^n .
\]
\end{definition}

\begin{center}
\begin{tikzpicture}[scale=0.85]
  \draw (0,0) ellipse (2.2 and 1.2); \node at (1.6,0.6) {$M$};
  \draw[fill=gray!15] (-0.5,0) circle (0.55); \node at (-0.5,0) {\small $U$};
  \draw[->] (-0.5,-1.0) -- (-0.5,-2.2) node[midway,left] {$\varphi_U$};
  \draw (-2.4,-3.3) -- (1.6,-3.3) -- (2.2,-2.5) -- (-1.8,-2.5) -- cycle;
  \node at (1.6,-2.8) {\small $\R^n$};
  \draw[fill=gray!15] (-0.3,-2.9) ellipse (0.6 and 0.28); \node at (-0.3,-2.9) {\small $V$};
\end{tikzpicture}
\end{center}

A smooth bijection with smooth inverse is a \emph{diffeomorphism}. So in the smooth
setting, coordinate changes should be diffeomorphisms.

A \emph{manifold} is a space $M$ with a collection of charts $\{\varphi_{U_i}\}$ (an
\emph{atlas}), indexed by a cover $\{U_i\}$ of $M$, together with \emph{transition maps}
that patch the charts together:
\[
  \tau_{ij}=\tau_{U_iU_j}=\varphi_{U_j}\circ\varphi_{U_i}^{-1}\colon
  \varphi_{U_i}(U_i\cap U_j)\to\varphi_{U_j}(U_i\cap U_j).
\]

\begin{center}
\begin{tikzpicture}[scale=0.85]
  \draw (0,0) ellipse (2.4 and 1.2);
  \draw (-0.55,0) circle (0.7); \draw (0.45,0) circle (0.7);
  \begin{scope} \clip (-0.55,0) circle (0.7); \fill[gray!35] (0.45,0) circle (0.7); \end{scope}
  \node at (-0.95,0) {\small $U_i$}; \node at (0.85,0) {\small $U_j$};
  \draw (-2,-3) circle (0.6); \node at (-2.2,-3) {\small $V_i$};
  \draw (2,-3) circle (0.6); \node at (2.2,-3) {\small $V_j$};
  \draw[->] (-0.9,-0.9) -- (-1.8,-2.3) node[midway,left] {$\varphi_{U_i}$};
  \draw[->] (0.8,-0.9) -- (1.8,-2.3) node[midway,right] {$\varphi_{U_j}$};
  \draw[->] (-1.3,-2.75) to[bend left=20] node[midway,above] {\small $\tau_{ij}$} (1.3,-2.75);
  \draw[->] (1.3,-3.25) to[bend left=20] node[midway,below] {\small $\tau_{ji}=\tau_{ij}^{-1}$} (-1.3,-3.25);
\end{tikzpicture}
\end{center}

\begin{definition}[Smooth structure]
\label{L07:def:smoothstructure}
A \emph{smooth differentiable structure} on $M$ is a collection of (smooth) coordinate
charts $\varphi_i\colon U_i\to V_i\subset\R^n$, $U_i\subset M$, such that
\begin{enumerate}
  \item $M=\bigcup_i U_i$;
  \item every transition map $\tau_{ij}=\varphi_j\circ\varphi_i^{-1}$ is smooth
    \added{(on its domain $\varphi_i(U_i\cap U_j)$)};
  \item it is a \emph{maximal} atlas: if a chart is compatible with all $\varphi_i$ as
    above, then it is in the collection.
\end{enumerate}
A topological space with a smooth structure is a \emph{smooth manifold} (of dimension $n$).
\end{definition}

Basically we can ignore maximality for computations, as maximal atlases are ``unique'':
\added{every atlas satisfying (1) and (2) is contained in exactly one maximal atlas, so it
suffices to write down any atlas.}

\section{Examples; smooth maps}
\label{L07:sec:examples}

\begin{example}
\label{L07:ex:manifolds}
\leavevmode
\begin{itemize}
  \item $\R^n$: choose the identity map as the chart \added{(one chart, $U=\R^n$)}.
  \item $\C^n\to\R^{2n}$, using $(\operatorname{Re},\operatorname{Im})$ on each $U_i$.
  \item $\Sph^n\to\R^n$: stereographic projection on $\Sph^n\setminus\{N\}$ and on
    $\Sph^n\setminus\{S\}$ \added{(north and south poles; the transition map is the
    inversion $y\mapsto y/\abs{y}^2$ on $\R^n\setminus\{0\}$, which is smooth)}.
  \item On tori, by products of the $\Sph^1$ charts.
  \item On open subsets of manifolds, by carrying over the subspace topology
    \added{and restricting the charts}.
\end{itemize}
\added{In particular $\Conf(N,X)$ is an open subset of $X^N$, so it is a smooth manifold of
dimension $N\dim X$ whenever $X$ is.}
\end{example}

We can define \emph{smooth maps} between manifolds.

\begin{definition}[Smooth map]
\label{L07:def:smoothmap}
A map $f\colon M\to N$ is \emph{smooth} if for each $p\in M$ and some charts $\varphi$ on
$U\ni p$ for $M$ and $\psi$ on $V\ni f(p)$ for $N$ (and then for all such charts),
\added{with $f(U)\subset V$,} the induced map
\[
  \begin{array}{ccc}
    U & \xrightarrow{\ f\ } & V\\[2pt]
    {\scriptstyle\varphi}\downarrow & & \downarrow{\scriptstyle\psi}\\[2pt]
    \varphi(U) & \xrightarrow{\ \psi f\varphi^{-1}\ } & \psi(V)
  \end{array}
\]
$\psi\circ f\circ\varphi^{-1}$ is smooth on its domain.
\added{(``Some $\Rightarrow$ all'' because the transition maps are smooth.)}
\end{definition}
\orig{p.\,12 (= scan p.\,25) repeats this diagram}

Similarly for maps between manifolds of other kinds. So we can define diffeomorphisms
(homeomorphisms) of manifolds: a (smooth) bijection $M\to N$ with (smooth) inverse,
written $M\cong N$.
If we change the underlying conditions on the maps \added{(e.g.\ continuous, $C^k$, or
piecewise linear transition maps)} we can still define coordinate charts, atlases, etc.
\orig{crossed out: ``Sadly, not all the maps we want to work with will be smooth''}

\section{Homotopy type}
\label{L07:sec:homotopy}

\orig{a first ``8.3 Morse theory'' is crossed out here}
We will want to talk about \emph{homotopy type} as well, since we may be dealing with
things that are not diffeomorphic but are topologically ``similar''.

\begin{definition}[Homotopy]
\label{L07:def:homotopy}
Two maps $f,g\colon M\to N$ are \emph{homotopic}, written $f\simeq g$, if there exists a
continuous function
\[
  H\colon M\times I\to N,\qquad I=[0,1],
\]
such that $H(x,0)=f(x)$ and $H(x,1)=g(x)$ for all $x\in M$.
\end{definition}

\begin{definition}[Homotopy equivalence]
\label{L07:def:hequiv}
Two (smooth) spaces $M$ and $N$ are \emph{homotopy equivalent} (have the same homotopy
type) if there are maps
\[
  f\colon M\to N,\qquad g\colon N\to M
  \qquad\text{such that}\qquad
  f\circ g\simeq\mathrm{id}_N,\quad g\circ f\simeq\mathrm{id}_M .
\]
\orig{letters $X,Y$ corrected to $M,N$ in blue on the page}
\end{definition}

The idea: the spaces can be continuously (smoothly) deformed into each other.

\begin{example}[$D^2\simeq\{0\}$]
\label{L07:ex:disk}
Let $f\colon D^2\to\{0\}$ be the radial retraction to $0$ (the constant map) and
$g\colon\{0\}\to D^2$ the inclusion. In polar coordinates put
\[
  h\big(t,(r,\theta)\big):=\big((1-t)r,\theta\big).
\]
Then $h(0,\cdot)=\mathrm{id}_{D^2}$ and $h(1,\cdot)=g\circ f$ (the map $D^2\to 0$), so
$g\circ f\simeq\mathrm{id}_{D^2}$; and $f\circ g=\mathrm{id}_{\{0\}}$.
\added{($h$ is continuous also at $r=0$, since $(1-t)\cdot 0=0$ regardless of $\theta$.)}
\end{example}

\begin{definition}[Relative homotopy; deformation retract]
\label{L07:def:relhom}
A \emph{homotopy relative to a subspace} is a homotopy which fixes the points of the
subspace: for $f,g\colon M\to N$ and $K\subset M$, $f$ and $g$ are \emph{homotopic
relative to $K$} if there is a homotopy $h\colon M\times I\to N$ between $f$ and $g$ such
that
\[
  h(k,t)=f(k)=g(k)\qquad\text{for all } k\in K,\ t\in I .
\]
If $f=\mathrm{id}_M$ and $g$ is a \emph{retraction}, i.e.\ a continuous map $M\to K$ with
$g|_K=\mathrm{id}$ \added{(followed by the inclusion $K\hookrightarrow M$)}, then this
homotopy defines a \emph{strong deformation retraction} of $M$ onto $K$.
\end{definition}
See the last example: $h$ is a strong deformation retraction of $D^2$ onto $\{0\}$.

\section{An aside about homotopy groups}
\label{L07:sec:homotopygroups}

We consider homotopy classes of maps to pointed spaces
\added{(with homotopies relative to the base point)}. Consider
$f\colon \Sph^1\to(X,x_0)$, or equivalently paths $f\colon I\to X$ with
$f(0)=f(1)=x_0$. These have a group structure under concatenation,
\[
  (f\cdot g)(t)=\begin{cases} f(2t), & t\in[0,\tfrac12],\\ g(2t-1), & t\in[\tfrac12,1].\end{cases}
\]
This is the fundamental group $\pi_1(X,x_0)$.

In general one can consider maps $f\colon\Sph^n\to(X,x_0)$, viewed as maps
$f\colon I^n\to X$ with $f(\partial I^n)=x_0$, with concatenation defined appropriately
\added{(in the first coordinate)}; this gives $\pi_n(X,x_0)$. If $X$ is path connected we
may ignore the base point $x_0$ \added{(up to isomorphism)}.

\begin{example}
Draw $(D^2,x_0)\setminus\{0\}$.
\orig{example left blank on the page}
\added{The punctured disk deformation retracts (radially) onto the circle of radius
$\abs{x_0}$, so $\pi_1(D^2\setminus\{0\},x_0)\cong\pi_1(\Sph^1)\cong\Z$, generated by
the loop going once around the puncture. Compare $\pi_1(D^2)=0$ from
Example~\ref{L07:ex:disk}: removing a point changes the homotopy type.}
\end{example}

We will not deal with this too much, but note that this is a nice invariant for spaces:
\added{homotopy equivalent spaces have isomorphic homotopy groups.} This is also the
source of the braid groups of Section~\ref{L07:sec:braid}.

\section{Morse theory}
\label{L07:sec:morse}
\orig{numbered ``8.2'' in the notes}

Morse theory is a fairly general way to study manifolds and functions on them. The idea is
that the structure of the sets of critical points of functions and the topological
structure of the manifold are related. (We will follow Milnor and others
\added{(see J.~Milnor, \emph{Morse Theory}, Annals of Math.\ Studies 51, Part~I)}.)

\begin{example}[The classic picture]
\label{L07:ex:torus}
Consider $f\colon\T^2\to\R$ given by
\[
  \T^2\hookrightarrow\R^3\xrightarrow{\ \mathrm{proj}_z\ }\R,
\]
the height function $h=\mathrm{proj}_z(x,y,z)=z$ on a torus standing upright, normalised
so that it takes values in $[0,1]$. It has four critical points: $p$ (minimum, height $0$),
$q$ and $r$ (the bottom and top of the hole) and $s$ (maximum, height $1$).
\end{example}

\begin{center}
\begin{tikzpicture}[scale=0.95]
  \draw[thick] (0,0) ellipse (1.25 and 1.9);
  \draw[thick] (-0.1,-0.75) .. controls (-0.38,-0.3) and (-0.38,0.3) .. (-0.1,0.75);
  \draw[thick] (-0.17,-0.6) .. controls (0.12,-0.25) and (0.12,0.25) .. (-0.17,0.6);
  \foreach \y/\n in {-1.9/p,-0.68/q,0.68/r,1.9/s} {
    \draw[dashed,gray] (-1.6,\y) -- (3.2,\y);
  }
  \fill (0,-1.9) circle (1.5pt) node[below left] {$p$};
  \fill (-0.13,-0.68) circle (1.5pt) node[left] {\small $q$};
  \fill (-0.13,0.68) circle (1.5pt) node[left] {\small $r$};
  \fill (0,1.9) circle (1.5pt) node[above left] {$s$};
  \node[right] at (3.2,-1.9) {\small $0$};
  \node[right] at (3.2,1.9) {\small $1$};
  \draw[->] (2.4,-2.3) -- (2.4,2.3) node[above] {\small $h=z$};
\end{tikzpicture}
\end{center}
\noindent\textit{Figure:} the torus standing on the plane $z=0$, seen from the side, with
the height function and its critical points.

Consider the \emph{sublevel sets}
\[
  M_a:=\{x\in M: f(x)<a\}.
\]
\added{(Milnor uses $M^a=\{f\le a\}$; for $a$ not a critical value the two have the same
homotopy type.)}
\begin{itemize}
  \item $a<0=f(p)$: $M_a=\emptyset$.
  \item $f(p)<a<f(q)$: $M_a$ is a $2$-cell, i.e.\ a $2$-ball (a disk).
  \item $f(q)<a<f(r)$: $M_a$ is a cylinder.
  \item $f(r)<a<f(s)$: $M_a$ is a torus minus a disk.
  \item $f(s)<a$: $M_a$ is the whole torus.
\end{itemize}

Up to homotopy, these are handle or \emph{cell attachments}.

\begin{definition}[Attaching a cell]
\label{L07:def:cell}
Let $e^k$ be a closed $k$-dimensional ball, with $\partial e^k=\Sph^{k-1}$, let $Y$ be a
topological space and $g\colon\Sph^{k-1}\to Y$ continuous. Then $Y\cup_g e^k$, the result
of \emph{attaching a $k$-cell to $Y$ by $g$}, is defined as
\[
  Y\cup_g e^k=\big(Y\sqcup e^k\big)/\sim,
  \qquad x\sim g(x)\ \text{ for } x\in\Sph^{k-1}.
\]
\end{definition}

\orig{scan pp.\,28--29 are in reverse order (circled 16, 15)}
The sublevel sets of the torus, up to homotopy, are:
\begin{itemize}
  \item a disk $\simeq$ a point ($0$-cell);
  \item the cylinder $\simeq$ a disk with a $1$-handle attached (a ``U'' shape)
    $\simeq$ a point with a $1$-cell attached, i.e.\ a circle;
  \item the torus minus a disk $\simeq$ a disk with two $1$-handles
    $\simeq$ a point with two $1$-cells attached, i.e.\ a figure eight
    $\Sph^1\vee\Sph^1$;
  \item the torus $\simeq$ the figure eight with a $2$-cell attached
    \added{along the loop $aba^{-1}b^{-1}$}.
\end{itemize}
\orig{p.\,15: sketches of the handle pictures; one equivalence marked ``?''}

\begin{center}
\begin{tikzpicture}[scale=0.85]
  % point
  \fill (0,0) circle (2pt);
  \node at (0,-1.3) {\small $\{p\}$};
  % circle
  \begin{scope}[xshift=2.4cm]
    \draw[thick] (0,0.5) circle (0.5); \fill (0,0) circle (2pt);
    \node at (0,-1.3) {\small $e^0\cup e^1$};
  \end{scope}
  % figure eight
  \begin{scope}[xshift=5.2cm]
    \draw[thick] (-0.5,0) circle (0.5); \draw[thick] (0.5,0) circle (0.5);
    \fill (0,0) circle (2pt);
    \node at (0,-1.3) {\small $e^0\cup e^1\cup e^1$};
  \end{scope}
  % torus as square
  \begin{scope}[xshift=8.4cm,yshift=-0.6cm]
    \fill[gray!15] (0,0) rectangle (1.2,1.2);
    \draw[thick,->] (0,0) -- (0.65,0); \draw[thick] (0.6,0) -- (1.2,0);
    \draw[thick,->] (0,1.2) -- (0.65,1.2); \draw[thick] (0.6,1.2) -- (1.2,1.2);
    \draw[thick,->] (0,0) -- (0,0.65); \draw[thick] (0,0.6) -- (0,1.2);
    \draw[thick,->] (1.2,0) -- (1.2,0.65); \draw[thick] (1.2,0.6) -- (1.2,1.2);
    \node[below] at (0.6,0) {\scriptsize $a$}; \node[above] at (0.6,1.2) {\scriptsize $a$};
    \node[left] at (0,0.6) {\scriptsize $b$}; \node[right] at (1.2,0.6) {\scriptsize $b$};
    \node at (0.6,-0.7) {\small $\cdots\cup e^2=\T^2$};
  \end{scope}
  \node at (1.2,0) {$\leadsto$};
  \node at (3.75,0) {$\leadsto$};
  \node at (7.0,0) {$\leadsto$};
\end{tikzpicture}
\end{center}
\noindent\textit{Figure:} homotopy types of the sublevel sets $M_a$ as $a$ passes
$p$, $q$, $r$, $s$.

The homotopy type (and the diffeomorphism type, outside of a neighbourhood) changes
exactly at the critical points of the function $f$: those points where the gradient is
zero, $\nabla f=0$.

If we look at the Hessian \added{$\Hess f_x=\big(\partial^2 f/\partial u_i\partial u_j(x)\big)$
in local coordinates $u$}, which here is non-degenerate at each critical point, we see
\added{(writing the signature as (\#positive, \#negative) eigenvalues)}:
\[
  \operatorname{sig}(\Hess f_p)=(2,0),\qquad
  \operatorname{sig}(\Hess f_q)=\operatorname{sig}(\Hess f_r)=(1,1),\qquad
  \operatorname{sig}(\Hess f_s)=(0,2).
\]
We call the negative part the \emph{index}, i.e.\ the number of decreasing directions at the
critical point, and we see that it corresponds to the dimension of the attached cell:
index $0$ at $p$ ($0$-cell), index $1$ at $q$ and $r$ ($1$-cells), index $2$ at $s$
($2$-cell).

This is why we want to start with smooth functions.

\begin{definition}[Morse function]
\label{L07:def:morse}
We will call a smooth function $f\colon M\to\R$ \emph{Morse} if it has no degenerate
critical points\fixed{notes say ``no nondegenerate critical pts''}, i.e.\ \added{at every
critical point $x$ (where $df_x=0$) the Hessian $\Hess f_x$ is a non-singular matrix}.
\end{definition}

We will formalise this next time.

\begin{addedblock}
For reference, the two basic theorems (Milnor, \emph{Morse Theory}, Thms.~3.1 and 3.2):
let $f\colon M\to\R$ be smooth and write $M^a=f^{-1}(-\infty,a]$.
(A) If $a<b$, $f^{-1}[a,b]$ is compact and contains no critical points, then $M^a$ is
diffeomorphic to $M^b$, and $M^a$ is a deformation retract of $M^b$.
(B) If $x$ is a non-degenerate critical point of index $\lambda$ with $f(x)=c$, and
$f^{-1}[c-\varepsilon,c+\varepsilon]$ is compact and contains no other critical point, then
for small $\varepsilon>0$, $M^{c+\varepsilon}$ has the homotopy type of
$M^{c-\varepsilon}\cup e^{\lambda}$.
\end{addedblock}

% ============================================================
% MAT670 -- Lectures 9 and 10 (source: MAT670_9X, pp. 1--29)
%   Lecture 9:  pp. 1--15
%   Lecture 10: pp. 16--29
% ============================================================
\chapter{Lecture 9: More on functions and their critical points}
\label{L9:chap}

Last time we defined a smooth function $f\colon M\to\R$ on a manifold $M$ to be
\emph{Morse} if it has no degenerate critical points.%
\fixed{notes say ``no non-degenerate critical points''}
\added{That is, every critical point of $f$ is non-degenerate.}

\section{What does this mean?}
\label{L9:sec:meaning}

To make sense of this definition we need to understand criticality, so we introduce
the tangent space $T_pM$ at a point $p\in M$.

\begin{center}
\begin{tikzpicture}[scale=0.8]
  \draw[thick] (0,0) ellipse (2 and 1.3);
  \node at (1.5,-0.6) {$M$};
  \draw[fill=blue!8,opacity=0.8] (-1.6,-0.1) -- (0.6,-0.5) -- (1.6,0.7) -- (-0.6,1.1) -- cycle;
  \fill (0,0.3) circle (1.5pt) node[below left] {$p$};
  \node at (1.9,1.0) {$T_pM$};
\end{tikzpicture}
\end{center}

We consider velocities of curves
\[
  \gamma\colon[-1,1]\to M,\qquad \gamma(0)=p,
\]
and, with respect to a chart $\varphi\colon U\to\R^n$ with $p\in U\subset M$, we ask
that $\varphi\circ\gamma\colon[-1,1]\to\R^n$ be differentiable at $0$.
The classes of such curves, where two curves are identified when their tangents
\added{$(\varphi\circ\gamma)'(0)$} at $0$ agree, define the \emph{tangent space} $T_pM$ at $p$.
\added{(This identification does not depend on the chart, since the transition maps
are diffeomorphisms; the tangent vector represented by $\gamma$ is written $\gamma'(0)$.)}

Given a smooth map $f\colon M\to N$, there is an induced \emph{linear} map
\[
  f_*\colon T_pM\to T_{f(p)}N,\qquad f_*\bigl(\gamma'(0)\bigr)=(f\circ\gamma)'(0),
\]
\fixed{notes: target $T_pN$}
called the \emph{differential} or \emph{push-forward} of $f$.

\begin{definition}
\label{L9:def:critical}
Let $f\colon M\to\R$ be smooth. A point $p\in M$ is a \emph{critical point} of $f$ if
\[
  f_*\colon T_pM\to T_{f(p)}\R
\]
is the zero map. The value $f(p)$ is then called a \emph{critical value}.
\end{definition}

In local coordinates $(x_1,\dots,x_n)$ \added{about $p$} this is written
\[
  \frac{\partial f}{\partial x_i}\Big|_p = 0\qquad\text{for all } i=1,\dots,n .
\]
Also, there is a bilinear form \added{on $T_pM$, well defined at a critical point $p$,} which can be
written similarly in coordinates as
\[
  H_f(p)=\Bigl(\frac{\partial^2 f}{\partial x_i\,\partial x_j}\Big|_p\Bigr)_{i,j},
  \qquad\text{also written } f_{**}(p).
\]
This is the \emph{Hessian} of $f$ at $p$.

\begin{definition}
\label{L9:def:nondeg}
A critical point $p$ of $f$ is \emph{non-degenerate} if the Hessian $H_f(p)$ is
non-degenerate, i.e.\ its matrix is non-singular. The \emph{index} of $f$ at $p$ is the
number of negative eigenvalues of $H_f(p)$ \added{(the dimension of a maximal subspace of
$T_pM$ on which $H_f(p)$ is negative definite)}.
\end{definition}

The notion of index is well defined in this case.
\added{Indeed, under a change of coordinates the Hessian at a critical point changes by
$H\mapsto J^{\mathsf T}HJ$ with $J$ the (invertible) Jacobian, and by Sylvester's law of inertia the number
of negative eigenvalues is unchanged.}

\section{Why do we care about non-degenerate critical points?}
\label{L9:sec:morselemma}
\orig{numbered 6.2--6.4 in the notes; read as 9.2--9.4}

\begin{lemma}[Morse Lemma]
\label{L9:lem:morse}
If $p$ is a non-degenerate critical point of a smooth function $f$, there is a local
coordinate system $(x_1,\dots,x_n)$ about $p$ \added{with $p$ at the origin} such that
\[
  f = \added{f(p)}\ - \sum_{i=1}^{\lambda} x_i^2 + \sum_{j=\lambda+1}^{n} x_j^2 ,
\]
where $\lambda$ is the index of $f$ at $p$.
\end{lemma}

\begin{proof}[Proof (sketch)]
We may shift so that $f(p)=0$ \added{and $p=0$ in some chart}. By the fundamental theorem of
calculus and the integral form of Taylor's theorem,
\[
  f(x_1,\dots,x_n)=\int_0^1 \frac{d}{dt}f(tx)\,dt
  =\int_0^1\sum_{i=1}^n \frac{\partial f}{\partial x_i}(tx_1,\dots,tx_n)\,x_i\,dt ,
\]
\added{so $f=\sum_i x_i g_i$ with $g_i$ smooth and $g_i(0)=\frac{\partial f}{\partial x_i}(0)=0$.
Applying the same argument to each $g_i$,}
and in fact
\[
  f(x_1,\dots,x_n)=\sum_{i,j} h_{ij}(x_1,\dots,x_n)\,x_ix_j
\]
for some smooth functions $h_{ij}$. We can symmetrize ($h_{ij}=h_{ji}$, replacing $h_{ij}$
by $\tfrac12(h_{ij}+h_{ji})$) and compute
\[
  h_{ij}(0)=\frac12\,\frac{\partial^2 f}{\partial x_i\,\partial x_j}(0)
\]
\fixed{notes: $\partial f$ for $\partial^2 f$}
in some coordinate system for $T_pM$.
This can be inductively diagonalized in a neighbourhood of $p$, for example via the
spectral theorem plus a continuity argument.

\emph{Example in dimension 2.} Write
\[
  f(x,y)=x^2h_{11}+2xy\,h_{12}+y^2h_{22},
  \qquad \frac{\partial^2 f}{\partial x^2}\Big|_{(0,0)}=2h_{11}(0,0),\ \text{etc.}
\]
By a \added{linear} coordinate change we may assume $h_{11}(0,0)\neq0$; by continuity there is
a neighbourhood on which $h_{11}\neq0$. Then we can define a local change of coordinates
\[
  \bar x=\sqrt{\abs{h_{11}}}\Bigl(x+\frac{h_{12}}{h_{11}}\,y\Bigr)
  \quad\Longrightarrow\quad
  \bar x^2=\abs{h_{11}}\Bigl(x^2+2\frac{h_{12}}{h_{11}}xy+\frac{h_{12}^2}{h_{11}^2}y^2\Bigr).
\]
In the case $h_{11}>0$,
\[
  \bar x^2=h_{11}x^2+2h_{12}xy+\frac{h_{12}^2}{h_{11}}y^2 ,
\]
and since
\[
  f=\underbrace{x^2h_{11}+2xy\,h_{12}}_{=\ \bar x^2-\frac{h_{12}^2}{h_{11}}y^2}+\,y^2h_{22},
  \qquad\text{we get}\qquad
  f=\bar x^2+\Bigl(h_{22}-\frac{h_{12}^2}{h_{11}}\Bigr)y^2 .
\]
\added{Non-degeneracy at $0$ means $h_{11}h_{22}-h_{12}^2\neq0$ there, so the coefficient of $y^2$
is non-zero near $0$, and $\bar y=\sqrt{\abs{h_{22}-h_{12}^2/h_{11}}}\;y$ finishes the job.
(The Jacobian of $(x,y)\mapsto(\bar x,\bar y)$ at $0$ is invertible, so this is a coordinate
change by the inverse function theorem.)}
The proof proceeds the same way for the other cases and for higher dimensions.
\begin{addedblock}
In general, after a linear change we have $h_{11}(0)\ne0$; put
$\bar x_1=\sqrt{\abs{h_{11}}}\,\bigl(x_1+\sum_{j>1}x_jh_{1j}/h_{11}\bigr)$, so that
$f=\pm\bar x_1^2+\sum_{i,j\ge2}x_ix_j\,h'_{ij}$ with $h'_{ij}$ smooth and symmetric and
$(h'_{ij}(0))$ non-singular; induct on $n$. (See Milnor, \emph{Morse Theory}, Lemma 2.2.)
\end{addedblock}
\end{proof}

\begin{corollary}
\label{L9:cor:isolated}
Non-degenerate critical points are isolated.
\end{corollary}
\begin{addedblock}
Indeed, in Morse coordinates $df=(-2x_1,\dots,-2x_\lambda,2x_{\lambda+1},\dots,2x_n)$ vanishes
only at the origin.
\end{addedblock}

\section{One-parameter groups of diffeomorphisms}
\label{L9:sec:oneparam}

\begin{definition}
\label{L9:def:oneparam}
A \emph{one-parameter group of diffeomorphisms} of a manifold $M$ is a smooth map
\[
  \Phi\colon\R\times M\to M
\]
such that
\begin{enumerate}[nosep]
  \item for each $t\in\R$, the map $\varphi_t\colon M\to M$ defined as
        $\varphi_t(p)=\Phi(t,p)$ is a diffeomorphism of $M$ onto $M$, and
  \item for all $t,s\in\R$, $\varphi_{t+s}=\varphi_t\circ\varphi_s$.
\end{enumerate}
\end{definition}
\orig{$\Phi$ vs.\ $\varphi_t$ distinction added in blue on the page}

This defines a vector field $X$ on $M$: for a smooth $f\colon M\to\R$, $X$ is characterized by
\[
  X_q(f)=\lim_{h\to0}\frac{f(\varphi_h(q))-f(q)}{h}.
\]
So this is an assignment of a direction to each point $q$ of $M$, and $X_q$ takes a
directional derivative of $f$. We say that $X$ \emph{generates} the group $\varphi$.

\begin{lemma}
\label{L9:lem:flow}
A smooth vector field on $M$ which vanishes outside of a compact set $K\subset M$ generates
a unique one-parameter group of diffeomorphisms of $M$.
\end{lemma}

\begin{proof}
Consider a smooth curve $\gamma\colon t\mapsto\gamma(t)\in M$ with velocity vector
$\frac{d\gamma}{dt}\in T_{\gamma(t)}M$, characterized by
\[
  \frac{d\gamma}{dt}f=\lim_{h\to0}\frac{f(\gamma(t+h))-f(\gamma(t))}{h}.
\]
Let $\varphi$ be a one-parameter group of diffeomorphisms generated by a vector field $X$ as in
the lemma. Then for fixed $q$, the curve $t\mapsto\varphi_t(q)$ satisfies
\begin{equation}\label{L9:eq:ode}
  \frac{d\varphi_t(q)}{dt}=X_{\varphi_t(q)},\qquad \varphi_0(q)=q ,
\end{equation}
since
\[
  \frac{d\varphi_t(q)}{dt}(f)=\lim_{h\to0}\frac{f(\varphi_{t+h}(q))-f(\varphi_t(q))}{h}
  =\lim_{h\to0}\frac{f\bigl(\varphi_h(\varphi_t(q))\bigr)-f(\varphi_t(q))}{h}
  =X_{\varphi_t(q)}(f).
\]
\added{This already gives uniqueness: any group generated by $X$ consists of solutions of
\eqref{L9:eq:ode}, which are unique.}

\emph{Existence.} By existence and uniqueness of solutions of ODEs (Picard--Lindel\"of, applied in
local coordinates), the system \eqref{L9:eq:ode} has a unique solution that depends smoothly on
the initial conditions. Hence for every $q\in M$ there is a neighbourhood $U$ of $q$ and
an $\varepsilon>0$ such that \eqref{L9:eq:ode} has a unique smooth solution for initial points in
$U$ and $\abs{t}<\varepsilon$.

Since $K$ is compact, \added{finitely many such neighbourhoods cover $K$, so} there is a minimal
$\varepsilon>0$ for such a cover of $K$. For $q\notin K$, let $\varphi_t(q)=q$
\added{(this solves \eqref{L9:eq:ode} since $X$ vanishes there)}. Hence there is a unique smooth
solution $\varphi_t(q)$, $\abs{t}<\varepsilon$, for all $q\in M$.
Note that
\[
  \varphi_{t+s}=\varphi_t\circ\varphi_s\qquad\text{for }\abs{t},\abs{s},\abs{t+s}<\varepsilon
\]
\added{(both sides solve \eqref{L9:eq:ode} in $t$ with the same initial value $\varphi_s(q)$)}.
So $\varphi_t$ is a diffeomorphism for $\abs{t}<\varepsilon$ \added{with inverse $\varphi_{-t}$}.

But we can iterate this process for $t\ge\varepsilon$:
\unclear{``by LNP\dots''; step size $\varepsilon$ vs.\ $\varepsilon/2$ unclear}
writing $t=k\,\tfrac{\varepsilon}{2}+r$ with $k=\lfloor 2t/\varepsilon\rfloor$ and
$0\le r<\tfrac{\varepsilon}{2}$, \added{define
$\varphi_t=\varphi_{\varepsilon/2}\circ\cdots\circ\varphi_{\varepsilon/2}\circ\varphi_r$ ($k$ factors),
and $\varphi_{-t}=(\varphi_t)^{-1}$. One checks that this is well defined, smooth, and satisfies the group law.}
\end{proof}

\section{The diffeomorphism theorem}
\label{L9:sec:diffeo}
\orig{title ``Homotopy Theorem'' crossed out, replaced by ``Diffeomorphism''}

For $f\colon M\to\R$ we have the \emph{sublevel sets}
\[
  M^a:=f^{-1}(-\infty,a]=\{p\in M: f(p)\le a\}.
\]
\orig{notes mix $M_a$ and $M^a$}

\begin{theorem}
\label{L9:thm:diffeo}
Let $f\colon M\to\R$ be a smooth function, let $a<b$, and suppose $f^{-1}[a,b]$ is compact and
contains no critical points of $f$. Then $M^a$ is diffeomorphic to $M^b$.
\added{Moreover, $M^a$ is a deformation retract of $M^b$.}
\end{theorem}

\emph{Idea.} Push $f^{-1}(b)$ down to $f^{-1}(a)$ along $-\nabla f$.

\begin{center}
\begin{tikzpicture}[scale=0.9]
  % lower cap region M^a
  \draw[thick] (-2,0) arc (180:360:2 and 1.3);
  \draw[thick] (-2,0) -- (-2,2.5);
  \draw[thick] (2,0) -- (2,2.5);
  \draw[thick] (0,2.5) ellipse (2 and 0.4);
  \draw (-2,0) arc (180:360:2 and 0.4);
  \draw[dashed] (2,0) arc (0:180:2 and 0.4);
  \draw[->,decorate,decoration={snake,amplitude=1pt,segment length=6pt}] (-0.6,2.3) -- (-0.6,-0.2);
  \draw[->,decorate,decoration={snake,amplitude=1pt,segment length=6pt}] (0.6,2.3) -- (0.6,-0.2);
  \node[left] at (-2.1,2.5) {$f^{-1}(b)$};
  \node[left] at (-2.1,0) {$f^{-1}(a)$};
  \node at (0,1.2) {\small $-\nabla f$};
  \draw[|-|] (3,-1.3) -- (3,0) node[midway,right] {$M^a$};
  \draw[|-|] (4,-1.3) -- (4,2.5) node[midway,right] {$M^b$};
\end{tikzpicture}
\end{center}

\begin{proof}
$M$ admits a \added{Riemannian} metric; we need it to define $\langle X,Y\rangle$, the inner product
of tangent vectors. So we can define the gradient $\nabla f$, which is the vector field
characterized by
\[
  \langle X,\nabla f\rangle = Xf
\]
for any vector field $X$; that is, $\langle X,\nabla f\rangle$ is the directional derivative of $f$
along $X$. Note that $\nabla f$ vanishes exactly at the critical points of $f$.
If $\gamma\colon\R\to M$ is a curve with velocity vector $\frac{d\gamma}{dt}$, then
\fixed{notes: $\gamma\colon\R\to\R$}
\[
  \Bigl\langle\frac{d\gamma}{dt},\nabla f\Bigr\rangle=\frac{d(f\circ\gamma)}{dt}.
\]

Define a smooth function $\rho\colon M\to\R$ that is
\[
  \rho=\frac{1}{\langle\nabla f,\nabla f\rangle}\quad\text{on } f^{-1}[a,b],
\]
and $0$ outside of a compact neighbourhood of $f^{-1}[a,b]$
\added{(chosen to contain no critical points, which is possible since the critical set is closed
and $f^{-1}[a,b]$ is compact)}.
Then the vector field $X$ defined as
\[
  X_q=\rho(q)\,(\nabla f)_q
\]
satisfies the hypotheses of Lemma~\ref{L9:lem:flow} (it is smooth and vanishes outside a compact
set), so there is a one-parameter group of diffeomorphisms $\varphi_t\colon M\to M$ generated by $X$.

For fixed $q$, consider $t\mapsto f(\varphi_t(q))$. If $\varphi_t(q)\in f^{-1}[a,b]$, then
\[
  \frac{d f(\varphi_t(q))}{dt}=\Bigl\langle\frac{d\varphi_t(q)}{dt},\nabla f\Bigr\rangle
  =\langle X,\nabla f\rangle = 1 .
\]
So $t\mapsto f(\varphi_t(q))$ is linear with derivative $1$ while $f(\varphi_t(q))$ lies in $[a,b]$.
Therefore $\varphi_{b-a}\colon M\to M$ is a diffeomorphism which carries $M^a$
\added{onto} $M^b$.

\begin{center}
\begin{tikzpicture}[scale=0.8]
  \draw[thick] (-1.5,0) -- (-1.5,3); \draw[thick] (1.5,0) -- (1.5,3);
  \draw[thick] (0,3) ellipse (1.5 and 0.3);
  \draw (0,1.5) ellipse (1.5 and 0.3);
  \draw[thick] (-1.5,0) arc (180:360:1.5 and 0.3);
  \draw[dashed] (1.5,0) arc (0:180:1.5 and 0.3);
  \draw[-{Latex}] (0,-0.3) -- (0,1.2); \draw[-{Latex}] (0,1.2) -- (0,2.7);
  \fill (0,-0.3) circle (1.5pt);
  \node[right] at (1.6,3) {$f=b$}; \node[right] at (1.6,0) {$f=a$};
  \node[left] at (-1.6,1.5) {$\varphi_t$};
  \node at (4.2,1.5) {\small $\varphi_{b-a}\colon M^a\xrightarrow{\ \cong\ }M^b$};
\end{tikzpicture}
\end{center}
\orig{p.\,15: sketch of this flow}

Furthermore, define $r_t\colon M^b\to M^b$, $0\le t\le 1$, by
\fixed{notes: $r_t\colon M_b\to M_a$}
\[
  r_t(q)=\begin{cases}
    q & \text{if } f(q)\le a,\\
    \varphi_{t(a-f(q))}(q) & \text{if } a\le f(q)\le b.
  \end{cases}
\]
Then $r_0=\mathrm{id}$, and $r_1$ is a retraction from $M^b$ to $M^a$; \added{$r_t$ is continuous
since the two formulas agree on $f^{-1}(a)$.} Hence $r_t$ is a deformation retraction
\added{of $M^b$ onto $M^a$.}
\added{(Cf.\ Milnor, \emph{Morse Theory}, Theorem 3.1.)}
\end{proof}

% ============================================================
\chapter{Lecture 10: Configuration spaces and Morse-like results}
\label{L10:chap}
\orig{p.\,16, boxed: ``need to be rewritten, out of order / missing pieces''}

\section{Recap and the handle attachment theorem}
\label{L10:sec:recap}

For smooth functions we have the diffeomorphism result (Theorem~\ref{L9:thm:diffeo}) that says:
if $f\colon M\to\R$, $a<b$, $f^{-1}[a,b]$ is compact, and there are no critical values in $[a,b]$,
then $M^a$ is diffeomorphic to $M^b$.

There is also a handle attachment theorem (not proved here) for smooth $f\colon M\to\R$:

\begin{theorem}[Handle attachment]
\label{L10:thm:handle}
Let $p$ be a non-degenerate critical point of $f$ with index $\lambda$ and $f(p)=c$.
Suppose there is $\varepsilon>0$ such that $f^{-1}[c-\varepsilon,c+\varepsilon]$ is compact and
contains no critical points other than $p$. Then $M^{c+\varepsilon}$ has the homotopy type of
$M^{c-\varepsilon}\cup_g e^\lambda$, \added{that is, $M^{c-\varepsilon}$ with a $\lambda$-cell attached
along a map $g\colon\partial e^\lambda\to M^{c-\varepsilon}$}.
\end{theorem}
\added{(Milnor, \emph{Morse Theory}, Theorem 3.2.)}

\subsection*{Morse-type inequalities}
Let $c_\lambda$ be the number of critical points of index $\lambda$ \added{of a Morse function on a
compact manifold $M$}. The homotopy type changes by attaching various cells of index $\lambda$ as we
``climb'' $M$ via the Morse function. Consequently
\[
  \sum_\lambda(-1)^\lambda c_\lambda=\chi(M),
\]
and in general we have the constraints
\[
  c_\lambda\ \ge\ b_\lambda(M),
\]
where $b_\lambda(M)$ is the $\lambda$-th Betti number of $M$ (the rank of the $\lambda$-th homology group).
\added{(These are the weak Morse inequalities.)}
The key point for us is that the Betti numbers can be computed for many spaces we like.

\section{Betti numbers for the configuration space of points on \texorpdfstring{$\Sph^2$}{S2}}
\label{L10:sec:betti}

Let $B\Conf(N)$ denote the configuration space of $N$ points on $\Sph^2$
\added{modulo rotations, $B\Conf(N)=\Conf(N,\Sph^2)/\SO(3)$.}
\unclear{meaning of ``$B$'' not on page; quotient by $\SO(3)$ inferred}
Via stereographic projection \added{from the first point}, for $N\ge3$ we have a
\added{diffeomorphism}
\[
  B\Conf(N)\ \xrightarrow{\ \cong\ }\ \Conf(N-1,\R^2)/\SO(2).
\]
\begin{addedblock}
Indeed, $\SO(3)$ acts transitively on $\Sph^2$, so we may rotate $u_1$ to the north pole; the
remaining freedom is the stabilizer $\SO(2)$ of the pole, and stereographic projection from the
pole identifies $\Sph^2\setminus\{u_1\}$ with $\R^2$, $\SO(2)$-equivariantly (rotations about the
origin). For $N\ge3$ the actions are free.
\end{addedblock}
That allows us to compute Betti numbers from known Betti numbers \added{via the
fibration $\SO(2)\to\Conf(N-1,\R^2)\to\Conf(N-1,\R^2)/\SO(2)$}: the Poincar\'e polynomial is
\[
  P\bigl(B\Conf(N)\bigr)=\frac{P\bigl(\Conf(N-1,\R^2)\bigr)}{P(\SO(2))},
\]
where
\[
  P\bigl(\Conf(N-1,\R^2)\bigr)=(1+t)(1+2t)\cdots\bigl(1+(N-2)t\bigr),\qquad P(\SO(2))=1+t .
\]
\begin{addedblock}
Here $P(X)(t)=\sum_i b_i(X)\,t^i$. The formula for $\Conf(n,\R^2)$ is due to Arnold (1969). The
division is justified because the $\SO(2)$-bundle is trivial: the map
$u\mapsto (u_2-u_1)/\abs{u_2-u_1}\in \Sph^1$ is $\SO(2)$-equivariant, so
$\Conf(N-1,\R^2)\cong \SO(2)\times\bigl(\Conf(N-1,\R^2)/\SO(2)\bigr)$, and the K\"unneth formula
applies. Hence
\[
  P\bigl(B\Conf(N)\bigr)=(1+2t)(1+3t)\cdots\bigl(1+(N-2)t\bigr).
\]
\end{addedblock}

The picture behind the formula for $\Conf(n,\R^2)$ is that of punctured planes, and the induction
up \added{in $n$: forgetting the last point gives a fibration
$\Conf(n,\R^2)\to\Conf(n-1,\R^2)$ whose fibre is the plane with $n-1$ punctures, with Poincar\'e
polynomial $1+(n-1)t$ (Fadell--Neuwirth)}.

So the Betti numbers $b_\lambda$ of $B\Conf(N)$ are:
\begin{center}
\begin{tabular}{c|cccc}
\toprule
 & $\lambda=0$ & $\lambda=1$ & $\lambda=2$ & $\lambda=3$ \\
\midrule
$N=3$ & $1$ & $0$ & $0$ & $0$\\
$N=4$ & $1$ & $2$ & $0$ & $0$\\
$N=5$ & $1$ & $5$ & $6$ & $0$\\
\added{$N=6$} & \added{$1$} & \added{$9$} & \added{$26$} & \added{$24$}\\
$\vdots$ & $1$ & & & \\
\bottomrule
\end{tabular}
\end{center}
In general $b_0=1$, and the top Betti number is $b_{N-3}=(N-2)!$.

One can use this idea to describe extensions and induction arguments for configuration spaces in
general; for example the $K(\pi,1)$ type of configuration spaces comes from the fact that such spaces
really change nicely \added{(the Fadell--Neuwirth fibrations: inductively, $\Conf(n,\R^2)$ is an
iterated fibration with aspherical fibres, so it is aspherical, with fundamental group the pure
braid group $P_n$)}.
\orig{some crossed-out words here}

\section{Min-type functions}
\label{L10:sec:mintype}
\orig{numbered 10.2 again in the notes}

\begin{definition}
\label{L10:def:mintype}
A \emph{min-type function} is given by a parametrized family of functions
\[
  f(x,p)\colon X\times M\to\R ,
\]
where $X$ is a parameter space, compact\dots; for our purposes, finite and discrete. We require
$f(x,\cdot)$ to be smooth (or smooth enough) for each parameter $x$. The min-type function is then
\[
  F(p)=\added{\min_{x\in X}f(x,p).}
\]
\end{definition}

Our function, the injectivity radius
\[
  \rho\colon\Conf(N)\to\R,\qquad \rho(u)=\min_{\substack{i\ne j\\ u_i,u_j\in u}} d(u_i,u_j),
\]
is of this type \added{(with $X$ the set of pairs $\{i,j\}$ and $f(\{i,j\},u)=d(u_i,u_j)$, the
geodesic distance on $\Sph^2$)}.
\added{(Here ``injectivity radius'' means the minimal distance; the largest radius of disjoint
spherical caps centred at the $u_i$ is $\rho(u)/2$.)}

In general, topological regularity is characterized by the non-emptiness of an intersection of
gradient cones of this type.
\begin{addedblock}
Precisely: for $p\in M$ let $X(p)=\{x\in X: f(x,p)=F(p)\}$ be the active parameters, and let
\[
  C_p=\{v\in T_pM:\ df_x|_p(v)>0\ \text{for all } x\in X(p)\}
\]
be the cone of directions that increase every active function, hence increase $F$ to first order.
Where $C_p\ne\emptyset$, a gradient-like vector field can be built locally, and the flow argument of
Theorem~\ref{L9:thm:diffeo} goes through for superlevel sets $\{F\ge c\}$
(cf.\ Baryshnikov--Bubenik--Kahle, \emph{Min-type Morse theory for configuration spaces of hard
spheres}, IMRN 2014).
\end{addedblock}

\begin{corollary}
\label{L10:cor:topchange}
The topology \added{of the superlevel sets of $F$} only changes at critical points $x$, which only
occur when the gradient cone is empty, i.e.\ when
\[
  \sum_{x} w_x\,df_x\big|_u=0\qquad\text{for some } w_x\ge0 \added{\text{ not all zero, } x\in X(u)}.
\]
\end{corollary}
\added{The equivalence of ``$C_u=\emptyset$'' and the existence of such weights is Gordan's theorem
of the alternative (a finite-dimensional Farkas/separation lemma).}

From our picture for the injectivity radius, this is when increasing some collection
\added{of active distances simultaneously} is not possible.
\unclear{p.\,21 right column only partly legible}

\begin{center}
\begin{tikzpicture}[scale=1]
  % regular: three points in a box, can all move apart
  \begin{scope}
    \draw[rounded corners] (-1.5,-1.1) rectangle (1.5,1.5);
    \coordinate (A) at (90:0.8); \coordinate (B) at (210:0.8); \coordinate (C) at (330:0.8);
    \draw (A) -- (B) -- (C) -- cycle;
    \foreach \P/\a in {A/90,B/210,C/330} { \fill (\P) circle (1.5pt); \draw[->,blue] (\P) -- ++(\a:0.5); }
    \node[below,align=center] at (0,-1.15) {\small regular:\\ \small all can move apart};
  \end{scope}
  % critical: three gradient directions at 120 degrees
  \begin{scope}[xshift=5.5cm,yshift=0.2cm]
    \fill (0,0) circle (1.5pt);
    \foreach \a in {90,210,330} \draw[->,thick] (0,0) -- (\a:1);
    \draw[->,dashed,gray] (0,0) -- (20:0.9) node[right] {$v$?};
    \node[below,align=center] at (0,-1.35) {\small critical: no $v$ in\\ \small the positive (dual) cone};
  \end{scope}
\end{tikzpicture}
\end{center}

\subsection{The contact graph and balanced configurations}
\orig{p.\,22 left column crossed out}

Let $u\in\Conf(N)$ with $\rho(u)\ge\theta$ \added{(we will take $\rho(u)=\theta$)}. There is a
\emph{contact graph}:
\begin{itemize}[nosep]
  \item vertices: the points $u_i\in u$;
  \item edges: $u_iu_j$ whenever $d(u_i,u_j)=\theta$.
\end{itemize}
Edges may be assigned positive weights $w_{ij}$, and if the resulting weighted graph is
\emph{force balanced}, this is equivalent to the cone being empty.
\begin{addedblock}
Write $e_{ij}\in T_{u_i}\Sph^2$ for the unit tangent vector at $u_i$ pointing along the geodesic
towards $u_j$. The differential of $d(u_i,u_j)$ has component $-e_{ij}$ at $u_i$ and $-e_{ji}$ at
$u_j$. So the condition of Corollary~\ref{L10:cor:topchange} reads: there are weights $w_{ij}\ge0$,
not all zero, on the edges of the contact graph with
\[
  \sum_{j\sim i} w_{ij}\,e_{ij}=0\qquad\text{at every vertex } u_i ,
\]
i.e.\ the ``forces'' along the edges balance at each point. The edges with $w_{ij}>0$ form the
\emph{balanced (stressed) subgraph}.
\end{addedblock}

So the only obstacles to topological regularity are exactly the force balanced configurations of
points.

\emph{Examples\dots}
\orig{left blank on p.\,23}

\begin{definition}
\label{L10:def:conftheta}
\added{$\Conf(N,\theta)=\{u\in\Conf(N,\Sph^2):\rho(u)\ge\theta\}$, and $\Conf(N,0):=\Conf(N,\Sph^2)$.}
\end{definition}

\begin{theorem}
\label{L10:thm:smalltheta}
Suppose $N\ge3$.
\unclear{``$N\ge3$'' or ``$N>3$''}
Then for $0<\theta<\pi/N$, $\Conf(N,\theta)$ is a strong deformation retract of, and hence homotopy
equivalent to, $\Conf(N,0)$.
\end{theorem}

\begin{proof}[Proof (sketch)]
First, there are no balanced configurations \added{with $\rho(u)=\theta'$ for $0<\theta'<\pi/N$};
then push each point apart \added{using the flow of a gradient-like vector field for $\rho$}\dots
\begin{addedblock}
\emph{No balanced configurations.} (i) A balanced subgraph is not contained in any open hemisphere:
if all its vertices lie in $\{y:\langle y,e\rangle>0\}$, take a vertex $v$ of the subgraph minimizing
$\langle v,e\rangle$. For each neighbour $u$, the vector $e_{vu}$ is a positive multiple of
$u-\langle u,v\rangle v$, whose $e$-component is $\langle u,e\rangle-\langle u,v\rangle\langle v,e\rangle
\ge\langle v,e\rangle(1-\langle u,v\rangle)>0$ (using $\langle u,e\rangle\ge\langle v,e\rangle$ and $\langle u,v\rangle<1$). So every force at $v$ has a positive $e$-component,
and they cannot sum to zero.
(ii) Each connected component of the balanced subgraph is itself balanced. If a component has $k$
vertices, a spanning tree has $k-1$ edges of length $\theta'$, total length $L=(k-1)\theta'$. A tree
has a centre $c$ within distance $L/2$ (along the tree) of all its points, so all vertices lie in the
spherical ball of radius $L/2$ about $c$. By (i), $L/2\ge\pi/2$, so $(k-1)\theta'\ge\pi$, so
$\theta'\ge\pi/(k-1)\ge\pi/(N-1)>\pi/N$.
(So this crude argument even gives the theorem for $\theta<\pi/(N-1)$. The sharp bound is
$\theta<2\pi/N$: a balanced graph has every valence $\ge2$, hence at least as many edges as
vertices, and Kusner--Kusner--Lagarias--Shlosman (Thm.~4.14, cited below) show its total length is at least
$2\pi$, with equality only for the $N$ points equally spaced on a great circle. Compare the
statement at the start of Lecture~11.)

\emph{Deformation.} So $\rho$ has no critical values in $(0,\theta]$, and the min-type version of
Theorem~\ref{L9:thm:diffeo} lets us push $\{\rho\ge\eta\}$ onto $\{\rho\ge\theta\}$ for every
$0<\eta<\theta$; some care is needed near $\rho=0$, where $\Conf(N,0)$ is not compact. (Cf.\
Kusner--Kusner--Lagarias--Shlosman, \emph{Configuration spaces of equal spheres touching a given
sphere: the twelve spheres problem}, arXiv:1611.10297 (2016); in \emph{New Trends in Intuitive
Geometry}, Bolyai Soc.\ Math.\ Stud.\ 27, Springer, 2018.)
\end{addedblock}
\end{proof}

For $N=4$ \dots, $\Conf(N,\theta)$ for small $\theta$ is therefore path connected.
\unclear{``$N=4$'' or ``$N>4$''}
\added{(Indeed $\Conf(N,\Sph^2)$ is connected, since the diagonals removed have real codimension~2.)}

\begin{example}[$N=4$]
\label{L10:ex:N4}
\orig{left blank in the notes}
\added{Balanced subgraphs need every vertex of valence $\ge2$, and a vertex of valence $2$ must have
its two edges in opposite directions, so it lies on a great circle through its two neighbours. One
finds: the square on a great circle ($\theta=\pi/2$), and the regular tetrahedron
($\theta=\arccos(-1/3)\approx109.47^\circ$, the maximum of $\rho$). An equilateral triangle on a great
circle ($\theta=2\pi/3$) plus a free fourth point is \emph{not} realizable, since every point of
$\Sph^2$ is within $\pi/2$ of one of the three. This is the editor's check; compare the question marks on p.\,29.}
\end{example}

\subsection{Outlook}
\orig{pp.\,24--25: headings only}
\begin{itemize}[nosep]
  \item Strata for Tammes maxima: table\dots
  \added{(The Tammes problem asks for $\max\rho$ over $\Conf(N,\Sph^2)$; it is solved for
  $N\le14$ and $N=24$, with $N=13,14$ due to Musin--Tarasov, 2012/2015.)}
  \item An algorithm for searching \added{for critical configurations} by evolution of force
  configurations\dots
  \item Local maxima; rings\dots
  \item Rigidity?, critical conditions and weights\dots
  \item Candidate graph sorting\dots
\end{itemize}

\section{Topological regularity revisited}
\label{L10:sec:revisit}
\orig{pp.\,26--29 seem to re-explain \S\ref{L10:sec:mintype}; possibly a later lecture}

What is this? This is a force balanced \added{configuration}.
\unclear{p.\,26 has only these two lines}

That is, for the injectivity radius, \added{the gradient cone} is the intersection of the convex
cones of motions that increase the distance between each \added{active} pair of points.

\emph{Picture: three points.} The configuration space of $3$ points on $\Sph^2$ is $6$-dimensional.
Fixing one point leaves a $4$-dimensional space; fixing also the rotation angle $\psi$ about it
leaves a $3$-dimensional space, which is \added{locally} parametrized by the distances
$d_{12},d_{13},d_{23}$. The cone for $d_{12}$ is the upper half space $\{d_{12}\text{ increasing}\}$,
etc.
\orig{angle called $\theta$ in the notes}

\begin{center}
\begin{tikzpicture}[scale=1]
  \begin{scope}
    \draw (0,0) circle (1);
    \fill (-0.3,0.15) circle (1.2pt); \fill (0.35,-0.05) circle (1.2pt); \fill (0.1,0.5) circle (1.2pt);
    \node at (0,-1.3) {\small 3 points on $\Sph^2$};
  \end{scope}
  \begin{scope}[xshift=3.5cm]
    \draw[dotted] (0,-1.3) -- (0,1.3);
    \coordinate (P3) at (0,-0.6); \coordinate (P1) at (-0.8,0.6); \coordinate (P2) at (0.8,0.6);
    \draw (P1) -- node[above]{\small $d_{12}$} (P2) -- node[right]{\small $d_{23}$} (P3)
          -- node[left]{\small $d_{13}$} cycle;
    \fill (P1) circle (1.3pt); \fill (P2) circle (1.3pt); \fill (P3) circle (2pt);
    \node[below right] at (P3) {\small fixed};
    \draw (0,-0.2) arc (90:55:0.4); \node at (0.22,-0.05) {\small $\psi$};
  \end{scope}
  \begin{scope}[xshift=7.5cm]
    \draw[->] (-1.3,0) -- (1.3,0) node[right]{\small $d_{13}$};
    \draw[->] (0,-1.1) -- (0,1.3) node[above]{\small $d_{12}$};
    \draw[->] (-0.9,-0.9) -- (0.9,0.9) node[above right]{\small $d_{23}$};
    \node at (0,-1.4) {\small coordinates near $u$};
  \end{scope}
\end{tikzpicture}
\end{center}
\added{When the three active differentials are linearly independent, the three half-spaces meet in
an open octant, so the cone is non-empty and the value is regular.}

So:
\[
  \text{$\theta$ is a topologically regular value}\iff
  \begin{minipage}[c]{0.5\linewidth}
  neighbourhoods of the value have superlevel sets that deformation retract to each other.
  \end{minipage}
\]
\added{(For $\theta'<\theta$ near $\theta$, $\{\rho\ge\theta\}\hookrightarrow\{\rho\ge\theta'\}$ is
a deformation retract.)} So critical values are where the topology changes\dots

For topologically regular values, this deformation retraction is defined by the flow of some vector
field defined by $f$, the gradient flow as in the last lecture (Theorem~\ref{L9:thm:diffeo}).
So the topology can only change at an obstruction to gradient flow.
For a min-type function, this obstruction is given by a convex cone, namely the intersection
of all variational cones that increase $f$ at the parameters $x\in X$.

\begin{exercise}
\label{L10:ex:classify}
Classify all distinct realizable balanced configurations of $N$ points on a sphere.
\end{exercise}

First, list the graphs. For $3$ points, the candidate contact graphs are:
\begin{center}
\begin{tikzpicture}[scale=0.7,every node/.style={circle,fill,inner sep=1.2pt}]
  \begin{scope}
    \node (a) at (0,0) {}; \node (b) at (1,0) {}; \node (c) at (2,0) {};
  \end{scope}
  \begin{scope}[xshift=3.5cm]
    \node (a) at (0,0) {}; \node (b) at (1,0) {}; \node (c) at (2,0) {}; \draw (a)--(b);
  \end{scope}
  \begin{scope}[xshift=7cm]
    \node (a) at (0,0) {}; \node (b) at (1,0) {}; \node (c) at (2,0) {}; \draw (a)--(b)--(c);
  \end{scope}
  \begin{scope}[xshift=10.5cm]
    \node (a) at (0,0) {}; \node (b) at (1,0) {}; \node (c) at (2,0) {}; \draw (a)--(b)--(c);
    \draw (a) to[bend right=40] (c);
  \end{scope}
\end{tikzpicture}
\end{center}
For example, to be force balanced every vertex needs valence $\ge2$
\unclear{reads ``valence $=2$''}
\added{(and a vertex of valence $2$ needs its two edges opposite, i.e.\ along a great circle). For
$N=3$ only the triangle survives, realized as three points at mutual distance $2\pi/3$ on a great
circle.}

Some graphs on $4$ points:
\begin{center}
\begin{tikzpicture}[scale=0.7,every node/.style={circle,fill,inner sep=1.2pt}]
  \begin{scope}
    \foreach \i in {0,1,2,3} \node (p\i) at (\i*0.7,0) {};
  \end{scope}
  \begin{scope}[xshift=3.5cm]
    \foreach \i in {0,1,2,3} \node (q\i) at (\i*0.7,0) {};
    \draw (q0)--(q1)--(q2)--(q3); \draw (q0) to[bend right=40] (q3);
  \end{scope}
  \begin{scope}[xshift=7cm]
    \node (a) at (0,0) {}; \node (b) at (1,0) {}; \node (c) at (0.5,0.85) {}; \node (d) at (1.7,0.4) {};
    \draw (a)--(b)--(c)--(a);
  \end{scope}
  \begin{scope}[xshift=10cm]
    \node (t) at (0.7,0.8) {}; \node (l) at (0,0) {}; \node (r) at (1.4,0) {}; \node (bt) at (0.7,-0.8) {};
    \draw (t)--(l)--(bt)--(r)--(t); \draw (l)--(r);
  \end{scope}
  \begin{scope}[xshift=13cm]
    \node (k1) at (0,0) {}; \node (k2) at (1.4,0) {}; \node (k3) at (0.2,1.0) {}; \node (k4) at (0.7,0.35) {};
    \draw (k1)--(k2)--(k3)--(k1); \draw (k4)--(k1); \draw (k4)--(k2); \draw (k4)--(k3);
  \end{scope}
\end{tikzpicture}
\end{center}
\orig{``?'' by the triangle${}+{}$point and the rhombus; ``opposite'', ``great circles''. One more sketch, apparently $K_4$, drawn at right.}
For the rhombus with a diagonal, the two valence-2 vertices need opposite edges, which forces great
circles; whether such configurations are realizable is a question of geometry\dots
\added{(It is not: the two valence-3 vertices would have to be at distance $\theta$ and at distance
$2\theta$ along a great circle, forcing $\theta=2\pi/3$, and then the two valence-2 vertices would
coincide.)}
\orig{p.\,29 ends with a sketch of two great circles}

% MAT670 -- Lecture XI (scans MAT670_XI, pp. 1--14)
\chapter{Lecture 11: Critical points, jamming and entropy}
\label{L11:ch}

\section{Critical points and jamming}
\label{L11:sec:critical}

\added{Throughout, a configuration is an ordered $N$-tuple $X=(x_1,\dots,x_N)$ of points on the sphere $\Sph^2$, and $d$ denotes the geodesic (angular) distance. Let $f(X)=\min_{i<j}d(x_i,x_j)$ and write $\Conf(N;\theta)=\{X: f(X)\ge\theta\}$. A configuration in $\Conf(N;\theta)$ is the same thing as a packing of $N$ spherical caps of angular radius $\rho=\theta/2$.}

Last time we stated the following theorem, which characterizes the regular values of $f$ in terms of the geometry of the contact graph.

\begin{definition}[Contact graph]\label{L11:def:contact}
Let $X\in\Conf(N;\theta)$ with $f(X)=\theta$, where $\theta=2\rho$. The \emph{contact graph} of $X$ is the set of edges on $\Sph^2$ \added{(geodesic arcs)} joining those points $x_i,x_j$ that are at distance exactly $\theta$.
\end{definition}
\orig{crossed out: ``geodesic skeleton of'', ``struts''}

\begin{addedblock}
For an edge $\{i,j\}$ of the contact graph let $u_{ij}\in T_{x_i}\Sph^2$ be the unit tangent vector at $x_i$ pointing along the geodesic towards $x_j$ (well defined since $0<\theta<\pi$). A tangent vector $v=(v_1,\dots,v_N)$, $v_i\in T_{x_i}\Sph^2$, changes the length of the edge at rate
\[
\frac{d}{dt}\Big|_{t=0} d(x_i+tv_i,\,x_j+tv_j) \;=\; -\ip{u_{ij}}{v_i}-\ip{u_{ji}}{v_j}.
\]
We say $v$ lies in the \emph{tangent cone} (of increasing directions) at $X$ if this rate is strictly positive for every contact edge, and that $\theta$ is a \emph{regular value} of $f$ if every $X$ with $f(X)=\theta$ admits such a $v$; otherwise $\theta$ is a \emph{critical value} and the offending $X$ are \emph{critical configurations}.
\end{addedblock}

\begin{theorem}\label{L11:thm:forcebalance}
A configuration $X$ with $f(X)=\theta$ has no element in its tangent cone if and only if there exist weights $w_{ij}=w_{ji}\ge0$ on the contact edges, not all zero, such that the contact graph is \emph{force balanced}:
\[
\sum_{j:\,\{i,j\}\text{ contact}} w_{ij}\,u_{ij} \;=\; 0 \qquad\text{for every vertex } i.
\]
In particular, $\theta$ is a regular value if and only if no configuration at level $\theta$ admits such a balanced (non-negative, not all zero) system of weights.
\end{theorem}
\orig{``(non-negative, not all 0)'' replaces crossed-out ``(positive)''}

\begin{addedblock}
\emph{Proof sketch.} Let $A$ be the linear map $v\mapsto\big(-\ip{u_{ij}}{v_i}-\ip{u_{ji}}{v_j}\big)_{\{i,j\}}$ from $\bigoplus_i T_{x_i}\Sph^2$ to $\R^{E}$, $E$ the set of contact edges. By Gordan's alternative (a form of Farkas' lemma), either there is $v$ with $Av>0$ componentwise, or there is $w\in\R^E$, $w\ge0$, $w\ne0$ with $A^{T}w=0$. Computing $A^Tw$ vertex by vertex gives exactly $-\sum_j w_{ij}u_{ij}=0$. Together with the deformation theorem from last lecture (if $[a,b]$ contains no critical values, then $\Conf(N;b)$ is a deformation retract of $\Conf(N;a)$), this is the basic tool for computing the topology of $\Conf(N;\theta)$ as $\theta$ increases.
\end{addedblock}

The simplest application is to points on a circle.

\begin{theorem}\label{L11:thm:circle}
\added{For $N$ points on the circle $\Sph^1$,} $\Conf(N;0)\simeq\Conf(N;2\pi/N)$.
\end{theorem}
\unclear{Blotted symbol before $2\pi/N$; read here as the circle case. Possibly ``$<$'': the $\Sph^2$ theorem, see remark.}
\begin{addedblock}
\emph{Alternative reading.} The blot may be ``$\theta<$'', i.e.\ the statement for $\Sph^2$:
$\Conf(N;\theta)$ is a strong deformation retract of $\Conf(N;0)$ for $0\le\theta<2\pi/N$. This is
Theorem~4.14 of Kusner--Kusner--Lagarias--Shlosman (cited in \S\ref{L11:sec:examples}). It sharpens
the bound $\pi/N$ of Lecture~10, since a balanced contact graph has total length at least $2\pi$,
and the smallest critical value is attained only by $N$ equally spaced points on a great circle. The
proof sketch below covers only the circle.
\end{addedblock}

\begin{addedblock}
\emph{Proof sketch.} Here $\Conf(N;0)$ means the usual ordered configuration space (distinct points). On $\Sph^1$ the vectors $u_{ij}$ at $x_i$ are $\pm$ the unit tangent, so force balance at a vertex forces it to have one contact on each side with equal weights. Following the weights along a chain of contacts, a chain with an endpoint is unbalanced at that endpoint; hence the contact graph of a critical configuration is the whole $N$-cycle, i.e.\ $X$ is a regular $N$-gon and $\theta=2\pi/N$. Thus there are no critical values in $(0,2\pi/N)$, so $\Conf(N;\theta)\simeq\Conf(N;0)$ for $\theta<2\pi/N$. At $\theta=2\pi/N$ itself the space consists of the labelled regular $N$-gons, a disjoint union of $(N-1)!$ circles (one for each cyclic order of the labels); and $\Conf(N;0)$ has the same homotopy type, each of its $(N-1)!$ components being a circle times an open simplex.
\end{addedblock}

This also gives an algorithm for exploring critical and maximal configurations: \emph{$\rho$-evolution}. Starting from any configuration, increase $\rho$ by moving the points in a direction of the tangent cone; update the contact graph whenever new contacts form, and continue until no increasing direction exists, i.e.\ until the contact graph becomes force balanced.
\unclear{Page ends with ``update\ldots'' and ``conflicting\ldots'' (or ``conf-flip''?); not expanded.}
\added{The configuration reached is then a critical configuration, generically a local maximum of $f$ (a \emph{jammed} configuration).}

\section{Examples with pictures}
\label{L11:sec:examples}

\orig{p.~2: ``$N=3$, $N=4$: full picture'' -- pictures were presumably drawn on the board.}
For $N=3$ and $N=4$ one can give the full picture of the critical points.
\begin{addedblock}
For $N=3$: a vertex with a single contact cannot be balanced, so every vertex of a critical contact graph has degree $2$, and the two tangent vectors at it must be opposite; thus the three points lie on a great circle, equally spaced. The only critical value is the maximum $\theta=2\pi/3$, and $\Conf(3;\theta)\simeq\Conf(3;0)$ for $\theta<2\pi/3$.
For $N=4$: the maximum is the regular tetrahedron, $\theta=\arccos(-1/3)\approx109.47^\circ$. The square inscribed in a great circle ($\theta=\pi/2$) is also critical (each vertex has two opposite contacts), but it is not a local maximum: pushing alternate vertices slightly north and south increases all four contact distances to second order, deforming the square towards the tetrahedron.
\end{addedblock}

For larger $N$:
\begin{itemize}
\item For various $N$ the maximum is the solution of the \emph{Tammes problem}. \added{(The Tammes problem is solved for $N\le 14$ and $N=24$; e.g.\ $N=13$ and $N=14$ by Musin--Tarasov, $N=24$ by Robinson.)}
\item $N=10$: local maxima? \unclear{``$N=10$ -- local max?'': whether there are non-global local maxima, presumably.}
\item $N=12$: local maxima? Higher critical points. The jitterbug?
\end{itemize}
\begin{addedblock}
For $N=12$ the global maximum is the icosahedron ($\theta=\arctan 2\approx63.43^\circ$). The cuboctahedron (the kissing configuration of the FCC lattice, $\theta=60^\circ$) is a critical configuration that is not a local maximum; Buckminster Fuller's ``jitterbug'' motion deforms the cuboctahedron into the icosahedron. (See R.~Kusner, W.~Kusner, J.~Lagarias, S.~Shlosman, \emph{Configuration spaces of equal spheres touching a given sphere: the twelve spheres problem}, arXiv:1611.10297 (2016); in \emph{New Trends in Intuitive Geometry}, Bolyai Soc.\ Math.\ Stud.\ 27, Springer, 2018. That paper shows the cuboctahedral (FCC) configuration is critical. That it is not a local maximum follows from the jitterbug: along it the triangle edges lengthen from $60^\circ$ to $\arctan 2$, while the square diagonals shrink from $90^\circ$ to $\arctan 2$.)
\end{addedblock}

\section{Local maxima and Edwards entropy}
\label{L11:sec:edwards}

\begin{definition}[Edwards' ansatz]\label{L11:def:edwards}
All jammed states of equal volume have equal probability.
\end{definition}
\orig{``Edw.'' crossed out and replaced by ``Ansatz''}
\begin{addedblock}
This is the hypothesis of S.~F.~Edwards and R.~B.~S.~Oakeshott (\emph{Theory of powders}, Physica A, 1989). The \emph{Edwards entropy} at volume $V$ is the logarithm of the number (or measure) of jammed states of volume $V$; it plays the role that the Boltzmann entropy plays for the microcanonical ensemble.
\end{addedblock}

This is related to the evolution of the $\rho$-flow: \added{if configurations are chosen at random and then evolved by $\rho$-evolution, the probability of a given jammed state is the measure of its basin of attraction, and Edwards' ansatz says these basins are (for equal volumes) equally large.} It is also related to the model \ldots\ \unclear{``Also the model\ldots''; crossed out: ``ex.\ annealed''} It is not always true\,? \added{Numerical studies of small disc packings have found that jammed states are in general \emph{not} equally likely, and that the frequencies depend on the preparation protocol.}
\added{\emph{Update:} later work (Martiniani, Schrenk, Ramola, Chakraborty, Frenkel, \emph{Nature Physics}, 2017) gave numerical evidence that the ansatz does hold for soft-sphere packings exactly at the jamming threshold.}

\subsection{Boltzmann measure: the discrete setting}
\label{L11:sec:boltzmann}

Consider $N$ particles, each in one of $k$ states. Each state has an energy $E_i$, $i\in\{1,\dots,k\}$. Let $n_i$ be the number of particles in state $i$. The system is closed, so
\[
E=\sum_{i=1}^k E_i\,n_i,\qquad N=\sum_{i=1}^k n_i
\]
\orig{he writes both $n_i$ and $N_i$; unified to $n_i$}
are fixed. The number of ways to place the particles with occupation numbers $\{n_1,\dots,n_k\}$ is the multinomial coefficient
\[
W=\frac{N!}{n_1!\,n_2!\cdots n_k!}.
\]
By the law of large numbers, the optimal partition (the one maximizing $W$) dominates for $N$ large.
\added{(More precisely: $W$, as a function of the frequencies $n_i/N$, is $e^{N H+o(N)}$ with $H=-\sum (n_i/N)\ln(n_i/N)$, so frequencies whose entropy is smaller than the maximum by a fixed amount are exponentially rare.)}
So we determine the partition that maximizes $W$; since $\log$ is monotone increasing we may replace $W$ by $\log W$. By Stirling's formula,
\[
\log W=(N\ln N-N)-\sum_{i=1}^k\,(n_i\ln n_i-n_i)\;+\;\text{error}\quad\added{(\text{error }=O(k\ln N))}.
\]
We maximize subject to the two constraints with Lagrange multipliers $\alpha,\beta$:
\[
f=\ln W+\alpha\Big(N-\sum_{i=1}^k n_i\Big)+\beta\Big(E-\sum_{i=1}^k E_i n_i\Big).
\]
Since $\frac{\partial}{\partial n}(n\ln n-n)=\ln n$,
\[
0=\frac{\partial f}{\partial n_i}=-\ln n_i-(\alpha+\beta E_i),
\qquad\text{so}\qquad n_i=e^{-\alpha-\beta E_i}.
\]
\added{Summing over $i$ gives $N=e^{-\alpha}Z$, which eliminates $\alpha$:}
\[
\frac{n_i}{N}=\frac{e^{-\beta E_i}}{\sum_{j=1}^k e^{-\beta E_j}},\qquad
Z=\sum_{j=1}^k e^{-\beta E_j}\quad\added{\text{(the \emph{partition function})}}.
\]

\begin{definition}\label{L11:def:boltzmann}
Here $\beta=\dfrac1T$ \ (or $\dfrac{1}{k_BT}$ \added{with Boltzmann's constant $k_B$}) is the inverse temperature, and
\[
F=-T\ln Z
\]
is the \emph{free energy}. The \emph{Boltzmann measure} is
\[
\Pr(\text{state } i)=\frac{e^{-\beta E_i}}{Z}.
\]
\end{definition}

The mean energy per particle is
\[
\bar E=\frac EN=\sum_{i=1}^k E_i\frac{n_i}{N}
=\sum_{i=1}^k\frac{E_ie^{-\beta E_i}}{Z}
=\frac{-\frac{\partial}{\partial\beta}Z}{Z}
=-\frac{\partial}{\partial\beta}\ln Z.
\]
\fixed{page has $e^{\beta E_i}$ and final ``$-\partial_\beta Z$''; signs and $\ln$ corrected}

\begin{lemma}\label{L11:lem:concave}
$-\ln Z$ is concave \added{as a function of $\beta$}. Equivalently, $\bar E=-\frac{\partial}{\partial\beta}\ln Z$ is monotone \added{(non-increasing in $\beta$, i.e.\ non-decreasing in $T$)}.
\end{lemma}
\begin{proof}[Proof idea]
We must show $\dfrac{\partial^2}{\partial\beta^2}\ln Z\ge0$. Why? Write it as the variance of $E$, which is $\ge0$.
\begin{addedblock}
Indeed, with $\langle g\rangle=\sum_i g(E_i)e^{-\beta E_i}/Z$,
\[
\frac{\partial^2}{\partial\beta^2}\ln Z=\frac{Z''}{Z}-\Big(\frac{Z'}{Z}\Big)^2=\langle E^2\rangle-\langle E\rangle^2=\operatorname{Var}(E)\ge0,
\]
with equality only if all $E_i$ are equal.
\end{addedblock}
\end{proof}
\fixed{page has ``$>0$''; it is $\ge 0$}

Thus $-\dfrac1\beta\ln Z$ is the free energy.
\fixed{page: ``$\frac1\beta\ln Z$''; sign restored as in $F=-T\ln Z$}
As $T\to\infty$ (i.e.\ $\beta\to0$), $\Pr(i)\to\dfrac1k$: all states become equally likely. \added{As $T\to0$ the measure concentrates on the states of minimal energy.}

\begin{remark}
\added{For hard particles the energy of a configuration is $0$ if the particles do not overlap and $+\infty$ if they do. The Boltzmann weight $e^{-\beta E}$ is then the indicator of the admissible configurations, so the Boltzmann measure is normalized volume on the configuration space, at every temperature. This is why volumes of configuration spaces appear as partition functions below.}
\end{remark}

\subsection{The hard disc model in an interval}
\label{L11:sec:hardrods}

Place hard discs \added{of diameter~$1$ (that is, hard rods of length~$1$)} in the interval $[0,N]$, and let $a_1,a_2,a_3,\dots$ be the gaps \added{in front of the first, second, third, \ldots\ disc}.

\begin{center}
\begin{tikzpicture}[scale=0.9]
\draw[thick] (0,0) -- (8,0);
\draw (0,-0.15)--(0,0.15) node[below=6pt] {$0$};
\draw (8,-0.15)--(8,0.15) node[below=6pt] {$N$};
\foreach \x in {1.2,3.6,5.3} { \fill[gray!50] (\x,-0.12) rectangle (\x+1,0.12); \draw (\x,-0.12) rectangle (\x+1,0.12); }
\draw[<->] (0,0.35) -- node[above] {$a_1$} (1.2,0.35);
\draw[<->] (2.2,0.35) -- node[above] {$a_2$} (3.6,0.35);
\draw[<->] (4.6,0.35) -- node[above] {$a_3$} (5.3,0.35);
\end{tikzpicture}
\end{center}

With three discs the configuration space is
\[
\Omega_{N,3}=\{(a_1,a_2,a_3):\ 0\le a_1,a_2,a_3,\ \ a_1+a_2+a_3\le N-3\},
\]
a simplex \orig{p.~7: sketch of the simplex with axes $a_1,a_2,a_3$} with
\[
\vol\Omega_{N,3}=\frac{(N-3)^3}{3!}.
\]
In general, \added{for $k\le N$ discs,}
\[
\vol(\Omega_{N,k})=\frac{(N-k)^k}{k!}.
\]
\added{(The gaps describe the discs up to relabelling; for labelled discs multiply by $k!$.)}

Now give each particle an energy $E_0$ and define the \emph{fugacity}
\[
z=e^{-\beta E_0}.
\]
\added{(In the usual notation $z=e^{\beta\mu}$ with chemical potential $\mu$; here $\mu=-E_0$.)} The \added{grand canonical} partition function is
\[
Z_N(z)=\sum_{k=0}^N z^k\,\vol(\Omega_{N,k})=\sum_{k=0}^N\frac{(N-k)^k}{k!}\,z^k.
\]
\added{This is the classical \emph{Tonks gas} (L.~Tonks, 1936), one of the few exactly solvable models of hard particles.}

\begin{exercise}\label{L11:ex:tonks}
\added{Using Stirling's formula on the largest term, show that $\frac1N\ln Z_N(z)\to P$, where $P>0$ solves $z=Pe^{P}$, and that the dominant number of discs satisfies $k/N\to P/(1+P)$. ($P$ is $\beta$ times the pressure; note that the density $k/N$ increases from $0$ to $1$ as $z$ goes from $0$ to $\infty$.)}
\end{exercise}

\section{A quasi-one-dimensional model with open ends}
\label{L11:sec:channel}

\orig{crossed out: ``Model''}
We study discs in a narrow channel, where a combinatorial analysis is possible. Take discs of radius $r=1$ in a channel whose width $w$ satisfies
\[
2<w<2+\sqrt3 .
\]
\begin{addedblock}
The centres then lie in a strip of width $h=w-2\in(0,\sqrt3)$. Discs cannot pass each other ($h<2$). If two consecutive discs touch while lying against opposite walls, their centres are horizontally $L=\sqrt{4-h^2}$ apart, and $1<L<2$; the condition $h<\sqrt3$ is exactly $L>1$, which guarantees that each disc can touch only its two neighbours in the row (discs $i$ and $i+2$ are always more than $2$ apart). Consecutive discs touching along the same wall are horizontally $2$ apart.
\end{addedblock}

\begin{figure}[ht]
\centering
\begin{tikzpicture}[scale=0.33]
% channel width 3.2, h = 1.2, L = 1.6
\draw[thick] (-1.5,0) -- (15.1,0); \draw[thick] (-1.5,3.2) -- (15.1,3.2);
\foreach \x/\y in {0/1,1.6/2.2,3.2/1,5.2/1,6.8/2.2,8.4/1,10.4/1,12.0/2.2,13.6/1} { \draw (\x,\y) circle (1); }
\node[right] at (15.3,1.6) {\small jammed: locally maximal density};
\end{tikzpicture}

\vspace{2ex}
\begin{tikzpicture}[scale=0.33]
\draw[thick] (-1.5,0) -- (8.7,0); \draw[thick] (-1.5,3.2) -- (8.7,3.2);
\foreach \x/\y in {0/1,1.6/2.2,3.6/2.2,5.6/2.2,7.2/1} { \draw (\x,\y) circle (1); }
\draw[->,thick] (3.6,2.2) -- (3.6,0.6);
\node[right] at (9,1.6) {\small critical saddle};
\draw[<->] (19,0) -- node[right] {\small $w$} (19,3.2);
\end{tikzpicture}
\caption{Discs in a narrow channel ($h=1.2$, $L=1.6$). Top: a jammed configuration, contact word $00100100$. Bottom: three consecutive discs on the same wall (word $011\,0$); the middle one is force balanced but can be pushed into the gap, so this is a critical point but not a local maximum.}
\label{L11:fig:channel}
\end{figure}
\orig{p.~9: three sketches; third shows a narrow channel with discs in a staircase}

Let the length of the configuration be $t$. \added{(Discs are confined only by the walls; the ends are open, so ``jammed'' means that $t$ is locally minimal for the given number of discs.)}

\subsection{\texorpdfstring{Combinatorics: how many maxima for given $t$?}{Combinatorics: how many maxima for given t?}}

The length is determined by the number of adjacent pairs of discs at opposite walls \added{versus at the same wall}. Record each contact between consecutive discs by a letter:

\begin{center}
\begin{tikzpicture}[scale=0.4]
\draw (0,0) circle (1); \draw (2,0) circle (1); \draw[->] (0,0)--(2,0);
\node at (1,-1.9) {$1$: same wall};
\draw (7,0) circle (1); \draw (8.6,-1.2) circle (1); \draw[->] (7,0)--(8.6,-1.2);
\node at (7.8,-3.1) {$0$: opposite walls};
\end{tikzpicture}
\end{center}

\added{Two consecutive $1$'s are impossible in a jammed configuration: the middle disc of three discs along one wall can move away from the wall (Figure~\ref{L11:fig:channel}, bottom).} So counting jammed configurations amounts to counting binary words with no two adjacent $1$'s. This is equivalent to a domino tiling problem, or a ``stars and bars'' count: such words are sequences formed of the blocks
\[
0=\square,\qquad 01=\square\!\square\ (\text{a domino}).
\]
\added{(Precisely: prefixing a $0$, words of length $\ell$ with no $11$ correspond bijectively to tilings of a $1\times(\ell+1)$ strip by squares and dominoes.)}

This satisfies the Fibonacci recursion \added{(we count configurations up to the up/down reflection of the channel)}: a tiling of length $t$ is a tiling of length $t-1$ followed by a square, or a tiling of length $t-2$ followed by a domino; there are $1,2,\dots$ tilings of lengths $1,2,\dots$
\unclear{``(not up/down symmetric)'' -- reading as: words count configurations modulo reflection.}
\orig{crossed out: ``\# words ending in $0$ of length $\ell$''}
\added{Hence the number of tilings of length $m$ is the Fibonacci number $F_{m+1}$, and the number of binary words of length $\ell$ with no $11$ is $F_{\ell+2}$ ($F_1=F_2=1$).}

\begin{center}
\begin{tikzpicture}[scale=0.45]
\draw (0,0) rectangle (6,1); \draw (6,0) rectangle (7,1); \node[left] at (-0.2,0.5) {$t-1$};
\draw (0,-2) rectangle (5,-1); \draw (5,-2) rectangle (7,-1); \node[left] at (-0.2,-1.5) {$t-2$};
\end{tikzpicture}
\end{center}

Given exactly $k$ ones in a word of length $\ell$, we are inserting $k$ ones into a sequence of $\ell-k$ zeros,
\[
{}_{\wedge}0_{\wedge}0_{\wedge}0_{\wedge}\cdots{}_{\wedge}0_{\wedge},
\]
at most one into each of the $\ell-k+1$ slots. So for each $k$ there are $\dbinom{\ell-k+1}{k}$ choices of order. \added{(Summing over $k$ gives $\sum_k\binom{\ell-k+1}{k}=F_{\ell+2}$.)}

So if the length of a ``$1$'' is $2$ and that of a ``$0$'' is $L$, the configuration with word of length $\ell$ \added{(i.e.\ $\ell+1$ discs)} containing $k$ ones has length
\[
t=2k+(\ell-k)L+2,
\]
the $+2$ coming from the two ends.

\begin{remark}
\unclear{Check: with open ends, a $1$ in the first or last slot seems not locally length-minimizing.}
\added{With truly open ends, an end disc lying against the same wall as its neighbour can roll around that neighbour towards the opposite wall, shortening $t$. If such end configurations are excluded, the count becomes $\binom{\ell-k-1}{k}$ instead of $\binom{\ell-k+1}{k}$; the asymptotics below are unchanged.}
\end{remark}

We can also compute the maximum of $\binom{n-k+1}{k}$ over $k$ (with $n=\ell$): as $n\to\infty$ the maximizing $k$ satisfies
\[
\frac kn\longrightarrow\frac5{10}-\frac{\sqrt5}{10}=\frac12-\frac{\sqrt5}{10}\approx0.276 .
\]
\begin{addedblock}
Indeed, with $x=k/n$, Stirling gives $\frac1n\ln\binom{n-k+1}{k}\to g(x)=(1-x)\ln(1-x)-x\ln x-(1-2x)\ln(1-2x)$, and $g'(x)=0$ becomes $(1-2x)^2=x(1-x)$, i.e.\ $5x^2-5x+1=0$, whose root in $(0,\tfrac12)$ is $x=\frac{5-\sqrt5}{10}$. At this point $g(x)=\ln\varphi$, $\varphi=\frac{1+\sqrt5}2$, which matches the growth $F_{n+2}\sim\varphi^{n}$: the maximal term carries the whole exponential growth.
\end{addedblock}

And there is a ``unique'' densest packing: \added{the zigzag with all contacts at opposite walls (the word $00\cdots0$, length $\ell L+2$), unique up to reflection. At the other extreme the least dense jammed configuration alternates pairs on the two walls (word $0101\cdots$, or $1010\cdots$ with ends allowed).}
\unclear{Second sketch annotated ``$\leftarrow$ \ldots\ last \ldots\ as well'': read as the least dense case.}

\begin{center}
\begin{tikzpicture}[scale=0.28]
\draw (-1.3,0)--(14.1,0); \draw (-1.3,3.2)--(14.1,3.2);
\foreach \i in {0,...,8} { \pgfmathsetmacro{\y}{mod(\i,2)==0 ? 1 : 2.2} \draw (1.6*\i,\y) circle (1); }
\node[right] at (14.5,1.6) {\small densest: $0000\cdots$};
\end{tikzpicture}

\vspace{1ex}
\begin{tikzpicture}[scale=0.28]
\draw (-1.3,0)--(15.7,0); \draw (-1.3,3.2)--(15.7,3.2);
\foreach \x/\y in {0/1,2/1,3.6/2.2,5.6/2.2,7.2/1,9.2/1,10.8/2.2,12.8/2.2,14.4/1} { \draw (\x,\y) circle (1); }
\node[right] at (16.1,1.6) {\small least dense: $10101010$};
\end{tikzpicture}
\end{center}

\begin{exercise}\label{L11:ex:typical}
Using the law of large numbers and \added{Stirling's} approximation, show that \added{if all jammed configurations of $n+1$ discs are equally likely (Edwards' ansatz without the volume constraint),} a typical configuration has a fraction
\[
\frac12-\frac{\sqrt5}{10}
\]
of $1$'s. \added{Deduce that its length per disc is close to $2x+(1-x)L$ with $x=\frac12-\frac{\sqrt5}{10}$.} Compare to coin flips, where the fraction would be $\tfrac12$.
\end{exercise}
\orig{``Ex.\ law of large \ldots\ \& approx.\ \ldots''; partly illegible}

\subsection{Outlook}
\begin{itemize}
\item Probabilities of critical points. \unclear{First word unclear (``prune''? ``prior''?).} Flow to jammed configurations\ldots\ The resulting distribution is probably concentrated on the maxima, since critical saddles should have probability near $0$. \added{(A saddle attracts only a lower-dimensional set of starting configurations under the flow.)}
\item $\Longrightarrow$ Question: compute these values!
\item Also\ldots\ for large $N$ this is hard, so we study low-order cases.
\end{itemize}
\unclear{last line read as ``so study low order cases''}


\end{document}
